
\documentclass[12pt,a4paper]{article}

\usepackage{fontspec}
\usepackage{xeCJK}
\usepackage{polyglossia}
\setmainlanguage{russian}
\setotherlanguage{english}
\setmainfont{DejaVu Serif}
\setsansfont{DejaVu Sans}
\setmonofont{DejaVu Sans Mono}
\setCJKmainfont{Noto Serif CJK SC}

\usepackage{amsmath,amssymb,amsthm,mathtools}
\usepackage{geometry}
\geometry{left=20mm,right=20mm,top=20mm,bottom=22mm}
\usepackage{enumitem}
\setlist{itemsep=0.15em,topsep=0.25em,leftmargin=1.4em}
\usepackage{xcolor}
\usepackage{array,booktabs,longtable}
\usepackage{microtype}
\usepackage{titlesec}
\usepackage{fancyhdr}
\usepackage{hyperref}
\hypersetup{colorlinks=true,linkcolor=black,urlcolor=blue}

\pagestyle{fancy}
\setlength{\headheight}{15pt}
\fancyhf{}
\lhead{Государственный экзамен 2026}
\rhead{Кафедра ОУ ВМК МГУ}
\cfoot{\thepage}

\setlength{\parindent}{0pt}
\setlength{\parskip}{0.45em}
\allowdisplaybreaks
\sloppy
\setcounter{secnumdepth}{0}
\setcounter{tocdepth}{2}
\titleformat{\section}{\Large\bfseries}{\thesection.}{0.5em}{}
\titleformat{\subsection}{\large\bfseries}{\thesubsection.}{0.5em}{}
\titleformat{\subsubsection}{\normalsize\bfseries}{\thesubsubsection.}{0.5em}{}

\newcommand{\R}{\mathbb{R}}
\newcommand{\eps}{\varepsilon}
\newcommand{\T}{\top}
\newcommand{\dd}{\,d}
\newcommand{\norm}[1]{\left\lVert #1\right\rVert}
\newcommand{\inner}[2]{\left\langle #1,#2\right\rangle}
\newcommand{\Lie}{\operatorname{Lie}}
\newcommand{\Vect}{\operatorname{Vec}}
\newcommand{\rank}{\operatorname{rank}}
\newcommand{\ad}{\operatorname{ad}}
\newcommand{\sgn}{\operatorname{sgn}}
\newcommand{\sng}{\mathrm{sng}}
\newcommand{\PV}{\mathrm{PV}}
\newcommand{\CV}{\mathrm{CV}}
\DeclareMathOperator*{\argmax}{arg\,max}
\DeclareMathOperator*{\argmin}{arg\,min}
\newcommand{\Argmax}{\operatorname*{Arg\,max}}
\newcommand{\ticket}[2]{\section*{Билет #1. #2}\addcontentsline{toc}{section}{Билет #1. #2}}
\newcommand{\official}[1]{\subsection*{Официальная формулировка}\textbf{#1}}
\newcommand{\fullanswer}{\subsection*{1. Полная версия ответа}}
\newcommand{\cnexplain}{\subsection*{2. 中文解释}}
\newcommand{\corrections}{\subsection*{3. Основные исправления по сравнению со студенческими версиями}}
\newcommand{\memory}{\subsection*{4. Краткая версия для запоминания}}
\newcommand{\blocktitle}[3]{%
  \clearpage
  \thispagestyle{fancy}%
  \begin{center}
    {\LARGE\bfseries Блок #1. #2}\\[0.35em]
    {\large #3}
  \end{center}
  \addcontentsline{toc}{section}{Блок #1. #2}
  \vspace{0.5em}
}

\begin{document}
\begin{titlepage}
\centering
\vspace*{3cm}
{\Huge\bfseries Полные ответы к государственному экзамену\\[0.4em]кафедры оптимального управления ВМК МГУ}\\[1.2cm]
{\Large Билеты 1--29}\\[0.8cm]
{\large Объединенная версия по блокам A--I}\\[2cm]
{\large 2026}
\vfill
\end{titlepage}

\tableofcontents

\blocktitle{A}{Обратные задачи и метод динамической регуляризации}{Билеты 1--3}
\noindent\textbf{Общее замечание.} Билеты 1--3 относятся к обратным задачам восстановления неизвестных характеристик динамической системы по наблюдениям её движения. Общая схема ответа:
\begin{center}
\begin{minipage}{0.92\textwidth}
\centering
некорректная обратная задача $\Rightarrow$ нормальное управление $\Rightarrow$ дискретные наблюдения\\
$\Rightarrow$ траектория-поводырь $\Rightarrow$ регуляризованная минимизация\\
$\Rightarrow$ сходимость к нормальному управлению.
\end{minipage}
\end{center}

\section{Билет 1}

\subsection*{Официальная формулировка}

\textbf{Обратные задачи, связанные с динамическими системами, описываемыми обыкновенными дифференциальными уравнениями. Метод динамической регуляризации для задачи восстановления характеристик систем обыкновенных дифференциальных уравнений.}

\subsection*{Ответ}

Рассмотрим управляемую систему обыкновенных дифференциальных уравнений
\begin{equation}
    \dot y(t)=f(t,y(t))u(t)+g(t,y(t)),\qquad y(0)=y_0,
    \qquad 0\le t\le T,
\end{equation}
где
\[
    y(t)\in \R^n,\qquad u(t)\in P\subset \R^r.
\]
Функции $f,g$, начальное состояние $y_0$, множество $P$ и горизонт $T$ считаются известными. Множество допустимых управлений имеет вид
\begin{equation}
    U=\{u(\cdot)\in L_2^r(0,T): u(t)\in P\ \text{для почти всех }t\in[0,T]\}.
\end{equation}

\textbf{Прямая задача} состоит в нахождении траектории $y(\cdot)$ по заданному управлению $u(\cdot)$. \textbf{Обратная задача} состоит в восстановлении управления $u(\cdot)$ по наблюдаемой траектории $y(\cdot)$.

Обратная задача, вообще говоря, является некорректной. Во-первых, управление может быть неединственным. Например, если
\[
    \dot y(t)=u_1(t)-u_2(t),\qquad y(0)=1,
\]
то наблюдаемой траектории $y(t)\equiv 1$ соответствует любое управление, для которого
\[
    u_1(t)=u_2(t)\quad \text{почти всюду}.
\]
Во-вторых, восстановление управления по зашумлённым наблюдениям может быть неустойчивым.

Чтобы устранить неединственность, вводят понятие \textbf{нормального управления}. Пусть $y^r(\cdot)$ --- реальная наблюдаемая траектория, порождённая некоторым допустимым управлением. Определим множество
\begin{equation}
    U_* = \{u(\cdot)\in U: y(t,u)=y^r(t),\ 0\le t\le T\}.
\end{equation}
Нормальным управлением называется управление минимальной $L_2$-нормы:
\begin{equation}
    u_* = \operatorname*{arg\,min}_{u\in U_*}\norm{u}_{L_2^r(0,T)}.
\end{equation}
При стандартных условиях, например при липшицевости $f,g$, компактности $P$, а также выпуклости $P$, прямая задача корректна, а нормальное управление существует и единственно.

На практике точная траектория неизвестна. Имеются только дискретные наблюдения
\[
    0=t_0<t_1<\cdots<t_N=T,
\]
\begin{equation}
    |\tilde y_i-y^r(t_i)|\le \delta,
    \qquad i=0,\ldots,N.
\end{equation}
Обозначим шаг сетки
\begin{equation}
    h=\max_{0\le i\le N-1}(t_{i+1}-t_i).
\end{equation}
Требуется по данным $f,g,T,P,\{t_i\},\{\tilde y_i\},\delta$ построить управление $\tilde u_{\delta,h}(\cdot)$ такое, чтобы при согласованном стремлении $\delta,h\to0$ оно сходилось к нормальному управлению $u_*$.

Метод динамической регуляризации строит кусочно-постоянное управление. На каждом промежутке $[t_i,t_{i+1})$ полагают
\begin{equation}
    \tilde u_{\delta,h}(t)=u_i,
    \qquad u_i\in P.
\end{equation}
Вводится вспомогательная траектория, называемая \textbf{поводырём}:
\begin{equation}
    z_0=\tilde y_0,
\end{equation}
\begin{equation}
    \frac{z_{i+1}-z_i}{t_{i+1}-t_i}
    =f(t_i,\tilde y_i)u_i+g(t_i,\tilde y_i),
    \qquad i=0,\ldots,N-1.
\end{equation}

Идея метода состоит в том, чтобы выбирать управление $u_i$ так, чтобы движение поводыря сближалось с наблюдаемой траекторией. Это реализует принцип экстремального сдвига Красовского. Для устойчивости к ошибкам наблюдений вводится стабилизатор Тихонова. На $i$-м шаге минимизируют функционал
\begin{equation}
    T_i(v)=
    2\inner{z_i-\tilde y_i}{f(t_i,\tilde y_i)v+g(t_i,\tilde y_i)}
    +\alpha |v|^2,
    \qquad v\in P,
\end{equation}
где $\alpha>0$ --- параметр регуляризации. Управление $u_i\in P$ выбирается приближённо:
\begin{equation}
    T_i(u_i)\le \inf_{v\in P}T_i(v)+\varepsilon,
\end{equation}
где $\varepsilon>0$ характеризует точность решения вспомогательной задачи минимизации.

Смысл члена $\alpha |v|^2$ состоит в том, что он подавляет неустойчивость обратной задачи и одновременно направляет построение к нормальному управлению, то есть к управлению минимальной нормы.

Основная теорема сходимости имеет следующий смысл. Если
\[
    \alpha\to0,
    \qquad \delta\to0,
    \qquad h\to0,
    \qquad \varepsilon\to0
\]
и параметры согласованы, например
\begin{equation}
    \frac{\delta+h+\varepsilon}{\alpha}\to0,
\end{equation}
то построенная траектория поводыря сходится к точной наблюдаемой траектории:
\begin{equation}
    \norm{z_h-y^r}_{C[0,T]}\to0,
\end{equation}
а построенные кусочно-постоянные управления сходятся к нормальному управлению:
\begin{equation}
    \norm{\tilde u_{\delta,h}-u_*}_{L_2^r(0,T)}\to0.
\end{equation}

Таким образом, метод динамической регуляризации заменяет неустойчивое прямое восстановление управления устойчивой пошаговой процедурой. На каждом шаге используется только уже поступившая информация, поэтому метод имеет динамический характер и может применяться при последовательной обработке наблюдений.

\subsection*{Краткая версия для запоминания}

\begin{enumerate}
    \item Модель: $\dot y=f(t,y)u+g(t,y)$, $u(t)\in P$.
    \item Прямая задача: по $u(\cdot)$ найти $y(\cdot)$.
    \item Обратная задача: по наблюдаемой траектории восстановить $u(\cdot)$.
    \item Обратная задача некорректна: возможны неединственность и неустойчивость.
    \item Вводится нормальное управление минимальной $L_2$-нормы.
    \item По дискретным наблюдениям строится поводырь.
    \item На каждом шаге решается регуляризованная задача минимизации с членом $\alpha |v|^2$.
    \item При согласовании параметров построенное управление сходится к нормальному управлению.
\end{enumerate}

\section{Билет 2}

\subsection*{Официальная формулировка}

\textbf{Обратные задачи, связанные с динамическими системами, описываемыми уравнениями с частными производными. Метод динамической регуляризации для задач восстановления характеристик параболических уравнений.}

\subsection*{Ответ}

Рассмотрим типичную обратную задачу для параболического уравнения на примере уравнения теплопроводности:
\begin{equation}
\begin{cases}
    y_t(t,x)=a^2y_{xx}(t,x)+u(t,x), &(t,x)\in Q=(0,T)\times(0,l),\\
    y(t,0)=y(t,l)=0,\\
    y(0,x)=y_0(x).
\end{cases}
\end{equation}
Здесь $y(t,x)$ --- состояние системы, $u(t,x)$ --- неизвестная правая часть или распределённое управление.

\textbf{Прямая задача} состоит в том, чтобы по заданному $u(t,x)$ найти решение $y(t,x)$. \textbf{Обратная задача} состоит в восстановлении $u(t,x)$ по наблюдениям состояния $y(t,x)$.

Такая обратная задача является некорректной: восстановление правой части по зашумлённым наблюдениям неустойчиво, а решение может быть неединственным. Поэтому, как и в случае ОДУ, вводится нормальное управление, то есть управление минимальной нормы среди всех управлений, порождающих наблюдаемую траекторию.

Перейдём к абстрактной операторной форме. Пусть
\[
    H=L_2(0,l),\qquad V=H_0^1(0,l).
\]
Имеем цепочку вложений
\begin{equation}
    V\subset H\simeq H^*\subset V^*,
\end{equation}
где вложение $V\subset H$ непрерывно, плотно и компактно.

Определим оператор
\[
    A:V\to V^*
\]
по формуле
\begin{equation}
    \langle Ay,v\rangle_{V^*,V}
    =a^2\int_0^l y_x(x)v_x(x)\dd x,
    \qquad y,v\in V.
\end{equation}
Оператор $A$ является симметричным и положительно определённым в энергетическом смысле.

Тогда уравнение теплопроводности записывается как
\begin{equation}
    \dot y(t)+Ay(t)=u(t),\qquad y(0)=y_0.
\end{equation}
Равенство понимается в пространстве $V^*$. При
\[
    u\in L_2(0,T;H),\qquad y_0\in H
\]
естественный класс слабых решений имеет вид
\begin{equation}
    y\in L_2(0,T;V),\qquad \dot y\in L_2(0,T;V^*).
\end{equation}
Отсюда следует
\begin{equation}
    y\in C([0,T];H),
\end{equation}
поэтому значения $y(t_i)\in H$ корректно определены.

Пусть допустимые управления задаются множеством
\begin{equation}
    U=\{u(\cdot)\in L_2(0,T;H): u(t)\in P\ \text{для почти всех }t\in[0,T]\},
\end{equation}
где $P\subset H$ --- выпуклое, замкнутое и ограниченное множество.

Пусть $y^r(t)$ --- точная наблюдаемая траектория. Множество управлений, порождающих эту траекторию, обозначим
\begin{equation}
    U_* = \{u\in U:y(t,u)=y^r(t),\ 0\le t\le T\}.
\end{equation}
Нормальное управление определяется как
\begin{equation}
    u_* = \operatorname*{arg\,min}_{u\in U_*}\norm{u}_{L_2(0,T;H)}.
\end{equation}

На практике точная траектория неизвестна. В дискретные моменты
\[
    0=t_0<t_1<\cdots<t_N=T
\]
имеются наблюдения
\begin{equation}
    \norm{y^r(t_i)-\tilde y_i}_H\le \delta,
    \qquad i=0,\ldots,N.
\end{equation}
Положим
\begin{equation}
    h=\max_{0\le i\le N-1}(t_{i+1}-t_i).
\end{equation}

Метод динамической регуляризации строит кусочно-постоянное управление
\begin{equation}
    \tilde u_{\delta,h}(t)=u_i,
    \qquad t\in[t_i,t_{i+1}),
    \qquad u_i\in P.
\end{equation}

Вводится траектория-поводырь $z_h(t)$, которая на каждом промежутке удовлетворяет уравнению
\begin{equation}
    \dot z_h(t)+Az_h(t)=u_i,
    \qquad t\in(t_i,t_{i+1}),
\end{equation}
с начальным условием, полученным на предыдущем шаге.

На $i$-м шаге управление выбирается из регуляризованной задачи минимизации:
\begin{equation}
    \varphi_i(v)
    =2\inner{z_h(t_i)-\tilde y_i}{v}_H
    +\alpha\norm{v}_H^2
    \to \min_{v\in P},
\end{equation}
где $\alpha>0$ --- параметр регуляризации. Управление $u_i\in P$ выбирается с точностью $\varepsilon>0$:
\begin{equation}
    \varphi_i(u_i)\le \inf_{v\in P}\varphi_i(v)+\varepsilon.
\end{equation}

Первое слагаемое реализует идею экстремального сдвига: поводырь должен сближаться с наблюдаемой траекторией. Второе слагаемое $\alpha\norm{v}_H^2$ является стабилизатором Тихонова и обеспечивает устойчивость восстановления.

Смысл теоремы сходимости следующий. Если
\[
    \delta\to0,
    \qquad h\to0,
    \qquad \varepsilon\to0,
    \qquad \alpha\to0
\]
и параметры согласованы, то
\begin{equation}
    z_h\to y^r \quad\text{в } C([0,T];H),
\end{equation}
а построенные управления сходятся к нормальному управлению:
\begin{equation}
    \tilde u_{\delta,h}\to u_* \quad\text{в } L_2(0,T;H).
\end{equation}
В студенческих материалах используется условие вида
\begin{equation}
    \frac{\sqrt h+\delta+\varepsilon}{\alpha}\to0,
\end{equation}
но точная форма оценки требует сверки с оригинальным текстом Осипова--Васильева--Потапова и поэтому помечается как требующая проверки.

Итак, метод динамической регуляризации для параболического уравнения заменяет неустойчивое восстановление правой части $u(t,x)$ устойчивой пошаговой процедурой, основанной на траектории-поводыре и регуляризованной минимизации.

\subsection*{Краткая версия для запоминания}

\begin{enumerate}
    \item Модель: $y_t=a^2y_{xx}+u$ с нулевыми граничными условиями.
    \item Обратная задача: восстановить $u(t,x)$ по наблюдениям $y(t,x)$.
    \item Используется тройка $V\subset H\simeq H^*\subset V^*$, где $V=H_0^1(0,l)$, $H=L_2(0,l)$.
    \item Оператор $A:V\to V^*$ задаётся формой $\langle Ay,v\rangle=a^2\int_0^l y_xv_x\dd x$.
    \item Уравнение: $\dot y+Ay=u$.
    \item Строится поводырь $\dot z_h+Az_h=u_i$.
    \item Управление $u_i\in P$ выбирается из регуляризованной задачи $2\langle z_h(t_i)-\tilde y_i,v\rangle_H+\alpha\norm{v}_H^2\to\min$.
    \item При согласовании параметров $z_h\to y^r$, $\tilde u\to u_*$.
\end{enumerate}

\section{Билет 3}

\subsection*{Официальная формулировка}

\textbf{Обратные задачи, связанные с динамическими системами, описываемыми уравнениями с частными производными. Метод динамической регуляризации для задач восстановления характеристик гиперболических уравнений.}

\subsection*{Ответ}

Рассмотрим типичную обратную задачу для гиперболического уравнения на примере волнового уравнения:
\begin{equation}
\begin{cases}
    y_{tt}(t,x)=a^2y_{xx}(t,x)+u(t,x), &(t,x)\in Q=(0,T)\times(0,l),\\
    y(t,0)=y(t,l)=0,\\
    y(0,x)=y_0(x),\\
    y_t(0,x)=y_1(x).
\end{cases}
\end{equation}
Здесь $u(t,x)$ --- неизвестная правая часть или распределённое управление.

\textbf{Прямая задача} состоит в том, чтобы по заданному $u(t,x)$ найти решение $y(t,x)$. \textbf{Обратная задача} состоит в восстановлении $u(t,x)$ по наблюдениям состояния $y$ или скорости $y_t$.

В отличие от параболического случая, уравнение является вторым порядком по времени, поэтому естественное состояние системы есть пара
\begin{equation}
    X(t)=(y(t),\dot y(t)).
\end{equation}

Введём пространства
\[
    H=L_2(0,l),\qquad V=H_0^1(0,l),
\]
и цепочку вложений
\begin{equation}
    V\subset H\simeq H^*\subset V^*.
\end{equation}
Оператор
\[
    A:V\to V^*
\]
определяется по формуле
\begin{equation}
    \langle Ay,v\rangle_{V^*,V}
    =a^2\int_0^l y_x(x)v_x(x)\dd x,
    \qquad y,v\in V.
\end{equation}
Тогда волновое уравнение записывается в абстрактной форме
\begin{equation}
    \ddot y(t)+Ay(t)=u(t),
\end{equation}
\begin{equation}
    y(0)=y_0,
    \qquad \dot y(0)=y_1.
\end{equation}
Равенство понимается в пространстве $V^*$.

При
\[
    y_0\in V,
    \qquad y_1\in H,
    \qquad u\in L_2(0,T;H)
\]
естественный энергетический класс решений имеет вид
\begin{equation}
    y\in C([0,T];V),
    \qquad \dot y\in C([0,T];H),
\end{equation}
а уравнение выполняется в $V^*$.

Множество допустимых управлений зададим как
\begin{equation}
    U=\{u(\cdot)\in L_2(0,T;H): u(t)\in P\ \text{для почти всех }t\in[0,T]\},
\end{equation}
где $P\subset H$ --- выпуклое, замкнутое и ограниченное множество.

Если $y^r$ --- точное наблюдаемое движение, то множество управлений, порождающих его, обозначается
\begin{equation}
    U_* = \{u\in U:y(t,u)=y^r(t),\ 0\le t\le T\}.
\end{equation}
Нормальное управление определяется как управление минимальной нормы:
\begin{equation}
    u_* = \operatorname*{arg\,min}_{u\in U_*}\norm{u}_{L_2(0,T;H)}.
\end{equation}

На практике точное движение неизвестно. Например, если наблюдается скорость, то в моменты
\[
    0=t_0<t_1<\cdots<t_N=T
\]
имеются данные
\begin{equation}
    \norm{\dot y^r(t_i)-\tilde w_i}_H\le \delta.
\end{equation}
Если наблюдается само состояние, то вместо этого пишут
\begin{equation}
    \norm{y^r(t_i)-\tilde y_i}_H\le \delta.
\end{equation}
В экзаменационном ответе важно указать, какой именно тип наблюдений используется.

Метод динамической регуляризации строит кусочно-постоянное управление
\begin{equation}
    \tilde u_{\delta,h}(t)=u_i,
    \qquad t\in[t_i,t_{i+1}),
    \qquad u_i\in P,
\end{equation}
где
\begin{equation}
    h=\max_i(t_{i+1}-t_i).
\end{equation}

Вводится траектория-поводырь $z_h(t)$, удовлетворяющая на каждом промежутке уравнению
\begin{equation}
    \ddot z_h(t)+Az_h(t)=u_i,
    \qquad t\in(t_i,t_{i+1}).
\end{equation}
Начальные данные $z_h(t_i)$, $\dot z_h(t_i)$ берутся из предыдущего шага.

Если наблюдается скорость, то естественная регуляризованная задача выбора управления имеет вид
\begin{equation}
    \varphi_i(v)
    =2\inner{\dot z_h(t_i)-\tilde w_i}{v}_H
    +\alpha\norm{v}_H^2
    \to \min_{v\in P},
\end{equation}
где $\alpha>0$ --- параметр регуляризации. Управление $u_i\in P$ выбирается с точностью $\varepsilon>0$:
\begin{equation}
    \varphi_i(u_i)\le \inf_{v\in P}\varphi_i(v)+\varepsilon.
\end{equation}

Если наблюдается состояние $y(t_i)$, то записывается эквивалентная форма алгоритма, согласованная с выбранным способом наблюдения; в любом случае управление выбирается пошагово из задачи экстремального сдвига с добавлением стабилизатора Тихонова $\alpha\norm{v}_H^2$.

Смысл этого стабилизатора тот же, что и в случае ОДУ и параболического уравнения: он подавляет неустойчивость обратной задачи и направляет решение к нормальному управлению минимальной нормы.

Сходимость формулируется в энергетических нормах. При согласованном стремлении
\[
    \delta\to0,
    \qquad h\to0,
    \qquad \varepsilon\to0,
    \qquad \alpha\to0
\]
получают
\begin{equation}
    (z_h,\dot z_h)\to (y^r,\dot y^r)
\end{equation}
в соответствующем пространстве траекторий, например в энергетическом смысле, и
\begin{equation}
    \tilde u_{\delta,h}\to u_* \quad\text{в } L_2(0,T;H).
\end{equation}
Точная форма оценки и условия согласования параметров требуют сверки с оригинальным текстом Осипова--Васильева--Потапова или с работами Осипова--Кряжимского--Максимова по системам с распределёнными параметрами.

Итак, для гиперболического уравнения метод динамической регуляризации также заменяет некорректное восстановление неизвестной правой части устойчивой динамической процедурой. Главное отличие от параболического случая состоит в том, что состояние включает как $y$, так и $\dot y$, а сходимость должна формулироваться в энергетическом пространстве.

\subsection*{Краткая версия для запоминания}

\begin{enumerate}
    \item Модель: $y_{tt}=a^2y_{xx}+u$ с нулевыми граничными условиями.
    \item Обратная задача: восстановить $u(t,x)$ по наблюдениям состояния или скорости.
    \item Поскольку уравнение второго порядка по времени, состояние есть пара $(y,\dot y)$.
    \item Энергетическое пространство: $y\in C([0,T];H_0^1)$, $\dot y\in C([0,T];L_2)$.
    \item Абстрактная форма: $\ddot y+Ay=u$, где $A:V\to V^*$.
    \item Строится поводырь $\ddot z_h+Az_h=u_i$.
    \item Управление выбирается из задачи экстремального сдвига с регуляризатором $\alpha\norm{v}_H^2$.
    \item Сходимость формулируется в энергетических нормах, а управление сходится к нормальному управлению.
\end{enumerate}

\section*{Замечания по надёжности блока A}

\begin{enumerate}
    \item Для билета 1 основная структура и формулы непосредственно соответствуют первой главе книги Осипова--Васильева--Потапова, посвящённой обратной задаче для систем ОДУ.
    \item Для билетов 2--3 структура ответа соответствует главам книги, указанным в оглавлении: восстановление правой части параболического и гиперболического уравнений. Однако точные оценки сходимости и условия согласования параметров требуют дополнительной сверки с полным текстом этих глав.
    \item В билетах 2--3 принципиально важно не смешивать сильную и слабую формы оператора $A$ и правильно использовать тройку $V\subset H\simeq H^*\subset V^*$.
    \item В билете 3 принципиально важно не смешивать наблюдение состояния $y$ и наблюдение скорости $\dot y$. Если в преподавательских материалах используется только один вариант наблюдения, итоговый ответ следует привести к этому варианту.
\end{enumerate}

\blocktitle{B}{Экономические модели и задачи оптимального управления на бесконечном горизонте}{Билеты 4--8}
\section*{Общие замечания по блоку B}
\addcontentsline{toc}{section}{Общие замечания по блоку B}

В блок входят пять вопросов:
\begin{itemize}
    \item Билет 4: одномерная модель Рамсея с производственной функцией Кобба--Дугласа;
    \item Билет 5: двумерная двухсекторная модель распределения ресурсов с функцией Кобба--Дугласа;
    \item Билет 6: принцип максимума Понтрягина на бесконечном интервале времени;
    \item Билет 7: метод Беллмана, ПМП и текущие сопряженные переменные;
    \item Билет 8: достаточные условия оптимальности на бесконечном горизонте.
\end{itemize}

Основные источники для блока: работы С.М. Асеева и А.В. Кряжимского по принципу максимума на бесконечном горизонте и задачам экономического роста, пособие Aseev \textit{Infinite-Horizon Optimal Control with Applications in Growth Theory}, классическая книга Л.С. Понтрягина и соавторов по математической теории оптимальных процессов, материал Ю.Н. Киселева по динамическому программированию, а также конспект по двумерной задаче распределения ресурсов с функцией Кобба--Дугласа.

\newpage

\section{Билет 4}

\subsection*{Официальная формулировка}

\textbf{4. Одномерная модель Рамсея на бесконечном горизонте планирования с производственной функцией в форме Кобба--Дугласа. Задача о наилучших пропорциях производства и потребления. Особый режим. Качественный характер решения.}

\subsection{Смысл задачи}

Рассматривается замкнутая экономика, в которой произведенный продукт распределяется между потреблением и инвестициями. Управление задает долю выпуска, направляемую на накопление капитала. Нужно найти оптимальную пропорцию между текущим потреблением и инвестициями на бесконечном горизонте планирования.

\subsection{Постановка модели}

Пусть
\[
    x(t)>0
\]
-- удельный капитал, а
\[
    u(t)\in[0,1]
\]
-- доля выпуска, направляемая в инвестиции. Тогда \(1-u(t)\) -- доля выпуска, направляемая на потребление.

Производственная функция Кобба--Дугласа в интенсивной форме:
\[
    F(x)=Ax^{\eps},\qquad A>0,\quad 0<\eps<1.
\]
Динамика капитала:
\[
    \dot x(t)=u(t)F(x(t))-\mu x(t),\qquad x(0)=x_0>0,
\]
где \(\mu>0\) -- коэффициент выбытия капитала. Если модель получена из переменных совокупного капитала \(K\) и труда \(L\), то \(\mu\) может включать как амортизацию, так и рост трудовых ресурсов.

Функционал:
\[
    J(u)=\int_0^\infty e^{-\nu t}(1-u(t))F(x(t))\,dt\to\max,
\]
где \(\nu>0\) -- коэффициент дисконтирования.

\subsection{Принцип максимума и функция переключения}

Гамильтониан в приведенных к начальному моменту переменных:
\[
    H=e^{-\nu t}(1-u)F(x)+\psi\bigl(uF(x)-\mu x\bigr).
\]
Выделим часть, зависящую от \(u\):
\[
    H=e^{-\nu t}F(x)-\psi\mu x+uF(x)\bigl(\psi-e^{-\nu t}\bigr).
\]
Так как \(F(x)>0\), функция переключения имеет вид
\[
    \Phi(t)=\psi(t)-e^{-\nu t}.
\]
Следовательно,
\[
    u^*(t)=
    \begin{cases}
        1, & \Phi(t)>0,\\
        0, & \Phi(t)<0,\\
        \text{не определяется из условия максимума}, & \Phi(t)=0.
    \end{cases}
\]

\subsection{Текущая сопряженная переменная и особый режим}

Введем текущую сопряженную переменную
\[
    p(t)=e^{\nu t}\psi(t).
\]
Тогда текущий гамильтониан:
\[
    M(x,u,p)=(1-u)F(x)+p\bigl(uF(x)-\mu x\bigr).
\]
Условие максимума зависит от знака \(p-1\):
\[
    u^*(t)=
    \begin{cases}
        1, & p(t)>1,\\
        0, & p(t)<1,\\
        \text{особый режим}, & p(t)=1.
    \end{cases}
\]

На особом режиме
\[
    p(t)\equiv1.
\]
Уравнение для текущей сопряженной переменной имеет вид
\[
    \dot p=\nu p-M_x(x,u,p).
\]
Так как
\[
    M_x=(1-u)F'(x)+p\bigl(uF'(x)-\mu\bigr),
\]
то при \(p=1\) получаем
\[
    M_x=F'(x)-\mu.
\]
Из \(\dot p=0\) следует
\[
    0=\nu+\mu-F'(x).
\]
Поэтому особое состояние \(x_s\) определяется уравнением
\[
    F'(x_s)=\mu+\nu.
\]
Для \(F(x)=Ax^\eps\):
\[
    A\eps x_s^{\eps-1}=\mu+\nu,
\]
откуда
\[
    x_s=\left(\frac{A\eps}{\mu+\nu}\right)^{\frac{1}{1-\eps}}.
\]
Особое управление определяется условием стационарности капитала:
\[
    \dot x=0,
\]
то есть
\[
    u_sF(x_s)=\mu x_s.
\]
Следовательно,
\[
    u_s=\frac{\mu x_s}{F(x_s)}.
\]
Для функции Кобба--Дугласа это дает
\[
    u_s=\frac{\mu\eps}{\mu+\nu}.
\]

\subsection{Частный случай \texorpdfstring{\(F(x)=\sqrt{x}\)}{F(x)=sqrt(x)}}

В часто используемом частном случае
\[
    A=1,
    \qquad
    \eps=\frac12,
\]
имеем
\[
    F(x)=\sqrt{x}.
\]
Тогда
\[
    F'(x)=\frac{1}{2\sqrt{x}}.
\]
Из условия \(F'(x_s)=\mu+\nu\):
\[
    \frac{1}{2\sqrt{x_s}}=\mu+\nu.
\]
Значит,
\[
    x_s=\frac{1}{4(\mu+\nu)^2}.
\]
Особое управление:
\[
    u_s=\frac{\mu x_s}{\sqrt{x_s}}=\mu\sqrt{x_s}=\frac{\mu}{2(\mu+\nu)}.
\]
На особом режиме
\[
    x(t)\equiv x_s,
    \qquad
    u(t)\equiv u_s.
\]

\subsection{Качественный характер решения}

Если начальный капитал меньше особого уровня,
\[
    x_0<x_s,
\]
то сначала выбирается максимальное инвестирование
\[
    u=1,
\]
чтобы увеличить капитал до \(x_s\), после чего траектория идет по особому режиму.

Если
\[
    x_0>x_s,
\]
то сначала выбирается
\[
    u=0,
\]
то есть весь выпуск направляется на потребление, капитал уменьшается до \(x_s\), затем начинается особый режим.

Если
\[
    x_0=x_s,
\]
то особый режим реализуется сразу:
\[
    x(t)\equiv x_s,
    \qquad
    u(t)\equiv u_s.
\]

Экономический смысл условия
\[
    F'(x_s)=\mu+\nu
\]
состоит в том, что на оптимальной магистрали предельная производительность капитала равна сумме дисконтирования и выбытия капитала. Это баланс между текущим потреблением и накоплением.

\subsection{Краткая версия для запоминания}

\[
    \dot x=uF(x)-\mu x,
    \qquad
    J=\int_0^\infty e^{-\nu t}(1-u)F(x)\,dt\to\max.
\]
\[
    M=(1-u)F(x)+p(uF(x)-\mu x).
\]
Функция переключения в текущих переменных: \(p-1\).
\[
    p>1\Rightarrow u=1,
    \qquad
    p<1\Rightarrow u=0.
\]
Особый режим:
\[
    p=1,
    \qquad
    F'(x_s)=\mu+\nu,
    \qquad
    u_sF(x_s)=\mu x_s.
\]
Для \(F(x)=\sqrt{x}\):
\[
    x_s=\frac{1}{4(\mu+\nu)^2},
    \qquad
    u_s=\frac{\mu}{2(\mu+\nu)}.
\]

\newpage

\section{Билет 5}

\subsection*{Официальная формулировка}

\textbf{5. Двумерная задача распределения ресурсов в двухсекторной экономической модели с производственной функцией Кобба--Дугласа на конечном отрезке времени с максимизацией второй фазовой координаты в конечный момент времени. Особый режим. Качественный характер решения.}

\subsection{Постановка задачи}

Рассмотрим задачу оптимального распределения ресурсов в двухсекторной экономической модели:
\[
    \dot x_1(t)=u(t)F(x(t)),
    \qquad
    x_1(0)=x_{10}>0,
\]
\[
    \dot x_2(t)=(1-u(t))F(x(t)),
    \qquad
    x_2(0)=x_{20}>0,
\]
\[
    0\le u(t)\le 1,
    \qquad
    0\le t\le T.
\]
Здесь
\[
    x(t)=(x_1(t),x_2(t))\in\R^2_+
\]
характеризует уровни развития двух секторов экономики. Управление \(u(t)\) есть доля выпуска, направляемая в первый сектор, а \(1-u(t)\) -- доля выпуска, направляемая во второй сектор.

Целевой функционал:
\[
    J[u]=x_2(T)\to\max.
\]
Требуется максимизировать уровень развития второго сектора в конечный момент времени \(T\).

Производственная функция Кобба--Дугласа:
\[
    F(x)=Ax_1^{\eps_1}x_2^{\eps_2},
\]
где
\[
    A>0,
    \qquad
    \eps_1>0,
    \qquad
    \eps_2>0,
    \qquad
    \eps_1+\eps_2=1.
\]
Постоянный множитель \(A>0\) можно устранить заменой масштаба времени, поэтому далее можно считать
\[
    A=1,
    \qquad
    F(x)=x_1^{\eps_1}x_2^{\eps_2}.
\]

Так как \(\eps_1+\eps_2=1\), функция \(F\) положительно однородна степени 1:
\[
    F(\lambda x)=\lambda F(x),
    \qquad
    \lambda>0.
\]
Кроме того, функция Кобба--Дугласа вогнута на \(\R^2_+\). Это свойство используется при доказательстве оптимальности экстремального процесса.

\subsection{Принцип максимума Понтрягина}

Так как функционал имеет тип Майера,
\[
    \Phi(x)=x_2,
    \qquad
    \Phi'(x)=
    \begin{pmatrix}
        0\\1
    \end{pmatrix},
\]
терминальное условие для сопряженных переменных:
\[
    \psi_1(T)=0,
    \qquad
    \psi_2(T)=1.
\]

Гамильтониан:
\[
    H(x,\psi,u)=\psi_1uF(x)+\psi_2(1-u)F(x).
\]
Преобразуем:
\[
    H(x,\psi,u)=F(x)\bigl(u(\psi_1-\psi_2)+\psi_2\bigr).
\]
Введем функцию переключения
\[
    \Pi(\psi)=\psi_1-\psi_2.
\]
Тогда
\[
    H(x,\psi,u)=F(x)\bigl(u\Pi(\psi)+\psi_2\bigr).
\]
Поскольку \(F(x)>0\) в положительном квадранте, условие максимума по \(u\in[0,1]\) дает
\[
    u^*(t)=
    \begin{cases}
        1, & \Pi(t)>0,\\
        0, & \Pi(t)<0,\\
        \text{не определяется однозначно}, & \Pi(t)=0.
    \end{cases}
\]
Если \(\Pi(t)\equiv0\) на некотором интервале, то возникает особый режим.

Сопряженная система:
\[
    \dot\psi_1(t)=-F_{x_1}'(x(t))\bigl(u(t)\Pi(t)+\psi_2(t)\bigr),
\]
\[
    \dot\psi_2(t)=-F_{x_2}'(x(t))\bigl(u(t)\Pi(t)+\psi_2(t)\bigr).
\]

\subsection{Особый режим}

Пусть на некотором интервале выполнено
\[
    \Pi(t)\equiv0.
\]
Тогда
\[
    \psi_1(t)=\psi_2(t),
    \qquad
    \dot\psi_1(t)=\dot\psi_2(t).
\]
Из сопряженной системы получаем
\[
    F_{x_1}'(x(t))=F_{x_2}'(x(t)).
\]
Для функции Кобба--Дугласа
\[
    F_{x_1}'(x)=\frac{\eps_1}{x_1}F(x),
    \qquad
    F_{x_2}'(x)=\frac{\eps_2}{x_2}F(x).
\]
Следовательно,
\[
    \frac{\eps_1}{x_1}F(x)=\frac{\eps_2}{x_2}F(x).
\]
Так как \(F(x)>0\), получаем
\[
    \frac{x_1}{\eps_1}=\frac{x_2}{\eps_2}.
\]
Значит, особая магистраль в исходных координатах:
\[
    L_{\sng}=\left\{x\in\R^2_+:
    \frac{x_1}{\eps_1}=\frac{x_2}{\eps_2}\right\}.
\]

Чтобы траектория оставалась на этой магистрали, необходимо
\[
    \frac{\dot x_1}{\eps_1}=\frac{\dot x_2}{\eps_2}.
\]
Подставляя уравнения движения,
\[
    \frac{u}{\eps_1}=\frac{1-u}{\eps_2}.
\]
Отсюда
\[
    \eps_2u=\eps_1(1-u),
\]
\[
    u(\eps_1+\eps_2)=\eps_1.
\]
Так как \(\eps_1+\eps_2=1\), получаем особое управление
\[
    u_{\sng}=\eps_1.
\]
Итак, особый режим задается условиями
\[
    \frac{x_1}{\eps_1}=\frac{x_2}{\eps_2},
    \qquad
    u(t)=\eps_1.
\]

\subsection{Нормированные координаты}

Для геометрического описания удобно ввести координаты
\[
    y_1=\frac{x_1}{\eps_1},
    \qquad
    y_2=\frac{x_2}{\eps_2}.
\]
Тогда
\[
    x_1=\eps_1y_1,
    \qquad
    x_2=\eps_2y_2.
\]
После устранения постоянного множителя в производственной функции система принимает вид
\[
    \dot y_1=\frac{u}{\eps_1}F(y),
\]
\[
    \dot y_2=\frac{1-u}{\eps_2}F(y).
\]
Цель \(x_2(T)\to\max\) эквивалентна цели \(y_2(T)\to\max\), так как
\[
    x_2(T)=\eps_2y_2(T).
\]
В координатах \(y\) особая магистраль имеет простой вид
\[
    L_{\sng}^y=\{y\in\R^2_+: y_1=y_2\}.
\]

В этих координатах удобно положить
\[
    q_1=\frac{\psi_1}{\eps_1},
    \qquad
    q_2=\frac{\psi_2}{\eps_2},
    \qquad
    \pi=q_1-q_2.
\]
Тогда
\[
    H(y,\psi,u)=F(y)(u\pi+q_2).
\]
Следовательно,
\[
    \pi>0\Rightarrow u=1,
    \qquad
    \pi<0\Rightarrow u=0,
\]
а на особой дуге
\[
    \pi\equiv0,
    \qquad
    y_1=y_2,
    \qquad
    u_{\sng}=\eps_1.
\]

\subsection{Обоснование оптимальности}

Рассмотрим максимизированный гамильтониан
\[
    M(x,\psi)=\max_{u\in[0,1]}H(x,\psi,u).
\]
Так как
\[
    H(x,\psi,u)=F(x)(u\Pi+\psi_2),
\]
получаем
\[
    M(x,\psi)=F(x)g(\psi),
\]
где
\[
    g(\psi)=\psi_2+\Pi^+,
    \qquad
    \Pi^+=\max\{\Pi,0\}.
\]
Эквивалентно,
\[
    g(\psi)=\max\{\psi_1,\psi_2\}
    =\frac{\psi_1+\psi_2+|\psi_1-\psi_2|}{2}.
\]
Для экстремальной тройки в рассматриваемой задаче получается
\[
    g(\psi(t))>0,
    \qquad
    0\le t\le T.
\]

Пусть \((x(t),u(t),\psi(t))\) -- экстремальная тройка, а \((\widehat x(t),\widehat u(t))\) -- произвольный допустимый процесс с тем же начальным условием. Обозначим
\[
    \Delta x(t)=\widehat x(t)-x(t),
    \qquad
    \Delta J=J[\widehat u]-J[u].
\]
Так как
\[
    J[u]=x_2(T),
    \qquad
    \psi(T)=
    \begin{pmatrix}
        0\\1
    \end{pmatrix},
\]
то
\[
    \Delta J=(\psi(T),\Delta x(T))-(\psi(0),\Delta x(0)).
\]
Поскольку \(\Delta x(0)=0\),
\[
    \Delta J=\int_0^T \frac{d}{dt}(\psi(t),\Delta x(t))\,dt.
\]
Используя сопряженную систему и условие максимума, получаем
\[
    \Delta J\le
    \int_0^T
    \left[ F(\widehat x(t))-F(x(t))-\bigl(F'(x(t)),\widehat x(t)-x(t)\bigr)\right]
    g(\psi(t))\,dt.
\]
Так как функция Кобба--Дугласа \(F\) вогнута на \(\R^2_+\),
\[
    F(\widehat x)-F(x)-\bigl(F'(x),\widehat x-x\bigr)\le0.
\]
Кроме того, \(g(\psi(t))>0\). Следовательно,
\[
    \Delta J\le0.
\]
То есть
\[
    J[\widehat u]\le J[u]
\]
для любого допустимого управления \(\widehat u\). Значит, экстремальный процесс является оптимальным.

\subsection{Качественный характер решения}

Для достаточно большого горизонта планирования оптимальная траектория имеет структуру
\[
    \text{граничный участок}
    \quad\longrightarrow\quad
    \text{особая магистраль}
    \quad\longrightarrow\quad
    \text{терминальный участок}.
\]
В координатах \(y\) особая магистраль есть прямая
\[
    y_1=y_2.
\]

Если начальная точка уже лежит на особой магистрали,
\[
    y^0\in L_{\sng},
\]
то сначала движение идет по особой магистрали:
\[
    u(t)=\eps_1,
\]
а затем на терминальном участке
\[
    u(t)=0,
\]
поскольку конечная цель состоит в максимизации второй координаты \(y_2(T)\).

Если начальная точка лежит ниже особой магистрали, то сначала надо увеличить относительный уровень второго сектора, поэтому начальный участок имеет вид
\[
    u(t)=0.
\]
Затем траектория выходит на особую магистраль и движется с управлением
\[
    u(t)=\eps_1.
\]
На последнем участке снова
\[
    u(t)=0.
\]

Если начальная точка лежит выше особой магистрали, то сначала надо увеличить относительный уровень первого сектора, поэтому начальный участок имеет вид
\[
    u(t)=1.
\]
Затем следует особая дуга
\[
    u(t)=\eps_1,
\]
а в конце -- терминальный участок
\[
    u(t)=0.
\]

Если \(\tau\) -- момент входа на особую магистраль, а \(\theta\) -- момент выхода с нее на терминальный участок, то структура управления записывается так:
\[
    u(t)=
    \begin{cases}
        0\ \text{или}\ 1, & 0\le t<\tau,\\
        \eps_1, & \tau<t<\theta,\\
        0, & \theta<t\le T.
    \end{cases}
\]
Для случая, когда начальная точка уже лежит на особой магистрали, \(\tau=0\), и структура упрощается:
\[
    u(t)=
    \begin{cases}
        \eps_1, & 0\le t<\theta,\\
        0, & \theta<t\le T.
    \end{cases}
\]

Экономический смысл: сначала ресурсы перераспределяются так, чтобы вывести экономику на сбалансированную магистраль
\[
    \frac{x_1}{\eps_1}=\frac{x_2}{\eps_2}.
\]
Затем экономика движется по особому режиму, где доля ресурсов, направляемая в первый сектор, равна его эластичности:
\[
    u_{\sng}=\eps_1.
\]
На заключительном участке все ресурсы направляются во второй сектор:
\[
    u=0,
\]
так как именно \(x_2(T)\) максимизируется.

\subsection{Краткая версия для запоминания}

\[
    \dot x_1=uF(x),
    \qquad
    \dot x_2=(1-u)F(x),
    \qquad
    J=x_2(T)\to\max.
\]
\[
    F(x)=Ax_1^{\eps_1}x_2^{\eps_2},
    \qquad
    \eps_1+\eps_2=1.
\]
\[
    H=F(x)(u(\psi_1-\psi_2)+\psi_2).
\]
\[
    \Pi=\psi_1-\psi_2.
\]
\[
    \Pi>0\Rightarrow u=1,
    \qquad
    \Pi<0\Rightarrow u=0.
\]
Особая дуга:
\[
    \Pi\equiv0\Rightarrow \psi_1=\psi_2,
    \quad
    \dot\psi_1=\dot\psi_2
    \Rightarrow
    F_{x_1}=F_{x_2}.
\]
\[
    \frac{\eps_1}{x_1}F=\frac{\eps_2}{x_2}F
    \Rightarrow
    \frac{x_1}{\eps_1}=\frac{x_2}{\eps_2}.
\]
\[
    \frac{\dot x_1}{\eps_1}=\frac{\dot x_2}{\eps_2}
    \Rightarrow
    \frac{u}{\eps_1}=\frac{1-u}{\eps_2}
    \Rightarrow
    u_{\sng}=\eps_1.
\]
Структура:
\[
    \text{boundary arc }(u=0\text{ or }1)
    \to
    \text{singular arc }(u=\eps_1)
    \to
    \text{terminal arc }(u=0).
\]

\newpage

\section{Билет 6}

\subsection*{Официальная формулировка}

\textbf{6. Принцип максимума Понтрягина для задач на бесконечном интервале времени (общий случай). Основные соотношения принципа максимума. Условия трансверсальности на бесконечности.}

\subsection{Постановка задачи}

Рассмотрим задачу оптимального управления на бесконечном полуинтервале времени:
\[
    \dot x(t)=f(x(t),u(t)),
    \qquad
    u(t)\in U,
    \qquad
    x(0)=x_0,
    \qquad
    t\ge0,
\]
\[
    J(x,u)=\int_0^\infty e^{-\rho t}g(x(t),u(t))\,dt\to\max.
\]
Здесь \(x(t)\in G\subset\R^n\), \(u(t)\in U\subset\R^m\), \(U\) -- компакт, \(\rho\ge0\) -- параметр дисконтирования.

Допустимое управление -- измеримая функция \(u:[0,\infty)\to U\), для которой существует соответствующая траектория \(x(\cdot)\), удовлетворяющая системе. Допустимая пара \((x,u)\) называется оптимальной, если функционал \(J\) принимает на ней наибольшее значение.

Так как горизонт времени бесконечен, функционал является несобственным интегралом. Поэтому обычно предполагают условия, обеспечивающие его сходимость, например равномерную оценку хвоста:
\[
    \int_T^\infty e^{-\rho t}|g(x(t),u(t))|\,dt\le\omega(T),
    \qquad
    \omega(T)\to0,
    \quad T\to\infty.
\]

\subsection{Основные соотношения принципа максимума}

Определим функцию Гамильтона--Понтрягина:
\[
    H(x,t,u,\psi,\psi^0)
    =\langle f(x,u),\psi\rangle+\psi^0e^{-\rho t}g(x,u),
\]
и гамильтониан
\[
    H(x,t,\psi,\psi^0)=\sup_{u\in U}H(x,t,u,\psi,\psi^0).
\]
Здесь \(\psi(t)\in\R^n\) -- сопряженная переменная, \(\psi^0\ge0\) -- числовой множитель.

Пара \((\psi,\psi^0)\) должна быть нетривиальной:
\[
    \|\psi(0)\|+\psi^0>0.
\]

Если \((x^*,u^*)\) -- оптимальная допустимая пара, то существует нетривиальная пара \((\psi,\psi^0)\), такая что выполняются основные соотношения ПМП:
\[
    \dot x^*(t)=f(x^*(t),u^*(t)),
\]
\[
    \dot\psi(t)=-\frac{\partial H}{\partial x}
    (x^*(t),t,u^*(t),\psi(t),\psi^0),
\]
то есть
\[
    \dot\psi(t)
    =-\bigl(f_x(x^*(t),u^*(t))\bigr)^T\psi(t)
    -\psi^0e^{-\rho t}g_x(x^*(t),u^*(t)).
\]
Условие максимума:
\[
    H(x^*(t),t,u^*(t),\psi(t),\psi^0)
    =\max_{u\in U}H(x^*(t),t,u,\psi(t),\psi^0)
\]
для почти всех \(t\ge0\).

Если \(\psi^0>0\), можно перейти к нормальной форме и положить
\[
    \psi^0=1.
\]
Если \(\psi^0=0\), имеет место анормальный случай, который нельзя заранее исключать без дополнительных предположений.

\subsection{Условия трансверсальности на бесконечности}

Для задач на конечном интервале обычно есть конечные условия трансверсальности в правом конце. В задаче на бесконечном интервале таких условий автоматически нет.

В приложениях часто используют условия вида
\[
    \lim_{t\to\infty}\psi(t)=0
\]
или
\[
    \lim_{t\to\infty}\langle x^*(t),\psi(t)\rangle=0.
\]
Однако эти условия не являются следствием основных соотношений ПМП в общем случае. Они требуют дополнительных предположений.

Главное отличие бесконечного горизонта от конечного состоит в том, что нужно отдельно контролировать:
\begin{itemize}
    \item сходимость функционала;
    \item поведение сопряженной переменной на бесконечности;
    \item корректность выбранного условия трансверсальности.
\end{itemize}

\subsection{Краткая версия для запоминания}

\[
    \dot x=f(x,u),
    \qquad
    J=\int_0^\infty e^{-\rho t}g(x,u)\,dt\to\max.
\]
\[
    H=\langle f(x,u),\psi\rangle+\psi^0e^{-\rho t}g(x,u).
\]
ПМП:
\[
    \dot\psi=-H_x,
    \qquad
    H(x^*,u^*,\psi,\psi^0)=\max_{u\in U}H(x^*,u,\psi,\psi^0).
\]
Нетривиальность:
\[
    \|\psi(0)\|+\psi^0>0.
\]
Нормальный случай:
\[
    \psi^0=1.
\]
Главная осторожность:
\[
    \lim_{t\to\infty}\psi(t)=0
\]
не является автоматическим условием в общем случае.

\newpage

\section{Билет 7}

\subsection*{Официальная формулировка}

\textbf{7. Метод динамического программирования Беллмана и принцип максимума Понтрягина для задач на бесконечном интервале времени. Текущие сопряженные переменные. Экономический смысл соотношений принципа максимума.}

\subsection{Метод Беллмана}

Рассмотрим стационарную задачу на бесконечном горизонте:
\[
    \dot x(t)=f(x(t),u(t)),
    \qquad
    u(t)\in U,
    \qquad
    x(0)=x,
\]
\[
    J(x,u)=\int_0^\infty e^{-\rho t}g(x(t),u(t))\,dt\to\max.
\]

В методе динамического программирования вводится функция цены:
\[
    V(x)=\sup_{u(\cdot)}\int_0^\infty e^{-\rho t}g(x(t),u(t))\,dt.
\]
Она показывает максимальный дисконтированный выигрыш, который можно получить, если начальное состояние равно \(x\).

Если функция \(V\) достаточно гладкая, то для стационарной задачи она удовлетворяет уравнению Беллмана:
\[
    \rho V(x)=\max_{u\in U}\{g(x,u)+\langle V_x(x),f(x,u)\rangle\}.
\]
Если существует управление
\[
    u^*(x)\in\argmax_{u\in U}\{g(x,u)+\langle V_x(x),f(x,u)\rangle\},
\]
то при выполнении условий проверочной теоремы Беллмана обратная связь
\[
    u(t)=u^*(x(t))
\]
является оптимальной.

\subsection{Связь с принципом максимума}

Связь с ПМП задается через текущую сопряженную переменную. Пусть
\[
    p(t)=V_x(x^*(t)).
\]
Тогда \(p(t)\) называется текущей сопряженной переменной.

Она связана с обычной, то есть приведенной к начальному моменту, сопряженной переменной \(\psi(t)\) формулами
\[
    \psi(t)=e^{-\rho t}p(t),
    \qquad
    p(t)=e^{\rho t}\psi(t).
\]

Введем гамильтониан в текущих ценах:
\[
    M(x,u,p)=g(x,u)+\langle p,f(x,u)\rangle.
\]
Тогда условия ПМП в текущих переменных имеют вид
\[
    \dot x^*(t)=f(x^*(t),u^*(t)),
\]
\[
    \dot p(t)=\rho p(t)-M_x(x^*(t),u^*(t),p(t)),
\]
то есть
\[
    \dot p(t)=\rho p(t)-\bigl(f_x(x^*(t),u^*(t))\bigr)^Tp(t)-g_x(x^*(t),u^*(t)).
\]
Условие максимума:
\[
    M(x^*(t),u^*(t),p(t))
    =\max_{u\in U}M(x^*(t),u,p(t))
\]
для почти всех \(t\ge0\).

На оптимальной траектории уравнение Беллмана дает
\[
    \rho V(x^*(t))=M(x^*(t),u^*(t),V_x(x^*(t))).
\]

\subsection{Экономический смысл}

Функция \(V(x)\) -- это внутренняя стоимость состояния \(x\), то есть максимальная будущая дисконтированная полезность. Градиент \(V_x(x)\) показывает предельную ценность дополнительной единицы состояния.

В экономике \(p(t)\) интерпретируется как текущая теневая цена состояния, а \(\psi(t)\) -- как та же цена, приведенная к начальному моменту времени.

Условие максимума означает, что оптимальное управление максимизирует сумму текущего выигрыша и прироста ценности состояния:
\[
    g(x,u)+\langle p,f(x,u)\rangle.
\]

В терминах текущей сопряженной переменной условия трансверсальности часто записываются как
\[
    \lim_{t\to\infty}e^{-\rho t}p(t)=0
\]
или
\[
    \lim_{t\to\infty}e^{-\rho t}\langle x^*(t),p(t)\rangle=0.
\]
Но такие условия не являются автоматическими в общем случае и требуют дополнительных предположений.

\subsection{Важное замечание}

Метод Беллмана и ПМП имеют разный логический статус. ПМП обычно дает необходимые условия оптимальности. Проверочная теорема Беллмана при выполнении своих условий дает достаточное условие оптимальности.

\subsection{Краткая версия для запоминания}

\[
    V(x)=\sup_u\int_0^\infty e^{-\rho t}g(x(t),u(t))\,dt.
\]
Уравнение Беллмана:
\[
    \rho V(x)=\max_{u\in U}\{g(x,u)+\langle V_x(x),f(x,u)\rangle\}.
\]
Текущая сопряженная переменная:
\[
    p(t)=V_x(x^*(t)).
\]
Связь переменных:
\[
    \psi(t)=e^{-\rho t}p(t).
\]
Текущий гамильтониан:
\[
    M(x,u,p)=g(x,u)+\langle p,f(x,u)\rangle.
\]
ПМП в current-value форме:
\[
    \dot p=\rho p-M_x,
    \qquad
    M(x^*,u^*,p)=\max_{u\in U}M(x^*,u,p).
\]

\newpage

\section{Билет 8}

\subsection*{Официальная формулировка}

\textbf{8. Достаточные условия оптимальности для задач на бесконечном интервале времени.}

\subsection{Смысл вопроса}

Принцип максимума Понтрягина дает, вообще говоря, только необходимые условия оптимальности. Поэтому выполнение сопряженной системы и условия максимума само по себе не гарантирует оптимальность допустимого процесса. Для достаточности нужны дополнительные условия, обычно связанные с вогнутостью и корректным условием на бесконечности.

\subsection{Постановка}

Рассмотрим задачу
\[
    \dot x(t)=f(x(t),u(t)),
    \qquad
    u(t)\in U,
    \qquad
    x(0)=x_0,
\]
\[
    J(x,u)=\int_0^\infty e^{-\rho t}g(x(t),u(t))\,dt\to\max.
\]
Пусть \((x^*,u^*)\) -- допустимая пара. Предположим, что существует сопряженная переменная \(\psi(t)\), такая что в нормальной форме выполняются соотношения ПМП:
\[
    \dot\psi(t)=-H_x(x^*(t),t,u^*(t),\psi(t)),
\]
\[
    H(x^*(t),t,u^*(t),\psi(t))
    =\max_{u\in U}H(x^*(t),t,u,\psi(t))
\]
для почти всех \(t\ge0\), где
\[
    H(x,t,u,\psi)=\langle f(x,u),\psi\rangle+e^{-\rho t}g(x,u).
\]

Введем максимизированный гамильтониан:
\[
    \mathcal H(x,t,\psi(t))
    =\max_{u\in U}\{\langle f(x,u),\psi(t)\rangle+e^{-\rho t}g(x,u)\}.
\]

\subsection{Достаточное условие оптимальности}

Одно из стандартных достаточных условий имеет следующий вид. Если:
\begin{enumerate}
    \item множество состояний \(G\) выпукло;
    \item для каждого \(t\ge0\) функция
    \[
        x\mapsto \mathcal H(x,t,\psi(t))
    \]
    вогнута на \(G\);
    \item для любой допустимой траектории \(x(\cdot)\) выполнено предельное условие
    \[
        \liminf_{t\to\infty}\langle \psi(t),x(t)-x^*(t)\rangle\ge0;
    \]
\end{enumerate}
то допустимая пара \((x^*,u^*)\) является оптимальной.

\subsection{Идея доказательства}

Из условия максимума имеем
\[
    H(x^*(t),t,u^*(t),\psi(t))=\mathcal H(x^*(t),t,\psi(t)).
\]
Из вогнутости \(\mathcal H\) по \(x\) следует оценка первого порядка:
\[
    \mathcal H(x(t),t,\psi(t))-\mathcal H(x^*(t),t,\psi(t))
    \le
    \langle \mathcal H_x(x^*(t),t,\psi(t)),x(t)-x^*(t)\rangle.
\]
С учетом сопряженного уравнения эта оценка приводит к неравенству
\[
    J_T(x,u)-J_T(x^*,u^*)
    \le
    -\langle \psi(T),x(T)-x^*(T)\rangle,
\]
где
\[
    J_T(x,u)=\int_0^T e^{-\rho t}g(x(t),u(t))\,dt.
\]
Переходя к пределу при \(T\to\infty\) и используя условие на бесконечности, получаем
\[
    J(x,u)\le J(x^*,u^*).
\]
Следовательно, \((x^*,u^*)\) оптимальна.

\subsection{Достаточность через метод Беллмана}

Другой способ получить достаточные условия -- метод Беллмана. Если существует гладкая функция \(V\), удовлетворяющая уравнению
\[
    \rho V(x)=\max_{u\in U}\{g(x,u)+\langle V_x(x),f(x,u)\rangle\},
\]
и управление \(u^*(x)\) достигает максимума в правой части, то при выполнении соответствующего условия на бесконечности обратная связь
\[
    u(t)=u^*(x(t))
\]
является оптимальной.

Таким образом, в задачах на бесконечном интервале достаточность обычно обеспечивается либо проверочной теоремой Беллмана, либо Arrow-type условием: нормальная форма ПМП, вогнутость максимизированного гамильтониана по фазовой переменной и корректное условие трансверсальности на бесконечности.

\subsection{Краткая версия для запоминания}

ПМП сам по себе -- необходимое условие, не достаточное.

Достаточность:
\[
    \text{нормальная форма ПМП}
    +
    \text{вогнутость }\mathcal H(x,t,\psi(t))\text{ по }x
    +
    \text{условие на бесконечности}
    \Rightarrow
    \text{оптимальность}.
\]
Главное предельное условие:
\[
    \liminf_{t\to\infty}\langle \psi(t),x(t)-x^*(t)\rangle\ge0.
\]
Идея доказательства:
\[
    J_T(x,u)-J_T(x^*,u^*)
    \le
    -\langle\psi(T),x(T)-x^*(T)\rangle.
\]
После перехода к пределу:
\[
    J(x,u)\le J(x^*,u^*).
\]

\newpage

\section*{Итоговая краткая схема блока B}
\addcontentsline{toc}{section}{Итоговая краткая схема блока B}

\begin{longtable}{p{0.14\textwidth}p{0.78\textwidth}}
\toprule
\textbf{Билет} & \textbf{Что помнить}\\
\midrule
4 & Модель Рамсея: \(\dot x=uF(x)-\mu x\), \(J=\int_0^\infty e^{-\nu t}(1-u)F(x)dt\to\max\). Особый режим: \(F'(x_s)=\mu+\nu\), \(u_sF(x_s)=\mu x_s\).\\
\midrule
5 & Двухсекторная модель: \(\dot x_1=uF(x)\), \(\dot x_2=(1-u)F(x)\), \(J=x_2(T)\to\max\). Особая магистраль: \(x_1/\eps_1=x_2/\eps_2\), особое управление: \(u_{\sng}=\eps_1\).\\
\midrule
6 & Бесконечный горизонт: ПМП = сопряженная система + условие максимума + нетривиальность. Нельзя без условий писать \(\lim\psi(t)=0\).\\
\midrule
7 & Bellman: \(\rho V=\max\{g+\langle V_x,f\rangle\}\). Current-value variable: \(p=V_x\), \(\psi=e^{-\rho t}p\).\\
\midrule
8 & ПМП не является достаточным условием. Достаточность: вогнутость максимизированного гамильтониана + корректное предельное условие.\\
\bottomrule
\end{longtable}

\blocktitle{C}{Геометрическая теория управления, сопряженные точки и особые траектории}{Билеты 9--11}
\section*{Билет 9}
\addcontentsline{toc}{section}{Билет 9. Геометрическая теория управления}

\subsection*{Официальная формулировка}
\textbf{9. Введение в геометрическую теорию управления: гладкие многообразия, обыкновенные дифференциальные уравнения, векторные поля, коммутаторы, алгебра Ли векторных полей. Условие скобочной порождаемости для систем линейных по управлению.}

\subsection*{Смысл вопроса}
В этом билете надо перейти от записи системы в $\mathbb R^n$ к геометрическому языку. Состояние системы является точкой гладкого многообразия, правая часть дифференциального уравнения является векторным полем, а управляемая система представляет собой семейство векторных полей. Главная идея состоит в том, что коммутаторы векторных полей дают новые инфинитезимальные направления движения, которые могут отсутствовать среди мгновенно допустимых скоростей. Поэтому условие скобочной порождаемости связывает алгебру Ли векторных полей с управляемостью.

\subsection*{Гладкие многообразия и касательные пространства}
Гладкое $n$-мерное многообразие $M$ --- это пространство, которое локально устроено как $\mathbb R^n$ и снабжено согласованными гладкими координатными картами. Точка многообразия обозначается через $q\in M$, а ее координаты в локальной карте --- через $x=(x^1,\ldots,x^n)$.

Касательное пространство $T_qM$ в точке $q\in M$ --- это линейное пространство касательных векторов к гладким кривым, проходящим через $q$. Если
\[
\gamma:(-\varepsilon,\varepsilon)\to M,\qquad \gamma(0)=q,
\]
то вектор скорости
\[
\dot\gamma(0)\in T_qM
\]
является касательным вектором в точке $q$. В локальных координатах элементы $T_qM$ имеют вид
\[
\xi=\sum_{i=1}^n \xi^i\frac{\partial}{\partial x^i}\bigg|_q.
\]
Касательное расслоение обозначается
\[
TM=\bigcup_{q\in M}T_qM.
\]

\subsection*{Векторные поля, ОДУ и потоки}
Гладкое векторное поле на $M$ --- это гладкое отображение
\[
X:M\to TM,\qquad X(q)\in T_qM.
\]
Множество всех гладких векторных полей на $M$ обозначают $\Vect(M)$.

Векторное поле $X$ задает обыкновенное дифференциальное уравнение на многообразии:
\[
\dot q=X(q),\qquad q\in M.
\]
Решением является кривая $q(t)$, такая что
\[
\frac{dq(t)}{dt}=X(q(t)).
\]
Если решение зависит от начальной точки $q_0$, то отображение
\[
P_X^t(q_0)=q(t;q_0)
\]
называется локальным потоком поля $X$. Если поле полно, то поток определен для всех $t\in\mathbb R$ и образует однопараметрическую группу диффеоморфизмов:
\[
P_X^{t+s}=P_X^t\circ P_X^s,
\qquad P_X^0=\operatorname{Id}.
\]

\subsection*{Управляемые системы на многообразиях}
В геометрической теории управления управляемая система рассматривается как семейство векторных полей
\[
\dot q=V_u(q),\qquad q\in M,\quad u\in U,
\]
где $M$ --- пространство состояний, $U$ --- множество управляющих параметров.

Если управление зависит от времени, $u=u(t)$, то получаем неавтономное уравнение
\[
\dot q(t)=V_{u(t)}(q(t)).
\]
Траектории такой системы можно приближать траекториями, соответствующими кусочно-постоянным управлениям, то есть последовательным движениям вдоль потоков различных векторных полей.

Особенно важный класс --- системы, аффинные или линейные по управлению:
\[
\dot q=f_0(q)+\sum_{i=1}^m u_i f_i(q),\qquad u=(u_1,\ldots,u_m).
\]
Если $f_0=0$, система называется бездрейфовой:
\[
\dot q=\sum_{i=1}^m u_i f_i(q).
\]
Для бездрейфовых систем условие скобочной порождаемости имеет прямую связь с управляемостью в смысле теоремы Чоу--Рашевского.

\subsection*{Коммутатор векторных полей}
Векторное поле можно понимать как дифференциальный оператор первого порядка на алгебре гладких функций:
\[
X\varphi(q)=d\varphi(q)X(q),\qquad \varphi\in C^\infty(M).
\]
Коммутатором, или скобкой Ли, двух векторных полей $X,Y\in\Vect(M)$ называется векторное поле
\[
[X,Y](\varphi)=X(Y\varphi)-Y(X\varphi),
\qquad \varphi\in C^\infty(M).
\]
В локальных координатах, если
\[
X=\sum_{i=1}^n X^i(x)\frac{\partial}{\partial x^i},
\qquad
Y=\sum_{i=1}^n Y^i(x)\frac{\partial}{\partial x^i},
\]
то
\[
[X,Y]=\sum_{i=1}^n
\left(
\sum_{j=1}^n
\left(
X^j\frac{\partial Y^i}{\partial x^j}
-
Y^j\frac{\partial X^i}{\partial x^j}
\right)
\right)
\frac{\partial}{\partial x^i}.
\]
Основные свойства скобки Ли:
\[
[X,Y]=-[Y,X],\qquad [X,X]=0,
\]
\[
[X,[Y,Z]]+[Y,[Z,X]]+[Z,[X,Y]]=0.
\]
Последнее равенство называется тождеством Якоби.

\subsection*{Геометрический смысл коммутатора}
Коммутатор показывает некоммутативность потоков. Если за малое время двигаться последовательно вдоль $X$, затем вдоль $Y$, затем назад вдоль $X$ и назад вдоль $Y$, то главный нетривиальный член смещения имеет направление $[X,Y]$:
\[
P_Y^{-\varepsilon}P_X^{-\varepsilon}P_Y^{\varepsilon}P_X^{\varepsilon}(q)
= q+\varepsilon^2[X,Y](q)+o(\varepsilon^2)
\]
в локальных координатах, с возможным изменением знака в зависимости от порядка композиции потоков.

Поэтому даже если мгновенные скорости системы лежат только в подпространстве
\[
\Delta(q)=\operatorname{span}\{f_1(q),\ldots,f_m(q)\}\subset T_qM,
\]
переключения управлений позволяют получать направления, связанные со скобками
\[
[f_i,f_j],\qquad [f_i,[f_j,f_k]],\qquad \ldots
\]
Это ключевая идея неголономной управляемости: ограничения на мгновенные скорости не обязательно означают такие же ограничения на достижимые состояния.

\subsection*{Алгебра Ли векторных полей}
Пусть $\mathcal F=\{f_1,\ldots,f_m\}\subset\Vect(M)$. Алгеброй Ли, порожденной этими векторными полями, называется наименьшее линейное пространство векторных полей, содержащее $f_1,\ldots,f_m$ и замкнутое относительно скобки Ли:
\[
\Lie\{f_1,\ldots,f_m\}.
\]
В каждой точке $q\in M$ можно рассмотреть подпространство
\[
\Lie_q\{f_1,\ldots,f_m\}
=
\operatorname{span}\{X(q):X\in\Lie\{f_1,\ldots,f_m\}\}\subset T_qM.
\]
Оно порождается значениями исходных полей и всех повторных коммутаторов:
\[
f_i(q),\quad [f_i,f_j](q),\quad [f_i,[f_j,f_k]](q),\quad \ldots
\]

\subsection*{Условие скобочной порождаемости}
Для бездрейфовой системы
\[
\dot q=\sum_{i=1}^m u_i f_i(q)
\]
условие скобочной порождаемости имеет вид
\[
\Lie_q\{f_1,\ldots,f_m\}=T_qM
\qquad \text{для всех }q\in M.
\]
Иначе говоря, исходные поля и их повторные коммутаторы должны порождать все касательное пространство в каждой точке.

Если $\dim M=n$, то это условие эквивалентно ранговому условию
\[
\rank\{f_i(q),[f_i,f_j](q),[f_i,[f_j,f_k]](q),\ldots\}=n.
\]

В случае системы с дрейфом
\[
\dot q=f_0(q)+\sum_{i=1}^m u_i f_i(q)
\]
аналогичное условие
\[
\Lie_q\{f_0,f_1,\ldots,f_m\}=T_qM
\]
обычно является условием локальной достижимости или доступности, но не всегда полной управляемости. Наличие дрейфа $f_0$ делает задачу тоньше: из ранга алгебры Ли не всегда следует возможность двигаться назад по времени вдоль дрейфа.

\subsection*{Теорема Чоу--Рашевского и управляемость}
Стандартная формулировка для бездрейфовой системы такова. Пусть $M$ --- связное гладкое многообразие, а система имеет вид
\[
\dot q=\sum_{i=1}^m u_i f_i(q),
\]
где управления позволяют двигаться в обе стороны вдоль полей $f_i$. Если для всех $q\in M$
\[
\Lie_q\{f_1,\ldots,f_m\}=T_qM,
\]
то любые две точки одной связной компоненты можно соединить допустимой траекторией системы. В локальной форме это означает, что множество достижимости имеет полный размер.

Экзаменационный смысл: геометрическая теория управления переводит вопрос об управляемости в вопрос о ранге системы векторных полей и их скобок Ли. Если мгновенных направлений мало, но их скобки порождают всё касательное пространство, то система может быть управляема.

\subsection*{Пример: модель автомобиля}
Простейшая модель автомобиля на $M=\mathbb R^2\times S^1$ задается векторными полями
\[
f_1(q)=
\begin{pmatrix}
\cos\theta\\
\sin\theta\\
0
\end{pmatrix},
\qquad
f_2(q)=
\begin{pmatrix}
0\\
0\\
1
\end{pmatrix},
\qquad q=(x_1,x_2,\theta).
\]
Здесь $f_1$ соответствует движению вперед-назад, а $f_2$ --- вращению автомобиля. Их коммутатор равен
\[
[f_1,f_2](q)=
\begin{pmatrix}
\sin\theta\\
-\cos\theta\\
0
\end{pmatrix}
\]
с точностью до знака в зависимости от принятого порядка скобки. Векторы
\[
f_1(q),\quad f_2(q),\quad [f_1,f_2](q)
\]
порождают всё пространство $T_qM$ размерности 3. Следовательно, хотя автомобиль не может двигаться мгновенно вбок, боковое направление возникает через маневр, соответствующий коммутатору.

\subsection*{Краткая схема ответа}
\begin{enumerate}
    \item Геометрическая теория управления рассматривает системы на гладких многообразиях.
    \item Векторное поле $X$ задает ОДУ $\dot q=X(q)$ и поток $P_X^t$.
    \item Управляемая система --- семейство векторных полей: $\dot q=V_u(q)$.
    \item Для системы, линейной по управлению: $\dot q=f_0(q)+\sum u_i f_i(q)$.
    \item Скобка Ли: $[X,Y](\varphi)=X(Y\varphi)-Y(X\varphi)$.
    \item Коммутаторы дают направления, достижимые за счет быстрых переключений управлений.
    \item Условие скобочной порождаемости: $\Lie_q\{f_1,\ldots,f_m\}=T_qM$.
    \item Для бездрейфовых систем это ведет к управляемости по теореме Чоу--Рашевского; для систем с дрейфом обычно дает доступность, а не автоматически полную управляемость.
\end{enumerate}

\newpage
\section*{Билет 10}
\addcontentsline{toc}{section}{Билет 10. Сопряженные точки и слабый минимум}

\subsection*{Официальная формулировка}
\textbf{10. Понятие сопряженной точки. Достаточные условия слабого минимума в классической задаче вариационного исчисления.}

\subsection*{Смысл вопроса}
Вопрос требует раскрыть два понятия: что такое сопряженная точка и какие условия обеспечивают слабый минимум в классической задаче вариационного исчисления. Основная логика:
\[
\text{уравнение Эйлера}
\Rightarrow
\text{вторая вариация}
\Rightarrow
\text{условие Лежандра}
\Rightarrow
\text{уравнение Якоби}
\Rightarrow
\text{сопряженные точки}
\Rightarrow
\text{слабый минимум}.
\]

\subsection*{Постановка задачи}
Рассмотрим простейшую классическую задачу вариационного исчисления:
\[
J[x]=\int_{t_0}^{t_1} f(t,x(t),\dot x(t))\,dt\to\min,
\]
\[
x(t_0)=a,\qquad x(t_1)=b,
\]
где
\[
x(t)\in\mathbb R^n,\qquad f\in C^2.
\]

Кривая $\widehat x(\cdot)$ называется слабым локальным минимумом, если существует $\varepsilon>0$ такое, что для любой допустимой кривой $x(\cdot)$, удовлетворяющей тем же граничным условиям и условию
\[
\|x-\widehat x\|_{C^1}<\varepsilon,
\]
выполнено
\[
J[x]\ge J[\widehat x].
\]
Иными словами, слабый минимум --- это минимум относительно малых $C^1$-возмущений.

\subsection*{Уравнение Эйлера--Лагранжа}
Первое необходимое условие экстремума --- уравнение Эйлера--Лагранжа:
\[
\frac{d}{dt}f_{\dot x}(t,\widehat x,\dot{\widehat x})
-
f_x(t,\widehat x,\dot{\widehat x})=0.
\]
Кривая, удовлетворяющая этому уравнению и граничным условиям, называется экстремалью.

\subsection*{Вторая вариация}
Для исследования минимума рассматриваются вариации
\[
x_\lambda(t)=\widehat x(t)+\lambda h(t),
\]
где
\[
h(t_0)=h(t_1)=0.
\]
На экстремали первая вариация равна нулю. Главную роль играет вторая вариация:
\[
\delta^2J[h]
=
\int_{t_0}^{t_1}
\left(
\langle A(t)\dot h,\dot h\rangle
+
2\langle C(t)\dot h,h\rangle
+
\langle B(t)h,h\rangle
\right)dt,
\]
где все матрицы вычисляются на экстремали $(t,\widehat x(t),\dot{\widehat x}(t))$:
\[
A(t)=f_{\dot x\dot x}(t,\widehat x,\dot{\widehat x}),
\]
\[
B(t)=f_{xx}(t,\widehat x,\dot{\widehat x}),
\]
\[
C(t)=f_{x\dot x}(t,\widehat x,\dot{\widehat x}).
\]
Матрицы $A(t)$ и $B(t)$ симметричны, а $C(t)$ вообще говоря не обязана быть симметричной.

\subsection*{Условие Лежандра}
Необходимое условие Лежандра для минимума имеет вид
\[
A(t)=f_{\dot x\dot x}(t,\widehat x,\dot{\widehat x})\ge0.
\]
Усиленное условие Лежандра:
\[
A(t)>0,
\]
то есть существует $\alpha>0$ такое, что
\[
\langle A(t)\xi,\xi\rangle\ge \alpha|\xi|^2
\]
для всех $\xi\in\mathbb R^n$ и всех $t\in[t_0,t_1]$.

\subsection*{Уравнение Якоби и сопряженные точки}
Рассмотрим присоединенную задачу для второй вариации. Ее уравнение Эйлера имеет вид
\[
-\frac{d}{dt}\left(A(t)\dot h(t)+C^T(t)h(t)\right)
+
C(t)\dot h(t)+B(t)h(t)=0.
\]
Это уравнение называется уравнением Якоби для исходной вариационной задачи.

Точка $\tau\in(t_0,t_1]$ называется сопряженной с точкой $t_0$ вдоль экстремали $\widehat x(\cdot)$, если существует нетривиальное решение $h(t)\not\equiv0$ уравнения Якоби, такое что
\[
h(t_0)=0,\qquad h(\tau)=0.
\]
Иначе говоря, сопряженная точка --- это момент, в котором возникает нетривиальная вариация, обращающаяся в нуль в двух концах и удовлетворяющая уравнению Якоби.

Геометрически наличие сопряженной точки означает, что семейство близких экстремалей, выходящих из одной начальной точки, начинает иметь пересечение или фокусировку. Аналитически появление сопряженной точки связано с потерей положительной определенности второй вариации.

Условие Якоби для слабого минимума состоит в отсутствии точек, сопряженных с $t_0$, на интервале
\[
(t_0,t_1].
\]

\subsection*{Достаточное условие слабого минимума}
Пусть $\widehat x(\cdot)$ --- экстремаль, удовлетворяющая уравнению Эйлера--Лагранжа и заданным граничным условиям. Если:
\begin{enumerate}
    \item выполнено усиленное условие Лежандра:
    \[
    A(t)=f_{\dot x\dot x}(t,\widehat x,\dot{\widehat x})>0;
    \]
    \item на интервале $(t_0,t_1]$ нет точек, сопряженных с $t_0$,
\end{enumerate}
то вторая вариация положительно определена:
\[
\delta^2J[h]>0
\]
для всех ненулевых допустимых вариаций $h$, $h(t_0)=h(t_1)=0$. Следовательно, экстремаль $\widehat x(\cdot)$ доставляет строгий слабый локальный минимум функционала $J$.

Таким образом, для слабого минимума недостаточно только уравнения Эйлера. Нужно проверить вторые условия: усиленное условие Лежандра и отсутствие сопряженных точек. Эти условия обеспечивают положительную определенность второй вариации, а значит, малые $C^1$-отклонения от экстремали увеличивают значение функционала.

\subsection*{Краткая схема ответа}
\[
\text{Euler}+
\text{усиленное условие Лежандра}+
\text{нет сопряженных точек}
\Rightarrow
\text{строгий слабый минимум}.
\]

\newpage
\section*{Билет 11}
\addcontentsline{toc}{section}{Билет 11. Особые траектории и особые управления}

\subsection*{Официальная формулировка}
\textbf{11. Особые траектории. Метод построения особых управлений.}

\subsection*{Смысл вопроса}
Вопрос требует объяснить, что такое особое управление и особая траектория, почему в особом режиме условие максимума Понтрягина не определяет управление, и как управление строится из условий
\[
H_u=0,\qquad \frac{d}{dt}H_u=0,\qquad \frac{d^2}{dt^2}H_u=0,\ldots
\]
до появления управления в явном виде.

\subsection*{Система, линейная по управлению}
Рассмотрим задачу оптимального управления, в которой управление входит в систему линейно. Для простоты сначала возьмем скалярное управление:
\[
\dot x(t)=f_0(t,x(t))+u(t)f_1(t,x(t)),
\qquad u(t)\in [a,b].
\]
В автономном случае:
\[
\dot x=f_0(x)+u f_1(x).
\]
Пусть Hamiltonian имеет вид
\[
H(x,\psi,u)=H_0(x,\psi)+uH_1(x,\psi),
\]
где
\[
H_0(x,\psi)=\langle \psi,f_0(x)\rangle,
\qquad
H_1(x,\psi)=\langle \psi,f_1(x)\rangle.
\]
Если имеется интегральный функционал, линейный по управлению, то соответствующий член также включается в $H_0$ и $H_1$. Главное свойство сохраняется: Hamiltonian является линейной функцией управления $u$.

\subsection*{Функция переключения}
Функция
\[
\Phi(t)=H_u(x(t),\psi(t),u(t))
\]
называется функцией переключения. В случае
\[
H=H_0+uH_1
\]
имеем
\[
\Phi(t)=H_1(x(t),\psi(t)).
\]

По принципу максимума Понтрягина оптимальное управление должно максимизировать Hamiltonian по $u\in[a,b]$:
\[
H(x(t),\psi(t),u(t))=
\max_{v\in[a,b]}H(x(t),\psi(t),v).
\]
Если
\[
\Phi(t)>0,
\]
то максимум достигается при
\[
u(t)=b.
\]
Если
\[
\Phi(t)<0,
\]
то максимум достигается при
\[
u(t)=a.
\]
Это обычные bang-bang участки.

\subsection*{Особый режим}
Особый режим возникает, если на некотором интервале $I=[t_1,t_2]$ выполнено тождество
\[
\Phi(t)\equiv0,
\qquad t\in I.
\]
В этом случае условие максимума первого порядка не определяет управление, потому что Hamiltonian не зависит от $u$ в первом порядке. Соответствующее управление называется особым управлением, а соответствующая траектория $x(t)$ называется особой траекторией.

Итак, особая траектория --- это проекция на фазовое пространство такой экстремали принципа максимума, на которой условие максимума не задает управление однозначно на некотором интервале.

\subsection*{Метод построения особого управления}
Метод состоит в последовательном дифференцировании функции переключения. На особом интервале должно выполняться
\[
\Phi(t)=0.
\]
Так как это тождество по времени, то также
\[
\dot\Phi(t)=0,\qquad \ddot\Phi(t)=0,
\]
и так далее. Дифференцирование производится вдоль экстремальной системы:
\[
\dot x=H_\psi,
\qquad
\dot\psi=-H_x.
\]

Если управление $u$ впервые явно появляется в производной порядка $k$, то из условия
\[
\left(\frac{d}{dt}\right)^k\Phi(t)=0
\]
можно выразить особое управление:
\[
u_{\sng}(t)=\cdots.
\]

В типичной теории особых экстремалей для задач, линейных по управлению, первый порядок, в котором управление появляется существенно, является четным. Если управление впервые появляется в производной порядка $2q$, то число $q$ называется порядком особого управления или особой дуги.

Формально это записывают так:
\[
\frac{\partial}{\partial u}
\left(\frac{d}{dt}\right)^k H_u=0,
\qquad k=0,1,\ldots,2q-1,
\]
но
\[
\frac{\partial}{\partial u}
\left(\frac{d}{dt}\right)^{2q}H_u\ne0.
\]
Тогда особое управление находится из уравнения
\[
\left(\frac{d}{dt}\right)^{2q}H_u=0.
\]
После нахождения $u_{\sng}$ необходимо проверить:
\begin{enumerate}
    \item допустимость: $u_{\sng}(t)\in[a,b]$;
    \item выполнение исходных условий PMP;
    \item выполнение дополнительных необходимых условий второго порядка, например условия Kelley или generalized Legendre--Clebsch.
\end{enumerate}

\subsection*{Условие Kelley}
При максимизационной форме принципа максимума часто используется условие Kelley:
\[
(-1)^q
\frac{\partial}{\partial u}
\left(\frac{d}{dt}\right)^{2q}H_u
<0
\]
на особой дуге. Это условие является необходимым условием локальной оптимальности особой дуги. Однако знак этого условия зависит от соглашения о Hamiltonian и от того, рассматривается задача максимизации или минимизации. Поэтому при использовании другой конвенции знак может измениться.

\subsection*{Запись через скобки Пуассона}
В автономном случае удобно использовать скобки Пуассона. Пусть
\[
h_0(x,\psi)=\langle \psi,f_0(x)\rangle,
\qquad
h_1(x,\psi)=\langle \psi,f_1(x)\rangle.
\]
Тогда
\[
H=h_0+u h_1,
\]
а switching function равна
\[
\Phi=h_1.
\]
Производная функции $h$ вдоль гамильтоновой системы выражается через скобку Пуассона:
\[
\dot h=\{H,h\}.
\]
Поэтому
\[
\dot\Phi=\{H,h_1\}.
\]
Так как
\[
\{h_1,h_1\}=0,
\]
получаем
\[
\dot\Phi=\{h_0,h_1\}.
\]
Следующая производная:
\[
\ddot\Phi
=
\{H,\{h_0,h_1\}\}
=
\{h_0,\{h_0,h_1\}\}
+
u\{h_1,\{h_0,h_1\}\}.
\]
Если
\[
\{h_1,\{h_0,h_1\}\}\ne0,
\]
то особое управление первого порядка находится из условия $\ddot\Phi=0$:
\[
u_{\sng}
=
-
\frac{\{h_0,\{h_0,h_1\}\}}
{\{h_1,\{h_0,h_1\}\}}.
\]
После этого нужно проверить
\[
u_{\sng}\in[a,b]
\]
и вторые необходимые условия оптимальности.

\subsection*{Векторное управление}
Если система имеет вид
\[
\dot x=f_0(x)+\sum_{i=1}^m u_i f_i(x),
\]
то
\[
H=h_0+\sum_{i=1}^m u_i h_i,
\qquad
h_i(x,\psi)=\langle \psi,f_i(x)\rangle.
\]
Функции $h_i(t)$ являются функциями переключения. Если на некотором интервале
\[
h_i(t)\equiv0
\]
для одной или нескольких компонент управления, то соответствующие компоненты являются особыми. Для их построения последовательно дифференцируют тождества $h_i=0$ до тех пор, пока соответствующие компоненты управления не появятся явно.

\subsection*{Chattering / Fuller phenomenon}
В окрестности особых дуг могут возникать участки с очень частыми переключениями bang-bang управления. В некоторых задачах появляется бесконечное число переключений на конечном интервале времени; это явление называют chattering или явлением Фуллера. Для основного экзаменационного ответа достаточно знать, что оно связано с сопряжением bang-bang участков с особыми дугами и с условиями высокого порядка.

\subsection*{Итог}
Особое управление не угадывается и не задается произвольно. Оно выводится из тождества функции переключения:
\[
H_u=0
\]
и его последовательных производных:
\[
H_u=0,
\quad \frac{d}{dt}H_u=0,
\quad
\frac{d^2}{dt^2}H_u=0,
\quad \ldots
\]
до появления управления. Затем найденное управление проверяется на допустимость и на выполнение дополнительных необходимых условий оптимальности.

\subsection*{Краткая схема ответа}
\[
\Phi\equiv0
\Rightarrow
\Phi=\dot\Phi=\ddot\Phi=\cdots=0
\Rightarrow
u_{\sng}
\Rightarrow
\text{проверка допустимости и Kelley}.
\]

\blocktitle{D}{Граничное управление, наблюдение и вариационный метод}{Билеты 12--14}
\section*{Замечание о надежности}
Настоящий текст является экзаменационной версией, составленной по официальным формулировкам билетов, студенческим конспектам 2017 и 2026 гг., материалам по функциональному анализу и методам оптимизации, а также по общим фактам теории управляемости волновых уравнений. Полный текст книги \textit{Васильев Ф.П., Куржанский М.А., Потапов М.М., Разгулин А.В. Приближенное решение двойственных задач управления и наблюдения} пока не использован, поэтому численные константы в неравенствах наблюдаемости и отдельные детали конечномерной аппроксимации желательно сверить по этому источнику.

\section*{Билет 12}
\addcontentsline{toc}{section}{Билет 12}

\subsection*{Официальная формулировка}
\textbf{12. Постановка задач граничного управления и наблюдения для волнового уравнения. Выбор функциональных пространств. Двойственность. Свойства управляемости и наблюдаемости и их связь с критическим моментом времени.}

\subsection*{Ответ}
Рассмотрим волновое уравнение на отрезке $(0,l)$:
\[
 y_{tt}(t,x)=y_{xx}(t,x)-q(x)y(t,x),
 \qquad 0<t<T,\quad 0<x<l.
\]
Задача граничного управления состоит в выборе граничных функций
\[
 y(t,0)=u_0(t),\qquad y(t,l)=u_1(t),
\]
так, чтобы решение, выходя из нулевого начального состояния
\[
 y(0,x)=0,\qquad y_t(0,x)=0,
\]
достигло в момент $T$ заданного конечного состояния
\[
 y(T,x)=f_0(x),\qquad y_t(T,x)=f_1(x).
\]
Обозначим
\[
 u=(u_0,u_1).
\]
Естественное пространство граничных управлений:
\[
 H=L_2(0,T)\times L_2(0,T).
\]

Так как управление задано на границе, решение имеет более слабую регулярность, чем при распределенном управлении. В слабой постановке удобно брать пространство конечных состояний
\[
 F=L_2(0,l)\times H^{-1}(0,l),
\]
а сопряженное пространство
\[
 F^*=H_0^1(0,l)\times L_2(0,l).
\]
Тогда оператор управления задается формулой
\[
 A:H\to F,\qquad Au=(y(T,\cdot),y_t(T,\cdot)).
\]
Задача точного управления принимает операторный вид
\[
 Au=f,\qquad f=(f_0,f_1)\in F.
\]
Система называется вполне управляемой в момент $T$, если
\[
 R(A)=F,
\]
то есть для любого конечного состояния $f\in F$ существует управление $u\in H$, переводящее систему в это состояние.

Рассмотрим теперь сопряженную задачу наблюдения. Пусть
\[
 v=(v_0,v_1)\in F^*=H_0^1(0,l)\times L_2(0,l).
\]
Сопряженная волновая задача имеет вид
\[
 p_{tt}(t,x)=p_{xx}(t,x)-q(x)p(t,x),
\]
\[
 p(t,0)=p(t,l)=0,
\]
\[
 p(T,x)=v_0(x),\qquad p_t(T,x)=-v_1(x).
\]
Наблюдаемыми сигналами являются граничные производные:
\[
 g_0(t)=p_x(t,0),\qquad g_1(t)=-p_x(t,l).
\]
Иными словами,
\[
 A^*v=(p_x(\cdot,0),-p_x(\cdot,l)).
\]
Здесь
\[
 A^*:F^*\to H^*
\]
- оператор, сопряженный к оператору управления $A$.

Связь управления и наблюдения выражается тождеством двойственности:
\[
 \langle Au,v\rangle_{F,F^*}=\langle u,A^*v\rangle_H.
\]
В развернутом виде это означает
\[
 \left\langle (y(T),y_t(T)),(v_0,v_1)\right\rangle_{F,F^*}
 =
 \int_0^T\left[u_0(t)p_x(t,0)-u_1(t)p_x(t,l)\right]dt.
\]
Это тождество показывает, что задача управления и задача наблюдения являются двойственными: управление описывает достижение конечного состояния, а наблюдение описывает восстановление сопряженного состояния по граничному сигналу.

Полная наблюдаемость сопряженной системы означает, что по сигналу
\[
 A^*v=(p_x(\cdot,0),-p_x(\cdot,l))
\]
можно устойчиво восстановить данные $v$. В сильной форме это выражается неравенством наблюдаемости:
\[
 \|A^*v\|_{H^*}^2\ge \mu \|v\|_{F^*}^2,
 \qquad \mu>0.
\]
Это неравенство означает, что ненулевое сопряженное состояние не может иметь нулевой или слишком малый наблюдаемый сигнал.

Для волнового уравнения принципиальную роль играет конечная скорость распространения волн. Если время наблюдения или управления слишком мало, волна не успевает распространиться по всему отрезку, и точная управляемость невозможна. Для отрезка длины $l$ и двустороннего граничного управления критическое время обычно равно
\[
 T^*=l.
\]
Если управление или наблюдение проводится только с одного конца, критическое время обычно равно
\[
 T^*=2l.
\]
При
\[
 T>T^*
\]
выполняется неравенство наблюдаемости, и исходная система точно управляема. При
\[
 T<T^*
\]
точная управляемость и полная наблюдаемость нарушаются, поскольку часть энергии волны не успевает достичь наблюдаемой границы.

Основная схема билета:
\[
 \text{точная управляемость } A
 \quad\Longleftrightarrow\quad
 \text{наблюдаемость } A^*
 \quad\Longleftrightarrow\quad
 \text{неравенство наблюдаемости}.
\]

\subsection*{Что важно помнить}
\begin{enumerate}
\item Оператор управления: $A:H\to F$, $Au=(y(T),y_t(T))$.
\item Сопряженный оператор наблюдения: $A^*:F^*\to H^*$.
\item Главное тождество: $\langle Au,v\rangle=\langle u,A^*v\rangle$.
\item Управляемость и наблюдаемость связаны через неравенство наблюдаемости.
\item Критическое время возникает из-за конечной скорости распространения волн.
\end{enumerate}

\newpage
\section*{Билет 13}
\addcontentsline{toc}{section}{Билет 13}

\subsection*{Официальная формулировка}
\textbf{13. Неравенство наблюдаемости как критерий управляемости и наблюдаемости. Истокопредставимость нормальных решений задач управления и наблюдения при наличии неравенства наблюдаемости. Описание вариационного метода решения линейных операторных уравнений: условия применимости, алгоритм, сходимость.}

\subsection*{Ответ}
Пусть $H$ и $F$ - гильбертовы пространства, а
\[
 A:H\to F
\]
- линейный ограниченный оператор. Задача управления записывается как операторное уравнение
\[
 Au=f,\qquad f\in F,
\]
где $u\in H$ - искомое управление.

Сопряженный оператор имеет вид
\[
 A^*:F^*\to H^*.
\]
Задача наблюдения состоит в восстановлении элемента $v\in F^*$ по наблюдаемому сигналу
\[
 A^*v\in H^*.
\]

Система называется вполне управляемой, если
\[
 R(A)=F,
\]
то есть уравнение $Au=f$ разрешимо для любого $f\in F$.

Критерием точной управляемости является неравенство наблюдаемости:
\[
 \|A^*v\|_{H^*}^2\ge \mu\|v\|_{F^*}^2
 \qquad \forall v\in F^*,\quad \mu>0.
\]
Это неравенство означает, что ненулевое состояние сопряженной системы не может быть невидимым при наблюдении.

В гильбертовых пространствах это связано с теоремой о замкнутом образе. Если
\[
 \|A^*v\|_{H^*}\ge c\|v\|_{F^*},
\]
то $A^*$ инъективен и его образ замкнут. Это эквивалентно сюръективности оператора $A$, то есть полной управляемости:
\[
 R(A)=F.
\]
Если же известно только
\[
 N(A^*)=\{0\},
\]
то это обычно дает лишь плотность образа $R(A)$, то есть приближенную управляемость, но не обязательно точную управляемость. Для точной управляемости нужна именно оценка снизу, то есть неравенство наблюдаемости.

Если уравнение
\[
 Au=f
\]
разрешимо, то среди всех его решений выделяют нормальное решение:
\[
 u_*=
 \operatorname*{arg\,min}\{\|u\|_H: Au=f\}.
\]
Нормальное решение имеет минимальную норму в пространстве управлений.

При выполнении неравенства наблюдаемости нормальное решение является истокопредставимым. Это означает, что существует элемент $v_*\in F^*$ такой, что
\[
 u_*=J_HA^*v_*,
\]
где
\[
 J_H:H^*\to H
\]
- изоморфизм Рисса.

Элемент $v_*$ можно найти из вариационной задачи. Рассмотрим функционал
\[
 I(v)=\|A^*v\|_{H^*}^2-2\langle f,v\rangle_{F,F^*},
 \qquad v\in F^*.
\]
Если известно априорное ограничение
\[
 \|v_*\|_{F^*}\le r,
\]
то минимизацию проводят на шаре
\[
 B_r=\{v\in F^*: \|v\|_{F^*}\le r\}.
\]
Искомый элемент $v_*$ является решением задачи
\[
 I(v)\to\min,\qquad v\in B_r.
\]
После его нахождения нормальное управление восстанавливается по формуле
\[
 u_*=J_HA^*v_*.
\]

Смысл метода: вместо непосредственного решения уравнения $Au=f$ ищется сопряженная переменная $v$, для которой нормальное решение имеет источник
\[
 u_*=J_HA^*v.
\]

Если данные и оператор заданы приближенно, используют вариационный метод с конечномерными аппроксимациями. Пусть построены пространства
\[
 H_N\subset H,
 \qquad
 F_N^*\subset F^*,
\]
и операторы
\[
 A_N:H_N\to F_N,
 \qquad
 A_N^*:F_N^*\to H_N^*,
\]
так, чтобы сохранялась взаимная сопряженность:
\[
 \langle A_Nu_N,v_N\rangle=\langle u_N,A_N^*v_N\rangle.
\]
Это ключевое требование: нельзя независимо дискретизировать управление и наблюдение, иначе нарушается двойственность.

Конечномерная вариационная задача имеет вид
\[
 I_N(v_N)=\|A_N^*v_N\|_{H_N^*}^2-2\langle f_N,v_N\rangle
 \to\min,
 \qquad v_N\in B_r^N.
\]
После решения этой квадратичной задачи управление восстанавливается:
\[
 u_N=J_{H_N}A_N^*v_N.
\]

Условия применимости метода:
\begin{enumerate}
\item операторное уравнение $Au=f$ разрешимо;
\item нормальное решение истокопредставимо: $u_*=J_HA^*v_*$;
\item выполнено неравенство наблюдаемости: $\|A^*v\|_{H^*}^2\ge \mu\|v\|_{F^*}^2$;
\item построены согласованные взаимно сопряженные аппроксимации $A_N,A_N^*$;
\item приближенные правые части сходятся: $f_N\to f$;
\item точность минимизации стремится к нулю.
\end{enumerate}
При выполнении этих условий приближенные управления сходятся к нормальному решению:
\[
 \|u_N-u_*\|_H\to0.
\]

\subsection*{Что важно помнить}
\begin{enumerate}
\item $A:H\to F$, $A^*:F^*\to H^*$.
\item Точная управляемость требует оценки снизу $\|A^*v\|\ge c\|v\|$.
\item $N(A^*)=\{0\}$ само по себе дает только приближенную управляемость, а не точную.
\item Нормальное решение имеет вид $u_*=J_HA^*v_*$.
\item $v_*$ ищется из квадратичного функционала $I(v)=\|A^*v\|^2-2\langle f,v\rangle$.
\end{enumerate}

\newpage
\section*{Билет 14}
\addcontentsline{toc}{section}{Билет 14}

\subsection*{Официальная формулировка}
\textbf{14. Вывод неравенства наблюдаемости для простейшего уравнения колебаний однородной струны. Реализация вариационного метода: конструкции конечномерных взаимносопряжённых приближений к операторам управления и наблюдения; схема решения конечномерной задачи квадратичной минимизации.}

\subsection*{Ответ}
Рассмотрим простейшее уравнение колебаний однородной струны:
\[
 p_{tt}(t,x)=p_{xx}(t,x),
 \qquad 0<t<T,\quad 0<x<l,
\]
с однородными граничными условиями
\[
 p(t,0)=p(t,l)=0.
\]
Это сопряженная задача наблюдения для гранично управляемого волнового уравнения. Терминальные данные задаются в момент $T$:
\[
 p(T,x)=v_0(x),
 \qquad
 p_t(T,x)=-v_1(x).
\]
Энергия решения:
\[
 E(t)=\int_0^l\left(p_t^2(t,x)+p_x^2(t,x)\right)dx.
\]
Для однородного волнового уравнения энергия сохраняется:
\[
 E(t)=E(T)=E(0).
\]
Наблюдаемыми величинами являются граничные производные:
\[
 p_x(t,0),\qquad -p_x(t,l).
\]
Поэтому оператор наблюдения имеет вид
\[
 A^*v=(p_x(\cdot,0),-p_x(\cdot,l)).
\]

Для однородной струны используется энергетическое тождество или метод множителей. Идея состоит в том, что энергия внутри отрезка выражается через поток энергии через границу. Поскольку волны распространяются с конечной скоростью, за время
\[
 T>l
\]
при двустороннем наблюдении вся энергия успевает проявиться на концах струны.

В результате получается неравенство наблюдаемости:
\[
 \int_0^T\left(|p_x(t,0)|^2+|p_x(t,l)|^2\right)dt
 \ge
 \mu\left(\|v_0\|_{H_0^1(0,l)}^2+\|v_1\|_{L_2(0,l)}^2\right),
\]
где $\mu>0$.

В нормировке, используемой в студенческих материалах, для случая $q(x)=0$ указана оценка
\[
 \|A^*v\|^2\ge \mu\|v\|^2,
 \qquad
 \mu=\frac{2(T-l)}{l}>0,
 \qquad T>l.
\]
Точная числовая константа зависит от выбранной нормы энергии и нормировки оператора наблюдения, поэтому в окончательной версии ее желательно сверить с основной литературой 10. Существенно, что положительная константа существует именно при $T>l$.

Это неравенство означает, что энергия любой ненулевой сопряженной волны обязательно наблюдается на границе. Следовательно, исходная задача граничного управления является точно управляемой при $T>l$.

Если же
\[
 T<l,
\]
то существует часть волны, которая не успевает достичь наблюдаемых концов за время $T$. Поэтому неравенство наблюдаемости нарушается и точная управляемость невозможна.

\subsection*{Конечномерные взаимно сопряженные приближения}
Для численного решения операторного уравнения
\[
 Au=f
\]
строят конечномерные аппроксимации оператора управления и оператора наблюдения.

Пусть
\[
 H_N\subset H
\]
- конечномерное пространство управлений, а
\[
 F_N^*\subset F^*
\]
- конечномерное пространство сопряженных состояний. Нужно построить дискретные операторы
\[
 A_N:H_N\to F_N,
 \qquad
 A_N^*:F_N^*\to H_N^*,
\]
так, чтобы выполнялось точное дискретное тождество двойственности:
\[
 \langle A_Nu_N,v_N\rangle_{F_N,F_N^*}
 =
 \langle u_N,A_N^*v_N\rangle_{H_N,H_N^*}.
\]
Это и называется построением взаимно сопряженных приближений. Главное требование: оператор наблюдения $A_N^*$ должен быть действительно сопряженным к $A_N$ в выбранных конечномерных скалярных произведениях. Его нельзя выбирать независимо.

В матричной форме, если скалярные произведения заданы матрицами Грама $G_H$ и $G_F$, сопряженный оператор записывается как
\[
 A_N^*=G_H^{-1}A_N^T G_F.
\]
Если базисы ортонормированы, эта формула упрощается до обычного транспонирования.

\subsection*{Конечномерная квадратичная задача}
После построения взаимно сопряженных операторов вариационный метод сводится к минимизации квадратичного функционала:
\[
 I_N(v_N)=\|A_N^*v_N\|_{H_N^*}^2
 -2\langle f_N,v_N\rangle_{F_N,F_N^*}
 \to\min.
\]
Если известна априорная оценка
\[
 \|v_*\|\le r,
\]
то минимизация ведется на шаре:
\[
 \|v_N\|\le r.
\]
В координатах
\[
 v_N=\sum_{j=1}^m c_j\varphi_j
\]
функционал принимает вид
\[
 I_N(c)=c^TG_Nc-2b_N^Tc,
\]
где $G_N$ - симметричная неотрицательно определенная матрица, соответствующая оператору $A_NA_N^*$, а $b_N$ - вектор правой части, соответствующий $f_N$.

Если $G_N$ положительно определена, минимум находится из системы линейных уравнений
\[
 G_Nc=b_N.
\]
Если есть ограничение $\|v_N\|\le r$, то решается конечномерная задача квадратичного программирования на шаре.

После нахождения $v_N$ управление восстанавливается по формуле
\[
 u_N=J_{H_N}A_N^*v_N.
\]
При согласованности аппроксимаций, выполнении неравенства наблюдаемости, сходимости данных $f_N\to f$ и стремлении ошибки минимизации к нулю получаем
\[
 \|u_N-u_*\|_H\to0.
\]

\subsection*{Что важно помнить}
\begin{enumerate}
\item Для однородной струны наблюдение идет через $p_x(t,0)$ и $-p_x(t,l)$.
\item При двустороннем наблюдении ключевое условие: $T>l$.
\item Студенческая нормировка дает $\mu=2(T-l)/l$, но константу надо сверять с первоисточником.
\item В конечномерной схеме главное - сохранить сопряженность: $\langle A_Nu_N,v_N\rangle=\langle u_N,A_N^*v_N\rangle$.
\item Конечная задача сводится к квадратичной минимизации $I_N(c)=c^TG_Nc-2b_N^Tc$.
\end{enumerate}

\blocktitle{E}{Пространства Соболева и задача оптимального нагрева стержня}{Билеты 15--17}
\section*{Замечание об источниках и надежности}
Основная литература по конкретной задаче оптимального нагрева стержня и лекционный источник по билетам 16--17 в текущем наборе материалов не были найдены. Поэтому билеты 16--17 ниже составлены на основе официальной формулировки вопроса, студенческих конспектов, стандартной теории слабых решений параболических уравнений и общих методов оптимизации. Если в лекциях преподавателя уравнение или функционал записаны с иными коэффициентами, например \(x_t=a^2x_{ss}+u\) или \(J=\frac12\|x(T)-b\|^2\), соответствующие коэффициенты в сопряженной задаче и градиенте нужно изменить согласованно.

\newpage

\section*{Билет 15}
\subsection*{Официальная формулировка}
\textbf{15. Обобщенная производная по Соболеву. Пространства Соболева. Теоремы вложения.}

\subsection*{Ответ}
Пусть \(\Omega\subset\R^n\) -- открытая область. Обозначим через
\[
C_c^\infty(\Omega)
\]
пространство бесконечно дифференцируемых функций с компактным носителем в \(\Omega\). Эти функции называются пробными функциями.

Пусть
\[
u\in L_{1,\mathrm{loc}}(\Omega).
\]
Для мультииндекса
\[
\alpha=(\alpha_1,\ldots,\alpha_n),\qquad
|\alpha|=\alpha_1+\cdots+\alpha_n
\]
обозначим
\[
D^\alpha=
\frac{\partial^{|\alpha|}}{\partial x_1^{\alpha_1}\cdots \partial x_n^{\alpha_n}}.
\]
Функция
\[
v_\alpha\in L_{1,\mathrm{loc}}(\Omega)
\]
называется обобщенной, или слабой, производной функции \(u\) порядка \(\alpha\), если для любой пробной функции
\[
\varphi\in C_c^\infty(\Omega)
\]
выполнено интегральное тождество
\[
\int_\Omega u(x)D^\alpha\varphi(x)\,dx
=
(-1)^{|\alpha|}
\int_\Omega v_\alpha(x)\varphi(x)\,dx.
\]
В этом случае пишут
\[
D^\alpha u=v_\alpha
\]
в обобщенном смысле.

Это определение получается из формулы интегрирования по частям. Главное отличие состоит в том, что производная переносится с функции \(u\) на гладкую пробную функцию \(\varphi\). Поэтому обобщенная производная может существовать даже тогда, когда классической производной у функции нет.

В одномерном случае, если
\[
f\in L_{1,\mathrm{loc}}(a,b),
\]
то функция \(\omega\in L_{1,\mathrm{loc}}(a,b)\) называется обобщенной производной \(f\), если
\[
\int_a^b f(x)\varphi'(x)\,dx
=
-
\int_a^b \omega(x)\varphi(x)\,dx
\]
для всех
\[
\varphi\in C_c^\infty(a,b).
\]

Основные свойства обобщенной производной:
\begin{enumerate}[leftmargin=2em]
\item обобщенная производная единственна с точностью до равенства почти всюду;
\item операция взятия обобщенной производной линейна:
\[
D^\alpha(\lambda u+
\mu v)=\lambda D^\alpha u+\mu D^\alpha v;
\]
\item если классическая производная существует и локально интегрируема, то она совпадает с обобщенной производной;
\item если на связном интервале \(u'=0\) в обобщенном смысле, то \(u\) равна константе почти всюду;
\item обобщенная производная позволяет рассматривать негладкие решения дифференциальных уравнений.
\end{enumerate}

\subsection*{Пространства Соболева}
Пусть
\[
1\le p\le\infty,
\qquad
m\in\mathbb N.
\]
Пространством Соболева \(W_p^m(\Omega)\) называется множество функций
\[
u\in L_p(\Omega),
\]
у которых все обобщенные производные
\[
D^\alpha u,
\qquad |\alpha|\le m,
\]
существуют и принадлежат \(L_p(\Omega)\). То есть
\[
W_p^m(\Omega)=
\{u\in L_p(\Omega):D^\alpha u\in L_p(\Omega),\ |\alpha|\le m\}.
\]

При \(1\le p<\infty\) норма задается формулой
\[
\|u\|_{W_p^m(\Omega)}=
\left(
\sum_{|\alpha|\le m}
\|D^\alpha u\|_{L_p(\Omega)}^p
\right)^{1/p}.
\]
При \(p=\infty\):
\[
\|u\|_{W_\infty^m(\Omega)}=
\max_{|\alpha|\le m}
\|D^\alpha u\|_{L_\infty(\Omega)}.
\]

При \(p=2\) пространство Соболева обозначают
\[
H^m(\Omega)=W_2^m(\Omega).
\]
Это гильбертово пространство со скалярным произведением
\[
(u,v)_{H^m(\Omega)}=
\sum_{|\alpha|\le m}
\int_\Omega D^\alpha u(x)D^\alpha v(x)\,dx.
\]

Например,
\[
H^1(a,b)=\{u\in L_2(a,b):u'\in L_2(a,b)\},
\]
и
\[
\|u\|_{H^1(a,b)}^2=
\int_a^b\left(u^2(x)+(u'(x))^2\right)dx.
\]

Пространство
\[
H_0^1(\Omega)
\]
определяется как замыкание \(C_c^\infty(\Omega)\) в норме \(H^1(\Omega)\):
\[
H_0^1(\Omega)=\overline{C_c^\infty(\Omega)}^{\,H^1}.
\]
Для достаточно регулярной области это означает, что функции из \(H_0^1(\Omega)\) имеют нулевой след на границе:
\[
u|_{\partial\Omega}=0
\]
в смысле следа.

\subsection*{Одномерный случай}
Если
\[
u\in H^1(a,b),
\]
то у \(u\) есть абсолютно непрерывный представитель, и для любых
\[
\alpha,\beta\in[a,b]
\]
выполнено
\[
u(\beta)-u(\alpha)=\int_\alpha^\beta u'(x)\,dx.
\]
Отсюда следует непрерывное вложение
\[
H^1(a,b)\hookrightarrow C[a,b].
\]
То есть каждая функция из \(H^1(a,b)\) имеет непрерывного представителя, и существует константа \(C>0\), такая что
\[
\|u\|_{C[a,b]}\le C\|u\|_{H^1(a,b)}.
\]
Это важно для краевых задач: в одномерном случае значения функции Соболева в точках имеют корректный смысл.

\subsection*{Теоремы вложения Соболева}
Теоремы вложения Соболева показывают, что наличие обобщенных производных повышает регулярность функции. Пусть \(\Omega\subset\R^n\) -- ограниченная область с достаточно регулярной границей.

Если
\[
mp<n,
\]
то имеет место непрерывное вложение
\[
W_p^m(\Omega)\hookrightarrow L_q(\Omega)
\]
для подходящих \(q\), в частности для
\[
q\le p^*=\frac{np}{n-mp}.
\]
Частный случай при \(m=1\):
\[
W_p^1(\Omega)\hookrightarrow L_{p^*}(\Omega),
\qquad
p^*=\frac{np}{n-p},
\qquad p<n.
\]

Если
\[
mp>n,
\]
то функции из \(W_p^m(\Omega)\) имеют более регулярные представители. В частности, при достаточном превышении гладкости над размерностью возникают вложения в пространства непрерывных или гёльдеровых функций:
\[
W_p^m(\Omega)\hookrightarrow C^{0,\gamma}(\overline\Omega)
\]
для некоторого \(\gamma>0\). В частности, при \(m=1\), \(p>n\) имеем вложение
\[
W_p^1(\Omega)\hookrightarrow C(\overline\Omega).
\]

Для \(p=2\) часто используют следующую форму. Если
\[
m-\frac n2>0,
\]
то \(H^m(\Omega)\) вкладывается в пространство непрерывных функций при достаточной регулярности \(\Omega\).

Кроме непрерывных вложений существуют компактные вложения Реллиха--Кондрашова. В типичной форме
\[
W_p^m(\Omega)\Subset W_q^l(\Omega)
\]
при более строгом условии регулярности, например когда
\[
m>l,
\qquad
m-\frac np>l-\frac nq.
\]
Здесь знак \(\Subset\) означает компактное вложение: из любой ограниченной последовательности в более сильном пространстве можно выбрать сходящуюся подпоследовательность в более слабом пространстве.

Для задач математической физики эти теоремы имеют принципиальное значение. Они позволяют корректно задавать граничные условия через следы, доказывать существование слабых решений, получать компактность минимизирующих последовательностей и переходить к пределу в вариационных задачах и задачах оптимального управления.

\subsection*{Связь с параболическими задачами}
В задачах типа уравнения теплопроводности часто используется гильбертов тройник
\[
H_0^1(0,l)\subset L_2(0,l)\simeq L_2^*(0,l)\subset H^{-1}(0,l).
\]
Если
\[
y\in L_2(0,T;H_0^1(0,l)),
\qquad
y_t\in L_2(0,T;H^{-1}(0,l)),
\]
то
\[
y\in C([0,T];L_2(0,l)).
\]
Именно это позволяет корректно говорить о начальном условии \(y(0)\) и терминальном значении \(y(T)\) в задачах управления нагревом стержня.

\subsection*{Краткий итог}
Обобщенная производная по Соболеву -- это производная, определенная через интегральное тождество с пробными функциями. Пространства Соболева \(W_p^m(\Omega)\) состоят из функций, у которых обобщенные производные до порядка \(m\) лежат в \(L_p\). Теоремы вложения показывают, в какие более регулярные пространства попадают функции Соболева при заданных \(m,p,n\). Эти понятия являются основой слабой постановки краевых задач для уравнений в частных производных.

\newpage

\section*{Билет 16}
\subsection*{Официальная формулировка}
\textbf{16. Постановка задачи об оптимальном нагреве стержня. Существование и единственность обобщенного решения краевой задачи. Теорема Вейерштраса.}

\subsection*{Ответ}
Рассмотрим стержень длины \(l\). Пусть
\[
Q=(0,T)\times(0,l),
\]
где \(t\in(0,T)\) -- время, \(s\in(0,l)\) -- пространственная координата. Через
\[
x=x(s,t)
\]
обозначим температуру стержня, а через
\[
u=u(s,t)
\]
-- распределенное управление, то есть интенсивность нагрева.

Задача оптимального нагрева стержня имеет вид
\[
J(u)=\int_0^l\bigl(x(s,T;u)-b(s)\bigr)^2\,ds\to\inf_{u\in U},
\]
где \(b(s)\) -- заданное желаемое распределение температуры в конечный момент времени \(T\).

Состояние \(x\) определяется уравнением теплопроводности
\[
x_t(s,t)=x_{ss}(s,t)+u(s,t),
\qquad (s,t)\in Q,
\]
с краевыми условиями
\[
x(0,t)=x(l,t)=0,
\qquad 0<t<T,
\]
и начальным условием
\[
x(s,0)=\varphi(s),
\qquad 0<s<l.
\]
Множество допустимых управлений задается как замкнутый шар в \(L_2(Q)\):
\[
U=\{u\in L_2(Q):\|u\|_{L_2(Q)}\le R\}.
\]
Физический смысл задачи: нужно подобрать распределенный нагрев \(u(s,t)\), ограниченный по энергии, так чтобы конечная температура \(x(s,T)\) была как можно ближе к заданному профилю \(b(s)\) в норме \(L_2(0,l)\).

\subsection*{Функциональные пространства}
Положим
\[
H=L_2(0,l),
\qquad
V=H_0^1(0,l),
\qquad
V^*=H^{-1}(0,l).
\]
Тогда имеем гильбертов тройник
\[
V\subset H\simeq H^*\subset V^*.
\]
Оператор теплопроводности задается слабой формой
\[
a(y,v)=\int_0^l y_s(s)v_s(s)\,ds,
\qquad y,v\in H_0^1(0,l).
\]
Соответствующий оператор
\[
A:V\to V^*
\]
определяется равенством
\[
\langle Ay,v\rangle_{V^*,V}=\int_0^l y_s(s)v_s(s)\,ds.
\]
Тогда уравнение теплопроводности можно записать абстрактно:
\[
\dot x(t)+Ax(t)=u(t),
\qquad x(0)=\varphi.
\]
Здесь \(u(t)\in H=L_2(0,l)\), а также \(u(t)\) можно рассматривать как элемент \(V^*\).

\subsection*{Обобщенное решение}
Пусть
\[
u\in L_2(Q)=L_2(0,T;L_2(0,l)),
\qquad
\varphi\in L_2(0,l).
\]
Функция \(x=x(s,t)\) называется обобщенным решением краевой задачи, если
\[
x\in L_2(0,T;H_0^1(0,l)),
\]
\[
x_t\in L_2(0,T;H^{-1}(0,l)),
\]
\[
x(0)=\varphi\quad \text{в }L_2(0,l),
\]
и для почти всех \(t\in(0,T)\) выполнено тождество
\[
\langle x_t(t),v\rangle_{H^{-1},H_0^1}
+
\int_0^l x_s(s,t)v_s(s)\,ds
=
\int_0^l u(s,t)v(s)\,ds
\]
для всех
\[
v\in H_0^1(0,l).
\]
Это слабая форма уравнения \(x_t=x_{ss}+u\).

В интегральной форме, близкой к записи студенческого варианта, это можно записать так: для всех достаточно гладких пробных функций \(\Phi(s,t)\), обращающихся в нуль при \(s=0,l\),
\[
\int_0^l x(s,T)\Phi(s,T)\,ds
+
\iint_Q\left(-x(s,t)\Phi_t(s,t)+x_s(s,t)\Phi_s(s,t)\right)dsdt
\]
\[
=
\int_0^l \varphi(s)\Phi(s,0)\,ds
+
\iint_Q u(s,t)\Phi(s,t)\,dsdt.
\]
Главное преимущество такой формулировки состоит в том, что она не требует существования классических производных \(x_t\) и \(x_{ss}\) как обычных функций.

Из стандартной теоремы о функциях из пространства
\[
L_2(0,T;H_0^1(0,l)),
\qquad
x_t\in L_2(0,T;H^{-1}(0,l))
\]
следует вложение
\[
x\in C([0,T];L_2(0,l)).
\]
Поэтому значения \(x(0)\) и \(x(T)\) имеют корректный смысл в \(L_2(0,l)\). Именно это делает функционал
\[
J(u)=\|x(T;u)-b\|_{L_2(0,l)}^2
\]
строго определенным.

\subsection*{Существование обобщенного решения}
Для каждого
\[
u\in L_2(Q),
\qquad
\varphi\in L_2(0,l)
\]
существует обобщенное решение
\[
x\in L_2(0,T;H_0^1(0,l)),
\qquad
x_t\in L_2(0,T;H^{-1}(0,l)).
\]
Это следует из стандартной теории линейных параболических уравнений. Билинейная форма
\[
a(y,v)=\int_0^l y_s v_s\,ds
\]
непрерывна и коэрцитивна на \(H_0^1(0,l)\). Правая часть \(u\in L_2(0,T;L_2(0,l))\) принадлежит также \(L_2(0,T;H^{-1}(0,l))\). Поэтому задача
\[
\dot x+Ax=u,
\qquad x(0)=\varphi
\]
имеет решение в указанном энергетическом классе.

\subsection*{Единственность обобщенного решения}
Пусть \(x_1\) и \(x_2\) -- два обобщенных решения при одном и том же управлении \(u\) и одном и том же начальном условии \(\varphi\). Рассмотрим их разность
\[
z=x_1-x_2.
\]
Тогда \(z\) удовлетворяет однородной задаче
\[
z_t=z_{ss},
\qquad z(0,t)=z(l,t)=0,
\qquad z(s,0)=0.
\]
В слабой форме
\[
\langle z_t(t),v\rangle+
\int_0^l z_s(s,t)v_s(s)\,ds=0.
\]
Берем \(v=z(t)\). Тогда получаем энергетическое равенство
\[
\frac12\frac{d}{dt}\|z(t)\|_{L_2(0,l)}^2+
\|z_s(t)\|_{L_2(0,l)}^2=0.
\]
Интегрируя по времени от \(0\) до \(t\), получаем
\[
\frac12\|z(t)\|_{L_2(0,l)}^2+
\int_0^t\|z_s(\tau)\|_{L_2(0,l)}^2\,d\tau=0.
\]
Следовательно,
\[
z(t)=0
\]
для всех \(t\in[0,T]\) в \(L_2(0,l)\). Значит,
\[
x_1=x_2.
\]
Итак, обобщенное решение краевой задачи существует и единственно.

\subsection*{Теорема Вейерштрасса и существование оптимального управления}
Обозначим через
\[
S:u\mapsto x(T;u)
\]
оператор, который каждому управлению \(u\in L_2(Q)\) ставит в соответствие конечное состояние \(x(T;u)\in L_2(0,l)\). Поскольку уравнение теплопроводности линейно,
\[
x(T;u)=x^0(T)+S_Tu,
\]
где \(x^0\) -- решение при \(u=0\), а \(S_T\) -- линейный непрерывный оператор. Функционал можно записать так:
\[
J(u)=\|x^0(T)+S_Tu-b\|_{L_2(0,l)}^2.
\]
Следовательно, \(J\) является выпуклым и слабо полунепрерывным снизу на \(L_2(Q)\).

Множество допустимых управлений
\[
U=\{u\in L_2(Q):\|u\|_{L_2(Q)}\le R\}
\]
является непустым, выпуклым, замкнутым и ограниченным в гильбертовом пространстве \(L_2(Q)\). Так как \(L_2(Q)\) рефлексивно, множество \(U\) слабо компактно.

Пусть
\[
\{u_k\}\subset U
\]
-- минимизирующая последовательность:
\[
J(u_k)\to \inf_{u\in U}J(u).
\]
Из слабой компактности \(U\) можно выбрать подпоследовательность, которую снова обозначим \(\{u_k\}\), такую что
\[
u_k\rightharpoonup u_*
\quad \text{в }L_2(Q),
\]
причем
\[
u_*\in U.
\]
Так как \(J\) слабо полунепрерывен снизу,
\[
J(u_*)\le \liminf_{k\to\infty}J(u_k).
\]
Но \(\{u_k\}\) -- минимизирующая последовательность, поэтому
\[
\lim_{k\to\infty}J(u_k)=\inf_{u\in U}J(u).
\]
Следовательно,
\[
J(u_*)=\inf_{u\in U}J(u).
\]
Значит, оптимальное управление существует.

Это и есть применение теоремы Вейерштрасса в слабой форме: слабо полунепрерывный снизу функционал достигает минимума на слабо компактном множестве.

Важно: из этого рассуждения следует существование оптимального управления, но не обязательно его единственность. Единственность оптимального управления требует дополнительных условий, например строгой выпуклости функционала по \(u\).

\subsection*{Краткий итог}
Задача оптимального нагрева стержня сводится к минимизации функционала
\[
J(u)=\|x(T;u)-b\|_{L_2(0,l)}^2
\]
по управлению \(u\in U\subset L_2(Q)\). Для каждого \(u\in L_2(Q)\) краевая задача для уравнения теплопроводности имеет единственное обобщенное решение
\[
x\in L_2(0,T;H_0^1(0,l)),
\qquad
x_t\in L_2(0,T;H^{-1}(0,l)).
\]
Благодаря вложению
\[
x\in C([0,T];L_2(0,l))
\]
конечное состояние \(x(T)\) корректно определено. Существование оптимального управления следует из слабой компактности множества \(U\) и слабой полунепрерывности снизу функционала \(J\).

\newpage

\section*{Билет 17}
\subsection*{Официальная формулировка}
\textbf{17. Градиент целевого функционала в задаче о нагреве стержня. Критерий оптимальности. Регуляризованный метод проекции градиента.}

\subsection*{Ответ}
Рассмотрим задачу оптимального нагрева стержня:
\[
J(u)=\int_0^l (x(s,T;u)-b(s))^2\,ds\to\inf_{u\in U},
\]
где состояние \(x=x(s,t;u)\) удовлетворяет уравнению теплопроводности
\[
x_t=x_{ss}+u(s,t),
\qquad (s,t)\in Q=(0,T)\times(0,l),
\]
с краевыми условиями
\[
x(0,t)=x(l,t)=0,
\]
и начальным условием
\[
x(s,0)=\varphi(s).
\]
Множество допустимых управлений:
\[
U=\{u\in L_2(Q):\|u\|_{L_2(Q)}\le R\}.
\]

\subsection*{Градиент целевого функционала}
Пусть \(u\in U\), а \(x=x(s,t;u)\) -- соответствующее решение прямой задачи. Возьмем приращение управления
\[
u+\varepsilon v,
\qquad v\in L_2(Q).
\]
Тогда вариация состояния
\[
z=\frac{d}{d\varepsilon}x(u+\varepsilon v)\bigg|_{\varepsilon=0}
\]
удовлетворяет линеаризованной задаче
\[
z_t=z_{ss}+v,
\]
\[
z(0,t)=z(l,t)=0,
\]
\[
z(s,0)=0.
\]
Первая вариация функционала равна
\[
\delta J(u;v)=
2\int_0^l (x(s,T;u)-b(s))z(s,T)\,ds.
\]

Чтобы выразить ее через \(v\), вводят сопряженную функцию \(\psi=\psi(s,t)\), решающую задачу
\[
-\psi_t=\psi_{ss},
\qquad (s,t)\in Q,
\]
\[
\psi(0,t)=\psi(l,t)=0,
\]
\[
\psi(s,T)=2(x(s,T;u)-b(s)).
\]
Эквивалентно первое уравнение можно записать как
\[
-\psi_t-\psi_{ss}=0.
\]
Умножая уравнение для \(z\) на \(\psi\), интегрируя по \(Q\) и используя интегрирование по частям, получаем
\[
\delta J(u;v)=
\iint_Q \psi(s,t)v(s,t)\,dsdt.
\]
Следовательно, производная Фреше функционала имеет вид
\[
\langle J'(u),v\rangle_{L_2(Q)}=
\iint_Q \psi(s,t)v(s,t)\,dsdt.
\]
Так как пространство управлений есть \(L_2(Q)\), то градиент функционала равен
\[
J'(u)=\psi.
\]

Итак, для вычисления градиента нужно решить прямую задачу для \(x(s,t;u)\), задать терминальное условие
\[
\psi(s,T)=2(x(s,T;u)-b(s)),
\]
решить сопряженную задачу назад по времени и получить \(J'(u)=\psi\).

\subsection*{Критерий оптимальности}
Пусть
\[
U=\{u\in L_2(Q):\|u\|\le R\}
\]
-- замкнутый выпуклый шар. Если \(u_*\in U\) -- оптимальное управление, то для выпуклого дифференцируемого функционала необходимое и достаточное условие оптимальности имеет вид вариационного неравенства:
\[
\langle J'(u_*),u-u_*\rangle_{L_2(Q)}\ge0
\qquad \forall u\in U.
\]
Так как \(J'(u_*)=\psi_*\), это можно записать как
\[
\iint_Q \psi_*(s,t)(u(s,t)-u_*(s,t))\,dsdt\ge0
\qquad \forall u\in U.
\]
Эквивалентная проекционная форма:
\[
u_*=P_U\bigl(u_* -\beta J'(u_*)\bigr)
\qquad \forall \beta>0.
\]

Для шара \(U\) условие можно записать в форме Куна--Таккера. Рассмотрим ограничение
\[
\|u\|^2-R^2\le0
\]
и функцию Лагранжа
\[
L(u,\lambda)=J(u)+\lambda(\|u\|^2-R^2),
\qquad \lambda\ge0.
\]
Тогда оптимальность \(u_*\) эквивалентна существованию \(\lambda_*\ge0\), такого что
\[
J'(u_*)+2\lambda_*u_*=0,
\]
\[
\lambda_*(\|u_*\|^2-R^2)=0,
\]
\[
\|u_*\|\le R.
\]
Если ограничение не активно, то
\[
\|u_*\|<R,
\qquad
\lambda_*=0,
\qquad
J'(u_*)=0.
\]
Если ограничение активно, то \(\|u_*\|=R\). Из условия \(J'(u_*)+2\lambda_*u_*=0\) следует
\[
\lambda_*=\frac{\|J'(u_*)\|}{2R}.
\]
Если используется запись
\[
J'(u_*)+\mu_*u_*=0,
\]
то \(\mu_*=2\lambda_*\), и в активном случае
\[
\mu_*=\frac{\|J'(u_*)\|}{R}.
\]
Именно это различие объясняет разные коэффициенты в студенческих версиях ответа.

Критерий оптимальности можно также записать в седловой форме. Точка \((u_*,\lambda_*)\) является седловой точкой функции
\[
L(u,\lambda)=J(u)+\lambda(\|u\|^2-R^2),
\qquad \lambda\ge0,
\]
если
\[
L(u_*,\lambda)\le L(u_*,\lambda_*)\le L(u,\lambda_*)
\]
для всех \(u\in L_2(Q)\), \(\lambda\ge0\). В этом случае \(u_*\) является оптимальным управлением.

\subsection*{Регуляризованный метод проекции градиента}
Обычный метод проекции градиента для задачи
\[
J(u)\to\inf,
\qquad u\in U
\]
имеет вид
\[
u_{k+1}=P_U\left(u_k-\beta_kJ'(u_k)\right),
\qquad k=0,1,\ldots .
\]
Здесь \(P_U\) -- метрическая проекция на замкнутое выпуклое множество \(U\), а \(\beta_k>0\) -- шаг.

В регуляризованном методе вводится функция Тихонова
\[
T_k(u)=J(u)+\alpha_k\|u\|^2,
\qquad \alpha_k>0.
\]
Ее градиент:
\[
T_k'(u)=J'(u)+2\alpha_k u.
\]
Если градиент вычисляется приближенно, используют \(J_k'(u)\) и предполагают оценку ошибки вида
\[
\|J_k'(u)-J'(u)\|\le \delta_k(1+\|u\|).
\]
Тогда регуляризованный метод проекции градиента задается формулой
\[
u_{k+1}=P_U\left(u_k-\beta_k\bigl(J_k'(u_k)+2\alpha_k u_k\bigr)\right).
\]
Если градиент точный, то \(J_k'(u_k)=J'(u_k)\), и формула принимает вид
\[
u_{k+1}=P_U\left(u_k-\beta_k\bigl(J'(u_k)+2\alpha_k u_k\bigr)\right).
\]
В задаче нагрева стержня это означает
\[
u_{k+1}=P_U\left(u_k-\beta_k\bigl(\psi_k+2\alpha_k u_k\bigr)\right),
\]
где \(\psi_k\) -- решение сопряженной задачи, соответствующее управлению \(u_k\).

\subsection*{Проекция на шар}
Если
\[
U=\{u\in L_2(Q):\|u\|\le R\},
\]
то проекция имеет явный вид:
\[
P_U(w)=
\begin{cases}
w, & \|w\|\le R,\\[2mm]
\dfrac{R}{\|w\|}w, & \|w\|>R.
\end{cases}
\]
Поэтому один шаг метода можно записать так:
\[
w_k=u_k-\beta_k(\psi_k+2\alpha_k u_k),
\]
\[
u_{k+1}=P_U(w_k).
\]

\subsection*{Условия сходимости}
Типичные условия сходимости регуляризованного метода:
\begin{enumerate}[leftmargin=2em]
\item \(U\) -- выпуклое замкнутое множество в гильбертовом пространстве;
\item \(J\) выпукл и непрерывно дифференцируем по Фреше;
\item множество оптимальных решений непусто;
\item градиент \(J'\) удовлетворяет условию типа Липшица или коэрцитивности;
\item параметры регуляризации стремятся к нулю:
\[
\alpha_k\to0;
\]
\item ошибка приближенного градиента согласована с регуляризацией:
\[
\frac{\delta_k}{\alpha_k}\to0;
\]
\item шаги \(\beta_k\) выбраны так, чтобы обеспечивать устойчивость и убывание регуляризованного функционала.
\end{enumerate}
При этих условиях последовательность \(u_k\) сходится к нормальному оптимальному управлению, то есть к оптимальному управлению минимальной нормы:
\[
u_k\to u_*,
\]
где
\[
u_*=
\operatorname*{arg\,min}\{\|u\|:u\in U^*\}.
\]
Здесь \(U^*\) -- множество оптимальных управлений.

\subsection*{Краткий итог}
В задаче нагрева стержня градиент целевого функционала находится через сопряженную задачу:
\[
J'(u)=\psi.
\]
Оптимальность на выпуклом множестве \(U\) задается вариационным неравенством
\[
\langle J'(u_*),u-u_*\rangle\ge0
\qquad \forall u\in U.
\]
Для шара это эквивалентно условиям Куна--Таккера
\[
J'(u_*)+2\lambda_*u_*=0,
\]
\[
\lambda_*(\|u_*\|^2-R^2)=0,
\]
\[
\lambda_*\ge0,
\qquad
\|u_*\|\le R.
\]
Регуляризованный метод проекции градиента имеет вид
\[
u_{k+1}=P_U\left(u_k-\beta_k(J_k'(u_k)+2\alpha_k u_k)\right).
\]

\blocktitle{F}{Линейно-квадратичные дифференциальные игры и максимум по Парето}{Билеты 18--20}
\section*{Замечание об источниках и надежности}
\addcontentsline{toc}{section}{Замечание об источниках и надежности}
Официальная программа указывает в качестве основных источников для этого блока книги В.И.~Жуковского и соавторов по линейно-квадратичным дифференциальным играм, гарантированным решениям конфликтов, уравновешиванию конфликтов, а также дополнительные источники по векторной оптимизации и многокритериальным динамическим системам. Полные тексты некоторых книг в текущем наборе материалов отсутствуют. Поэтому ниже используются: официальные формулировки билетов, студенческие версии 2017 и 2026 гг., а также открытые статьи В.И.~Жуковского и соавторов, непосредственно относящиеся к билетам 18--20.

Для билета 18 главным источником является статья \emph{«Об одной нерешенной задаче в матричных обыкновенных дифференциальных уравнениях»}. Для билета 19 главным источником является статья \emph{«Berge--Vaisman equilibrium for one linear-quadratic differential game»}. Для билета 20 используются открытые материалы о максимуме по Парето в линейно-квадратичных динамических играх, но отдельные детали желательно сверить с книгами по векторной оптимизации динамических систем.

\ticket{18}{Условия существования и явный вид ситуации равновесия по Нэшу в дифференциальной линейно-квадратичной дифференциальной игре трех лиц}
\official{18. Условия существования и явный вид ситуации равновесия по Нэшу в дифференциальной линейно-квадратичной дифференциальной игре трех лиц.}

\fullanswer
Рассмотрим дифференциальную позиционную линейно-квадратичную игру трех лиц
\[
\Gamma=\langle \{1,2,3\},\Sigma,\{\mathcal A_i\}_{i=1}^3,\{J_i\}_{i=1}^3\rangle .
\]
Фазовый вектор и управляющие воздействия имеют размерность
\[
x(t)\in\R^n,\qquad u_i(t)\in\R^n,\quad i=1,2,3.
\]
Динамика задается системой
\[
\dot x(t)=A(t)x(t)+u_1(t)+u_2(t)+u_3(t),\qquad x(t_0)=x_0,
\]
где \(A(\cdot)\in C^{n\times n}[0,\vartheta]\), а \(\vartheta>t_0\) -- фиксированный момент окончания игры.

Игроки используют линейные позиционные стратегии
\[
U_i:\quad u_i(t,x)=Q_i(t)x,\qquad Q_i(\cdot)\in C^{n\times n}[0,\vartheta].
\]
Множество стратегий \(i\)-го игрока:
\[
\mathcal A_i=\{U_i:\ u_i(t,x)=Q_i(t)x,\ Q_i(\cdot)\in C^{n\times n}[0,\vartheta]\}.
\]
Ситуация игры есть тройка
\[
U=(U_1,U_2,U_3)\in \mathcal A_1\times\mathcal A_2\times\mathcal A_3.
\]
При фиксированной ситуации \(U\) движение удовлетворяет
\[
\dot x(t)=\bigl(A(t)+Q_1(t)+Q_2(t)+Q_3(t)\bigr)x(t),\qquad x(t_0)=x_0.
\]

Функция выигрыша \(i\)-го игрока задается квадратичным функционалом
\[
J_i(U,t_0,x_0)=x^{\top}(\vartheta)C_i x(\vartheta)+
\int_{t_0}^{\vartheta}\sum_{j=1}^3 u_j^{\top}(t)D_{ij}u_j(t)\,dt,
\qquad i=1,2,3,
\]
где \(C_i,D_{ij}\) -- постоянные симметричные \(n\times n\)-матрицы. Каждый игрок стремится максимизировать свой выигрыш.

\subsection*{Равновесие по Нэшу}
Ситуация \(U^e=(U_1^e,U_2^e,U_3^e)\) называется равновесием по Нэшу, если для любой начальной позиции
\[
(t_0,x_0)\in [0,\vartheta)\times\R^n,\qquad x_0\ne 0,
\]
выполнены условия
\[
\max_{U_1\in\mathcal A_1}J_1(U_1,U_2^e,U_3^e,t_0,x_0)=J_1(U^e,t_0,x_0),
\]
\[
\max_{U_2\in\mathcal A_2}J_2(U_1^e,U_2,U_3^e,t_0,x_0)=J_2(U^e,t_0,x_0),
\]
\[
\max_{U_3\in\mathcal A_3}J_3(U_1^e,U_2^e,U_3,t_0,x_0)=J_3(U^e,t_0,x_0).
\]
Смысл: ни один игрок не может увеличить свой выигрыш, если остальные два сохраняют свои равновесные стратегии.

\subsection*{Метод динамического программирования}
Для построения равновесия вводятся функции Беллмана
\[
V_i(t,x),\qquad i=1,2,3.
\]
Определим
\[
W_i(t,x,u_1,u_2,u_3,V_i)=
\frac{\partial V_i}{\partial t}+
\left(\frac{\partial V_i}{\partial x}\right)^{\top}(A(t)x+u_1+u_2+u_3)+
\sum_{j=1}^3u_j^{\top} D_{ij}u_j .
\]
В равновесии по Нэшу \(i\)-й игрок максимизирует \(W_i\) только по собственному управлению \(u_i\), считая управления остальных игроков фиксированными. Для конечного максимума по \(u_i\) требуется вогнутость по \(u_i\), то есть
\[
D_{ii}<0,
\qquad i=1,2,3.
\]
Условие первого порядка дает
\[
\frac{\partial W_i}{\partial u_i}=\frac{\partial V_i}{\partial x}+2D_{ii}u_i=0.
\]
Следовательно,
\[
u_i^e(t,x,V_i)=-\frac12D_{ii}^{-1}\frac{\partial V_i}{\partial x}.
\]

В линейно-квадратичном случае функции Беллмана ищутся в виде
\[
V_i(t,x)=x^{\top}\Theta_i(t)x,
\qquad \Theta_i^{\top}(t)=\Theta_i(t).
\]
Тогда
\[
\frac{\partial V_i}{\partial x}=2\Theta_i(t)x,
\]
и равновесные стратегии имеют вид
\[
u_i^e(t,x)=-D_{ii}^{-1}\Theta_i(t)x,
\qquad i=1,2,3.
\]
Итак,
\[
Q_i^e(t)=-D_{ii}^{-1}\Theta_i(t).
\]

\subsection*{Система матричных уравнений типа Риккати}
Подстановка \(V_i=x^{\top}\Theta_i x\) и \(u_i^e=-D_{ii}^{-1}\Theta_i x\) в уравнения динамического программирования приводит к связанной системе матричных уравнений типа Риккати. В симметричной компактной записи для \(i=1,2,3\):
\[
\begin{aligned}
0={}&\dot\Theta_i+\Theta_iA+A^{\top}\Theta_i-\Theta_iD_{ii}^{-1}\Theta_i \\
&-\sum_{\substack{j=1\\ j\ne i}}^3
\left(\Theta_iD_{jj}^{-1}\Theta_j+\Theta_jD_{jj}^{-1}\Theta_i\right)
+\sum_{\substack{j=1\\ j\ne i}}^3
\Theta_jD_{jj}^{-1}D_{ij}D_{jj}^{-1}\Theta_j .
\end{aligned}
\]
Терминальные условия:
\[
\Theta_i(\vartheta)=C_i,\qquad i=1,2,3.
\]

Эта связанная система является главным техническим объектом билета. В общем случае построение ее явного решения для трех игроков является нетривиальной задачей; поэтому корректная экзаменационная формулировка должна звучать условно: если такая система имеет продолжимое решение, то соответствующие стратегии образуют равновесие по Нэшу.

\subsection*{Достаточное условие существования}
Если выполнены условия
\[
D_{ii}<0,
\qquad i=1,2,3,
\]
и система матричных уравнений типа Риккати с терминальными условиями \(\Theta_i(\vartheta)=C_i\) имеет продолжимое на \([0,\vartheta]\) решение в классе симметричных матриц
\[
\Theta_i(t)=\Theta_i^{\top}(t),
\qquad i=1,2,3,
\]
то ситуация
\[
U^e=(U_1^e,U_2^e,U_3^e),
\qquad
U_i^e:\ u_i^e(t,x)=-D_{ii}^{-1}\Theta_i(t)x,
\]
является равновесием по Нэшу. Равновесные выигрыши равны
\[
J_i^e(t_0,x_0)=J_i(U^e,t_0,x_0)=x_0^{\top}\Theta_i(t_0)x_0,
\qquad i=1,2,3.
\]

\subsection*{Случай отсутствия равновесия}
Если хотя бы для одного игрока
\[
D_{ii}>0,
\]
то максимум по \(u_i\) в определении равновесия по Нэшу не существует: игрок может увеличивать выигрыш за счет неограниченного усиления своей линейной стратегии. Следовательно, в рассматриваемом классе линейных позиционных стратегий равновесие по Нэшу отсутствует.

\cnexplain
本题的核心不是单纯写出一个 Riccati 方程，而是说明三人线性二次位置微分博弈中的 Nash 均衡如何由动态规划法构造。每个玩家选择线性反馈策略
\[
u_i(t,x)=Q_i(t)x.
\]
Nash 均衡表示任意一个玩家单独改变策略时，不能提高自己的收益。由于每个玩家最大化自己的收益，所以关于本玩家控制 \(u_i\) 的二次项必须是负定的：
\[
D_{ii}<0.
\]
否则最大化问题可能无解。

设 Bellman 函数为
\[
V_i(t,x)=x^{\top}\Theta_i(t)x.
\]
对 \(u_i\) 最大化得到
\[
u_i^e(t,x)=-D_{ii}^{-1}\Theta_i(t)x.
\]
把这个形式代回 Bellman 方程，得到三个矩阵函数 \(\Theta_1,\Theta_2,\Theta_3\) 的耦合 Riccati 型方程组。只有当这个方程组在整个区间上有对称矩阵解时，才能得到 Nash 均衡。

最重要的谨慎点：三人情形下 Riccati 方程组不是自动有显式解，不能写成“总是存在”。正确表述是“若该系统有延拓到整个区间的解，则相应策略构成 Nash 均衡”。

\corrections
\begin{enumerate}
\item 不应把本题主要写成“两名玩家加一个不确定性”的模型。官方题目写的是三人游戏，核心模型应为
\[
\dot x=A(t)x+u_1+u_2+u_3.
\]
\item 必须说明 \(D_{ii}<0\) 的原因：本题是最大化问题，最大值存在需要关于本玩家控制的二次型严格凹。
\item Riccati 方程必须带终端条件 \(\Theta_i(\vartheta)=C_i\)。没有终端条件，方程组不完整。
\item 不能声称三人耦合 Riccati 系统一定显式可解。可靠表述是“若该系统存在延拓解，则得到均衡”。
\item 所有矩阵维度必须一致：\(x,u_i\in\R^n\)，\(C_i,D_{ij},\Theta_i\in\R^{n\times n}\)，二次型写作 \(x^{\top} C_i x\)、\(u_j^{\top} D_{ij}u_j\)。
\end{enumerate}

\memory
Модель:
\[
\dot x=A(t)x+u_1+u_2+u_3,
\qquad u_i(t,x)=Q_i(t)x.
\]
Выигрыши:
\[
J_i=x^{\top}(\vartheta)C_ix(\vartheta)+\int_{t_0}^{\vartheta}\sum_{j=1}^3u_j^{\top}D_{ij}u_j\,dt.
\]
Равновесие по Нэшу:
\[
J_i(U^e)\ge J_i(U_i,U_{-i}^e)\quad \forall U_i,\\quad i=1,2,3.
\]
Ищем
\[
V_i(t,x)=x^{\top}\Theta_i(t)x.
\]
Условие максимума по \(u_i\):
\[
V_{i,x}+2D_{ii}u_i=0.
\]
Если \(D_{ii}<0\), то
\[
u_i^e(t,x)=-D_{ii}^{-1}\Theta_i(t)x.
\]
Матрицы \(\Theta_i\) удовлетворяют связанной системе Риккати с условиями
\[
\Theta_i(\vartheta)=C_i.
\]
Если система имеет продолжимое симметричное решение на \([0,\vartheta]\), то найденные стратегии образуют равновесие по Нэшу, а
\[
J_i^e(t_0,x_0)=x_0^{\top}\Theta_i(t_0)x_0.
\]
Если хотя бы одно \(D_{ii}>0\), равновесие по Нэшу отсутствует.

\newpage
\ticket{19}{Условия существования и явный вид ситуации равновесия по Бержу в дифференциальной линейно-квадратичной игре двух лиц с малым влиянием одного из игроков на скорость изменения фазового вектора}
\official{19. Условия существования и явный вид ситуации равновесия по Бержу в дифференциальной линейно-квадратичной игре двух лиц с малым влиянием одного из игроков на скорость изменения фазового вектора.}

\fullanswer
Рассмотрим дифференциальную позиционную линейно-квадратичную игру двух лиц
\[
\Gamma=\langle \{1,2\},\Sigma,\{\mathcal A_i\}_{i=1}^2,\{J_i\}_{i=1}^2\rangle .
\]
Динамика фазового вектора \(x(t)\in\R^n\) задается системой
\[
\dot x(t)=A(t)x(t)+u_1(t)+\varepsilon u_2(t),
\qquad x(t_0)=x_0,
\]
где \(A(\cdot)\in C^{n\times n}[0,\vartheta]\), а \(\varepsilon\ge 0\) -- малый параметр. Содержательно \(\varepsilon u_2\) означает малое влияние второго игрока на скорость изменения фазового вектора.

Игроки используют линейные позиционные стратегии
\[
U_i:\quad u_i(t,x)=Q_i(t)x,
\qquad Q_i(\cdot)\in C^{n\times n}[0,\vartheta],
\quad i=1,2.
\]
При фиксированной ситуации \(U=(U_1,U_2)\) движение имеет вид
\[
\dot x(t)=\bigl(A(t)+Q_1(t)+\varepsilon Q_2(t)\bigr)x(t).
\]

Функции выигрыша:
\[
J_1(U,t_0,x_0)=x^{\top}(\vartheta)C_1x(\vartheta)+
\int_{t_0}^{\vartheta}\left(u_1^{\top}D_{11}u_1-u_2^{\top}D_{12}u_2+x^{\top}G_1x\right)dt,
\]
\[
J_2(U,t_0,x_0)=x^{\top}(\vartheta)C_2x(\vartheta)+
\int_{t_0}^{\vartheta}\left(-u_1^{\top}D_{21}u_1+u_2^{\top}D_{22}u_2+x^{\top}G_2x\right)dt.
\]
Матрицы \(C_i,G_i,D_{ij}\) симметричны. В данной записи отрицательные знаки перед чужими управлениями вынесены явно, поэтому условия вогнутости формулируются как \(D_{12}>0\), \(D_{21}>0\).

\subsection*{Равновесие по Бержу}
Равновесие по Бержу отличается от равновесия по Нэшу. В Нэше каждый игрок максимизирует свой выигрыш по собственной стратегии. В равновесии по Бержу выигрыш каждого игрока максимизируется за счет стратегии другого игрока.

Ситуация \(U^B=(U_1^B,U_2^B)\) называется равновесием по Бержу, если
\[
J_1(U_1^B,U_2,t_0,x_0)\le J_1(U^B,t_0,x_0)
\qquad \forall U_2\in\mathcal A_2,
\]
\[
J_2(U_1,U_2^B,t_0,x_0)\le J_2(U^B,t_0,x_0)
\qquad \forall U_1\in\mathcal A_1.
\]
В строгой версии источника используется равновесие по Бержу--Вайсману, где дополнительно учитывается индивидуальная рациональность.

\subsection*{Динамическое программирование}
Введем функции Беллмана \(V_1(t,x)\), \(V_2(t,x)\) и выражения
\[
\begin{aligned}
W_1={}&\frac{\partial V_1}{\partial t}+\left(\frac{\partial V_1}{\partial x}\right)^{\top}(Ax+u_1+\varepsilon u_2)
+u_1^{\top}D_{11}u_1-u_2^{\top}D_{12}u_2+x^{\top}G_1x,\\
W_2={}&\frac{\partial V_2}{\partial t}+\left(\frac{\partial V_2}{\partial x}\right)^{\top}(Ax+u_1+\varepsilon u_2)
-u_1^{\top}D_{21}u_1+u_2^{\top}D_{22}u_2+x^{\top}G_2x.
\end{aligned}
\]
Для равновесия по Бержу условия максимума имеют перекрестный вид:
\[
\max_{u_2}W_1(t,x,u_1^B,u_2,V_1)=W_1(t,x,u_1^B,u_2^B,V_1),
\]
\[
\max_{u_1}W_2(t,x,u_1,u_2^B,V_2)=W_2(t,x,u_1^B,u_2^B,V_2).
\]
Условия первого порядка:
\[
\frac{\partial W_1}{\partial u_2}=\varepsilon\frac{\partial V_1}{\partial x}-2D_{12}u_2=0,
\]
\[
\frac{\partial W_2}{\partial u_1}=\frac{\partial V_2}{\partial x}-2D_{21}u_1=0.
\]
Отсюда
\[
u_2^B(t,x,V_1)=\frac{\varepsilon}{2}D_{12}^{-1}\frac{\partial V_1}{\partial x},
\qquad
u_1^B(t,x,V_2)=\frac12D_{21}^{-1}\frac{\partial V_2}{\partial x}.
\]

Ищем функции Беллмана в виде
\[
V_i(t,x)=x^{\top}\Theta_i(t,\varepsilon)x,
\qquad \Theta_i^{\top}(t,\varepsilon)=\Theta_i(t,\varepsilon).
\]
Тогда
\[
\frac{\partial V_i}{\partial x}=2\Theta_i(t,\varepsilon)x,
\]
и равновесные стратегии имеют вид
\[
u_1^B(t,x)=D_{21}^{-1}\Theta_2(t,\varepsilon)x,
\]
\[
u_2^B(t,x)=\varepsilon D_{12}^{-1}\Theta_1(t,\varepsilon)x.
\]
Следовательно,
\[
Q_1^B(t,\varepsilon)=D_{21}^{-1}\Theta_2(t,\varepsilon),
\qquad
Q_2^B(t,\varepsilon)=\varepsilon D_{12}^{-1}\Theta_1(t,\varepsilon).
\]

\subsection*{Матричная система и малый параметр}
Подстановка найденных управлений в уравнения динамического программирования приводит к системе матричных уравнений для \(\Theta_1\), \(\Theta_2\). В форме, удобной для экзаменационного ответа, ее можно записать так:
\[
\begin{aligned}
0={}&\dot\Theta_1+
\Theta_1(A+D_{21}^{-1}\Theta_2)+(A^{\top}+\Theta_2D_{21}^{-1})\Theta_1+G_1 \\
&+\Theta_2D_{21}^{-1}D_{11}D_{21}^{-1}\Theta_2
+\varepsilon^2\Psi_1(\Theta_1,\Theta_2),
\qquad \Theta_1(\vartheta,\varepsilon)=C_1,
\end{aligned}
\]
\[
\begin{aligned}
0={}&\dot\Theta_2+
\Theta_2A+A^{\top}\Theta_2+G_2+\Theta_2D_{21}^{-1}\Theta_2
+\varepsilon^2\Psi_2(\Theta_1,\Theta_2),\\
&\hspace{8cm}\Theta_2(\vartheta,\varepsilon)=C_2.
\end{aligned}
\]
Здесь \(\Psi_1,\Psi_2\) обозначают слагаемые, квадратичные по матрицам \(\Theta_1,\Theta_2\), возникающие из членов с \(\varepsilon u_2\). Главное: при \(\varepsilon=0\) система упрощается, а при достаточно малом \(\varepsilon\) решение продолжается по теореме о зависимости решений от параметра.

\subsection*{Условия существования}
Если \(\varepsilon\ge 0\) достаточно мало и выполнены условия
\[
C_2\le 0,
\qquad G_2\le 0,
\qquad D_{ij}>0,
\quad i,j=1,2,
\]
то ситуация равновесия по Нэшу в рассматриваемой игре не существует, но существует ситуация равновесия по Бержу--Вайсману
\[
U^B=(U_1^B,U_2^B),
\]
причем она имеет линейный позиционный вид
\[
u_1^B(t,x)=D_{21}^{-1}\Theta_2(t,\varepsilon)x,
\qquad
u_2^B(t,x)=\varepsilon D_{12}^{-1}\Theta_1(t,\varepsilon)x.
\]
Равновесные выигрыши:
\[
J_i(U^B,t_0,x_0)=V_i(t_0,x_0)=x_0^{\top}\Theta_i(t_0,\varepsilon)x_0,
\qquad i=1,2.
\]

Малость \(\varepsilon\) важна не только содержательно, но и математически: при \(\varepsilon=0\) система решается проще; затем применяется метод малого параметра Пуанкаре, позволяющий представить решение в виде ряда по степеням \(\varepsilon\):
\[
\Theta_i(t,\varepsilon)=\Theta_i^{(0)}(t)+\varepsilon\Theta_i^{(1)}(t)+\varepsilon^2\Theta_i^{(2)}(t)+\cdots .
\]

\cnexplain
本题最重要的区别是：Berge 均衡不是 Nash 均衡。Nash 中每个玩家通过改变自己的策略来改善自己的收益；Berge 中每个玩家的最大收益由另一个玩家的策略来实现。因此本题的极大化方向是交叉的：
\[
J_1 \text{ по } u_2,
\qquad
J_2 \text{ по } u_1.
\]
模型为
\[
\dot x=A(t)x+u_1+\varepsilon u_2,
\]
其中 \(\varepsilon\) 表示第二个玩家对状态变化的影响较小。设
\[
V_i(t,x)=x^{\top}\Theta_i(t,\varepsilon)x.
\]
交叉最大化给出
\[
u_1^B=D_{21}^{-1}\Theta_2x,
\qquad
u_2^B=\varepsilon D_{12}^{-1}\Theta_1x.
\]
注意三个细节：第一，\(u_1^B\) 由 \(\Theta_2\) 决定；第二，\(u_2^B\) 由 \(\Theta_1\) 决定；第三，\(u_2^B\) 前面有 \(\varepsilon\)。

\corrections
\begin{enumerate}
\item 必须明确 \(\text{Бержу}\ne\text{Нэш}\)。Nash 是 \(J_i\) 对 \(u_i\) 最大化；Berge 是 \(J_i\) 对对方控制最大化。
\item 符号约定不能混用。本文采用公开文章的约定：\(D_{ij}>0\)，负号显式写在泛函里。
\item 均衡控制中不能漏掉小参数：
\[
u_2^B=\varepsilon D_{12}^{-1}\Theta_1x.
\]
\item 严格文献中常称为 равновесие по Бержу--Вайсману；考试题写 Бержу，可以简称，但最好说明严格版本。
\item 矩阵方程中 \(\varepsilon^2\Psi_i\) 项来自小参数影响，不能把题目简化成普通两人 LQ 博弈。
\end{enumerate}

\memory
Модель:
\[
\dot x=A(t)x+u_1+\varepsilon u_2,
\qquad 0\le \varepsilon\ll 1.
\]
Стратегии:
\[
u_i(t,x)=Q_i(t)x.
\]
Равновесие по Бержу:
\[
J_1(U_1^B,U_2)\le J_1(U^B)\quad \forall U_2,
\]
\[
J_2(U_1,U_2^B)\le J_2(U^B)\quad \forall U_1.
\]
Главное отличие от Нэша:
\[
J_1 \text{ максимизируется по } u_2,
\qquad
J_2 \text{ максимизируется по } u_1.
\]
Функции Беллмана:
\[
V_i(t,x)=x^{\top}\Theta_i(t,\varepsilon)x.
\]
Условия первого порядка:
\[
\frac{\partial W_1}{\partial u_2}=0,
\qquad
\frac{\partial W_2}{\partial u_1}=0.
\]
Отсюда
\[
u_1^B(t,x)=D_{21}^{-1}\Theta_2(t,\varepsilon)x,
\]
\[
u_2^B(t,x)=\varepsilon D_{12}^{-1}\Theta_1(t,\varepsilon)x.
\]
При достаточно малом \(\varepsilon\) и условиях
\[
C_2\le 0,
\quad G_2\le 0,
\quad D_{ij}>0,
\]
существует равновесие по Бержу--Вайсману. Выигрыши:
\[
J_i^B(t_0,x_0)=x_0^{\top}\Theta_i(t_0,\varepsilon)x_0.
\]

\newpage
\ticket{20}{Максимум по Парето в двухкритериальной линейно-квадратичной динамической задаче}
\official{20. Максимум по Парето в двухкритериальной линейно-квадратичной динамической задаче.}

\fullanswer
Рассматривается двухкритериальная задача принятия решения при неопределенности или ее динамический линейно-квадратичный аналог. Лицо, принимающее решение, выбирает стратегию или управление, а неопределенный фактор заранее неизвестен.

В статической записи стратегия обозначается
\[
x\in X\subset\R^n,
\]
а неопределенность
\[
y\in Y\subset\R^m.
\]
Есть два критерия:
\[
f(x,y)\quad \text{-- выигрыш},
\]
\[
-R(x,y)\quad \text{-- отрицательный риск}.
\]
Поэтому задача записывается как
\[
G_2=\langle X,Y,\{f(x,y),-R(x,y)\}\rangle.
\]
Содержательно нужно одновременно увеличить выигрыш и уменьшить риск:
\[
f(x,y)\to\max,
\qquad
R(x,y)\to\min.
\]
Так как векторная оптимизация обычно формулируется как максимум всех критериев, второй критерий берется в виде \(-R\).

\subsection*{Линейно-квадратичный выигрыш}
Функция выигрыша имеет вид
\[
f(x,y)=x^{\top}Ax+2x^{\top}By+y^{\top}Dy+2a^{\top}x+2d^{\top}y,
\]
где
\[
A=A^{\top},
\qquad
D=D^{\top}.
\]
Матрицы и векторы имеют согласованные размерности.

\subsection*{Функция риска}
Риск, или сожаление, определяется как потеря выигрыша по сравнению с наилучшим выбором стратегии при уже известной неопределенности \(y\):
\[
R(x,y)=\max_{z\in X}f(z,y)-f(x,y).
\]
Очевидно,
\[
R(x,y)\ge 0.
\]
Если \(R(x,y)=0\), то выбранная стратегия \(x\) оптимальна при данном \(y\).

\subsection*{Минимальная по Парето неопределенность}
Для фиксированной стратегии \(x\) рассматривается двухкритериальная задача по неопределенности
\[
G_2(x)=\langle Y,\{f(x,y),-R(x,y)\}\rangle.
\]
Неопределенность \(y^P(x)\in Y\) называется минимальной по Парето для фиксированного \(x\), если не существует такого \(y\in Y\), которое одновременно ухудшает оба критерия для лица, принимающего решение, то есть нельзя одновременно уменьшить выигрыш и увеличить риск. Формально несовместна система
\[
f(x,y)\le f(x,y^P(x)),
\]
\[
-R(x,y)\le -R(x,y^P(x)),
\]
причем хотя бы одно из неравенств строгое.

Эквивалентно: нельзя найти другую неопределенность, при которой выигрыш не больше, риск не меньше и хотя бы одно ухудшение строгое.

\subsection*{Скаляризация для построения неопределенности}
Обозначим
\[
f[y]=\max_{x\in X}f(x,y).
\]
Пусть существует \(\sigma\in(0,1)\) и функция \(y^P(x):X\to Y\), такая что при каждом \(x\)
\[
y^P(x)\in \arg\min_{y\in Y}\bigl(f(x,y)-\sigma f[y]\bigr).
\]
Тогда \(y^P(x)\) является минимальной по Парето неопределенностью для задачи \(G_2(x)\). Эта конструкция переводит двухкритериальное сравнение в скалярную вспомогательную задачу.

\subsection*{Максимальная по Парето стратегия}
После построения минимальной по Парето неопределенности \(y^P(x)\) рассматривается задача по стратегии:
\[
\langle X,\{f(x,y^P(x)),-R(x,y^P(x))\}\rangle.
\]
Стратегия \(x^P\in X\) называется максимальной по Парето, если не существует \(x\in X\), такого что
\[
f(x,y^P(x))\ge f(x^P,y^P(x^P)),
\]
\[
-R(x,y^P(x))\ge -R(x^P,y^P(x^P)),
\]
и хотя бы одно из этих неравенств строгое. Иначе говоря, нельзя одновременно получить больший выигрыш и меньший риск.

Тройку
\[
(x^P,f^P,R^P)
\]
можно назвать гарантированным по исходам и рискам решением, где
\[
f^P=f(x^P,y^P(x^P)),
\qquad
R^P=R(x^P,y^P(x^P)).
\]

\subsection*{Динамическая линейно-квадратичная версия}
Официальный билет говорит именно о динамической задаче. В динамическом варианте вместо \(f(x,y)\) рассматриваются функционалы
\[
J_1(U,t_0,x_0),
\qquad
J_2(U,t_0,x_0),
\]
заданные на управляемой линейной системе. Типичная форма:
\[
\dot x=A(t)x+\sum_j u_j,
\]
\[
J_i(U,t_0,x_0)=x^{\top}(\vartheta)C_i x(\vartheta)+
\int_{t_0}^{\vartheta}\sum_j u_j^{\top} D_{ij}u_j\,dt,
\qquad i=1,2.
\]
Альтернатива \(U^P\) называется максимальной по Парето, если не существует другой альтернативы \(U\), такой что
\[
J_1(U,t_0,x_0)\ge J_1(U^P,t_0,x_0),
\]
\[
J_2(U,t_0,x_0)\ge J_2(U^P,t_0,x_0),
\]
и хотя бы одно неравенство строгое.

Для построения Pareto-решения используется положительная скаляризация:
\[
J_\alpha(U)=J_1(U)+\alpha J_2(U),
\qquad \alpha>0.
\]
Если \(U^P\) максимизирует \(J_\alpha\), то \(U^P\) является максимальной по Парето альтернативой. Это стандартное достаточное условие Pareto-оптимальности: максимум положительной линейной комбинации критериев дает Pareto-максимальную альтернативу.

В линейно-квадратичном случае после скаляризации снова получается LQ-задача. Вводятся взвешенные матрицы
\[
C(\alpha)=C_1+\alpha C_2,
\]
\[
D_j(\alpha)=D_{1j}+\alpha D_{2j}.
\]
Ищется функция Беллмана
\[
V^P(t,x)=x^{\top}\Theta^P(t)x.
\]
Если соответствующие матрицы имеют нужную знакоопределенность, то Pareto-максимальные управления имеют линейный позиционный вид
\[
u_j^P(t,x)=-D_j^{-1}(\alpha)\Theta^P(t)x.
\]
Матрица \(\Theta^P(t)\) удовлетворяет уравнению типа Риккати
\[
\dot\Theta^P+
\Theta^P A+A^{\top}\Theta^P-
\Theta^P\left(\sum_jD_j^{-1}(\alpha)\right)\Theta^P=0,
\]
с терминальным условием
\[
\Theta^P(\vartheta)=C(\alpha).
\]

Таким образом, максимум по Парето в двухкритериальной LQ-динамической задаче строится как неулучшаемый компромисс между двумя критериями, а технически может быть найден через скаляризацию и решение соответствующей задачи динамического программирования.

\cnexplain
这一题的关键词是：收益 \(f\) 和风险 \(R\)。决策者希望
\[
f\text{ 尽可能大},\qquad R\text{ 尽可能小}.
\]
为了统一成 Pareto 最大化，写成两个准则
\[
\{f,-R\}.
\]
风险函数是 regret：
\[
R(x,y)=\max_{z\in X}f(z,y)-f(x,y).
\]
它表示如果不确定性 \(y\) 已经实现，事后最优选择能得到多少收益，而当前选择 \(x\) 损失了多少。

本题有两个层次。第一，对固定策略 \(x\)，找 Pareto 最小不确定性 \(y^P(x)\)，意思是不能找到另一个 \(y\) 同时让收益更低、风险更高。第二，在这些最坏/保证意义下的不确定性基础上找 Pareto 最大策略 \(x^P\)，意思是不能找到另一个策略同时提高收益并降低风险。

如果是动态 LQ 版本，则把准则写成两个动态泛函 \(J_1(U),J_2(U)\)。常用方法是正权重标量化：
\[
J_\alpha(U)=J_1(U)+\alpha J_2(U),\quad \alpha>0.
\]
若 \(U^P\) 最大化这个加权泛函，则它是 Pareto 最大候选。在线性二次情形下，计算会归结为 Bellman 函数和 Riccati 方程。

\corrections
\begin{enumerate}
\item 必须区分 \(R\) 和 \(-R\)。风险本身要最小化；为了写成 Pareto 最大化，才使用 \(-R\)。
\item 风险不能随便定义，应写成 regret：
\[
R(x,y)=\max_{z\in X}f(z,y)-f(x,y).
\]
\item Pareto 最小不确定性 \(y^P(x)\) 与 Pareto 最大策略 \(x^P\) 是两个不同层次，不能混为一谈。
\item 官方题目强调动态问题，所以不能只停留在静态 \(f(x,y)\)。必须补充动态泛函 \(J_1,J_2\) 和标量化/Riccati 思路。
\item Pareto 解一般不是唯一的。给定权重 \(\alpha\) 后通常得到 Pareto 前沿上的一个代表点。
\end{enumerate}

\memory
Два критерия:
\[
f(x,y)\to\max,
\qquad R(x,y)\to\min.
\]
Пишем как Pareto-максимизацию:
\[
\{f(x,y),-R(x,y)\}.
\]
ЛК-выигрыш:
\[
f(x,y)=x^{\top}Ax+2x^{\top}By+y^{\top}Dy+2a^{\top}x+2d^{\top}y.
\]
Риск:
\[
R(x,y)=\max_{z\in X}f(z,y)-f(x,y)\ge0.
\]
Минимальная по Парето неопределенность \(y^P(x)\): нельзя одновременно уменьшить \(f\) и увеличить \(R\).

Достаточная конструкция:
\[
f[y]=\max_{x\in X}f(x,y),
\]
\[
y^P(x)\in\arg\min_{y\in Y}\bigl(f(x,y)-\sigma f[y]\bigr),
\qquad \sigma\in(0,1).
\]
Стратегия \(x^P\) максимальна по Парето, если не существует \(x\), такого что
\[
f(x,y^P(x))\ge f(x^P,y^P(x^P)),
\]
\[
-R(x,y^P(x))\ge -R(x^P,y^P(x^P)),
\]
и хотя бы одно неравенство строгое.

Динамическая версия:
\[
J_1(U),\quad J_2(U).
\]
Скаляризация:
\[
J_\alpha(U)=J_1(U)+\alpha J_2(U),\qquad \alpha>0.
\]
Если \(U^P\) максимизирует \(J_\alpha\), то \(U^P\) Pareto-максимальна.
В LQ-случае:
\[
V^P(t,x)=x^{\top}\Theta^P(t)x,
\qquad u^P(t,x)=Q^P(t)x,
\]
а \(\Theta^P\) находится из уравнения Риккати.

\section*{Итоговая краткая схема блока F}
\addcontentsline{toc}{section}{Итоговая краткая схема блока F}
\begin{center}
\begin{tabular}{p{2cm}p{12cm}}
\textbf{Билет} & \textbf{Что помнить} \\
18 & Nash в трехлицевой LQ-игре: \(\dot x=Ax+u_1+u_2+u_3\). Ищем \(V_i=x^{\top}\Theta_i x\). При \(D_{ii}<0\): \(u_i^e=-D_{ii}^{-1}\Theta_i x\). Матрицы \(\Theta_i\) решают связанную систему Риккати. \\
19 & Berge в двухлицевой LQ-игре: \(\dot x=Ax+u_1+\varepsilon u_2\). Главное отличие: \(J_1\) максимизируется по \(u_2\), \(J_2\) по \(u_1\). \(u_1^B=D_{21}^{-1}\Theta_2x\), \(u_2^B=\varepsilon D_{12}^{-1}\Theta_1x\). \\
20 & Pareto максимум: одновременно увеличить выигрыш и уменьшить риск. Используем критерии \(\{f,-R\}\), где \(R=\max_z f(z,y)-f(x,y)\). В динамической LQ-задаче строится через скаляризацию \(J_\alpha=J_1+\alpha J_2\) и Riccati. \\
\end{tabular}
\end{center}

\section*{Список использованных материалов}
\addcontentsline{toc}{section}{Список использованных материалов}
\begin{enumerate}
\item Официальная программа государственного экзамена кафедры оптимального управления ВМК МГУ, 2026.
\item Жуковский В.И., Смирнова Л.В., Высокос М.И. \emph{Об одной нерешенной задаче в матричных обыкновенных дифференциальных уравнениях}.
\item Zhukovskiy V.I., Smirnova L.V. \emph{Berge--Vaisman equilibrium for one linear-quadratic differential game}.
\item Жуковский В.И., Мухина Ю.С., Романова В.Э. \emph{Дифференциальная игра N лиц, в которой существует паретовское равновесие угроз и контругроз, но отсутствует равновесие по Нэшу}.
\item Студенческие материалы 2017 г. и краткая версия 2026 г.; использовались как вспомогательные материалы, а не как основной источник.
\end{enumerate}

\blocktitle{G}{Риск, гарантированные решения и равновесия при неопределенности}{Билеты 21--23}
\section*{Общая логика блока G}
\addcontentsline{toc}{section}{Общая логика блока G}
Билеты 21--23 объединены одной схемой:
\[
\text{неопределенность} \;\longrightarrow\; \text{риск / гарантия} \;\longrightarrow\; \text{равновесие по Нэшу или Парето}.
\]

Основные понятия блока:
\begin{itemize}
  \item функция выигрыша $f(x,y)$ или $f_i(x,y)$;
  \item неопределенность $y\in Y$;
  \item риск, или сожаление,
  \[
  R(x,y)=\max_{z\in X} f(z,y)-f(x,y);
  \]
  \item гарантированное решение: выбор стратегии, учитывающий худшую реализацию неопределенности;
  \item Парето-подход: одновременно учитываются несколько критериев, например $f_1,-R_1,f_2,-R_2$.
\end{itemize}

В этом блоке важно не менять направления оптимизации: выигрыши максимизируются, риски минимизируются, поэтому в многокритериальной записи риски входят как $-R_i$.

\section{Билет 21. Официальная формулировка}
\textbf{21. Необходимые и достаточные условия существования стратегий с нулевым риском.}

\subsection{Полная версия ответа}
Рассмотрим однокритериальную задачу принятия решений при неопределенности
\[
\langle X,Y,f(x,y)\rangle,
\]
где $x\in X\subset\R^n$ --- стратегия лица, принимающего решение, $y\in Y\subset\R^m$ --- неопределенный фактор, а $f(x,y)$ --- функция выигрыша. Лицо, принимающее решение, стремится выбрать $x$, обеспечивая большой выигрыш при возможной реализации любого $y\in Y$.

\subsubsection{Функция риска}
Функция риска, или сожаления, определяется формулой
\[
\Phi(x,y)=\max_{z\in X} f(z,y)-f(x,y).
\]
Здесь $\max_{z\in X}f(z,y)$ --- наилучший выигрыш, который можно было бы получить, если бы неопределенность $y$ была известна заранее. Поэтому $\Phi(x,y)$ показывает потерю выигрыша из-за выбора стратегии $x$ при реализации $y$.

Всегда
\[
\Phi(x,y)\ge 0,
\]
так как
\[
\max_{z\in X}f(z,y)\ge f(x,y).
\]
Риск равен нулю тогда и только тогда, когда $x$ оптимальна при данном $y$:
\[
\Phi(x,y)=0
\quad\Longleftrightarrow\quad
f(x,y)=\max_{z\in X}f(z,y).
\]

\subsubsection{Гарантированное по риску решение}
Для фиксированной стратегии $x$ естественно рассматривать наихудший риск:
\[
\max_{y\in Y}\Phi(x,y).
\]
Гарантированная по риску стратегия определяется задачей
\[
\varphi^r=\min_{x\in X}\max_{y\in Y}\Phi(x,y),
\]
\[
x^r\in \argmin_{x\in X}\max_{y\in Y}\Phi(x,y).
\]
Число $\varphi^r$ называется гарантированным риском. При компактности $X,Y$ и непрерывности $f$ соответствующие максимумы и минимумы существуют.

\subsubsection{Стратегия с нулевым риском}
Стратегия $x^0\in X$ называется стратегией с нулевым риском, если
\[
\Phi(x^0,y)=0\qquad \forall y\in Y.
\]
Эквивалентно,
\[
f(x^0,y)=\max_{z\in X}f(z,y)\qquad \forall y\in Y.
\]
То есть одна и та же стратегия $x^0$ должна быть оптимальной при всех возможных реализациях неопределенности.

Главное необходимое и достаточное условие:
\[
\bigcap_{y\in Y}\Argmax_{z\in X} f(z,y)\ne\varnothing.
\]
Если это пересечение непусто, то любой его элемент является стратегией с нулевым риском.

\subsubsection{Теорема об эквивалентных условиях}
Пусть для каждого $y\in Y$ максимум $\max_{z\in X}f(z,y)$ достигается. Тогда следующие условия эквивалентны:
\begin{enumerate}
  \item существует стратегия с нулевым риском:
  \[
  \exists x^0\in X:\quad \Phi(x^0,y)=0\quad \forall y\in Y;
  \]
  \item пересечение множеств оптимальных стратегий непусто:
  \[
  \bigcap_{y\in Y}\Argmax_{z\in X} f(z,y)\ne\varnothing;
  \]
  \item минимаксное значение риска равно нулю:
  \[
  \min_{x\in X}\max_{y\in Y}\Phi(x,y)=0;
  \]
  \item функция риска имеет седловую точку нулевого значения:
  \[
  \max_{y\in Y}\Phi(x^0,y)=\Phi(x^0,y^0)=\min_{x\in X}\Phi(x,y^0)=0.
  \]
\end{enumerate}

Стандартное минимаксное неравенство имеет направление
\[
\max_{y\in Y}\min_{x\in X}\Phi(x,y)\le
\min_{x\in X}\max_{y\in Y}\Phi(x,y).
\]
Это направление важно: обратное неравенство в общем случае неверно.

\subsubsection{Доказательство}
Если существует $x^0$ с нулевым риском, то
\[
\Phi(x^0,y)=0\quad \forall y\in Y.
\]
Так как $\Phi\ge 0$, имеем
\[
\max_{y\in Y}\Phi(x^0,y)=0,
\]
следовательно,
\[
\min_{x\in X}\max_{y\in Y}\Phi(x,y)=0.
\]
Кроме того,
\[
f(x^0,y)=\max_{z\in X}f(z,y)\quad \forall y,
\]
поэтому $x^0$ принадлежит всем множествам $\Argmax_{z\in X}f(z,y)$.

Обратно, если
\[
\bigcap_{y\in Y}\Argmax_{z\in X}f(z,y)\ne\varnothing,
\]
возьмем $x^0$ из этого пересечения. Тогда для любого $y\in Y$
\[
f(x^0,y)=\max_{z\in X}f(z,y),
\]
значит
\[
\Phi(x^0,y)=0\quad \forall y\in Y.
\]
Следовательно, $x^0$ имеет нулевой риск.

\subsection{中文解释}
这一题的核心是：风险不是任意惩罚项，而是“事后最优收益”和“实际收益”的差：
\[
\Phi(x,y)=\max_{z\in X}f(z,y)-f(x,y).
\]
如果某个策略 $x^0$ 对所有可能的 $y$ 都没有后悔值，那么
\[
\Phi(x^0,y)=0\quad \forall y.
\]
这等价于同一个策略 $x^0$ 对所有不确定状态都是最优策略：
\[
x^0\in\Argmax_{z\in X}f(z,y)\quad \forall y\in Y.
\]
所以最本质的条件是：所有不确定状态下的最优策略集合有公共元素。

\subsection{与学生版相比的主要修正}
\begin{enumerate}
  \item 最小最大不等式方向必须写为
  \[
  \max_y\min_x\Phi(x,y)\le \min_x\max_y\Phi(x,y).
  \]
  \item $\Phi$ 的连续性不能简单说成“连续函数之差”。值函数 $\max_{z\in X}f(z,y)$ 连续需要 $X$ 紧、$f$ 连续等条件。
  \item 零风险的核心条件应写成交集非空：
  \[
  \bigcap_y\Argmax_z f(z,y)\ne\varnothing.
  \]
  \item 不能只写“有鞍点”，而要写“有零值鞍点”。
\end{enumerate}

\subsection{Краткая версия}
\[
\Phi(x,y)=\max_{z\in X}f(z,y)-f(x,y)\ge 0.
\]
Стратегия $x^0$ имеет нулевой риск, если
\[
\Phi(x^0,y)=0\quad \forall y\in Y.
\]
Эквивалентно,
\[
f(x^0,y)=\max_{z\in X}f(z,y)\quad \forall y\in Y.
\]
Главное условие:
\[
\bigcap_{y\in Y}\Argmax_{z\in X}f(z,y)\ne\varnothing.
\]
Эквивалентная минимаксная форма:
\[
\min_{x\in X}\max_{y\in Y}\Phi(x,y)=0.
\]
Правильное неравенство:
\[
\max_y\min_x\Phi(x,y)\le \min_x\max_y\Phi(x,y).
\]

\newpage
\section{Билет 22. Официальная формулировка}
\textbf{22. Гарантированное по Парето равновесие в многошаговом варианте дуополии Курно с учетом импорта.}

\subsection{Полная версия ответа}
Рассматривается бескоалиционная игра двух фирм на рынке одного продукта. Игроки выбирают объемы выпуска или корректирующие управления, а импорт выступает как неопределенный фактор.

\subsubsection{Статическая модель Курно с импортом}
В статической модели объемы выпуска фирм обозначаются
\[
q_1,q_2\ge 0,
\]
а объем импорта ---
\[
y\in Y=[0,+\infty).
\]
Суммарный объем товара на рынке:
\[
\bar q=q_1+q_2+y.
\]
Цена линейно зависит от совокупного предложения:
\[
p(\bar q)=a-b\bar q,
\qquad a>0,\quad b>0.
\]
Издержки $i$-й фирмы:
\[
C_i(q_i)=cq_i+d.
\]
Прибыль первой фирмы:
\[
\psi_1(q_1,q_2,y)=\bigl[a-b(q_1+q_2+y)\bigr]q_1-(cq_1+d),
\]
то есть
\[
\psi_1=aq_1-bq_1^2-bq_1q_2-byq_1-cq_1-d.
\]
Прибыль второй фирмы:
\[
\psi_2(q_1,q_2,y)=\bigl[a-b(q_1+q_2+y)\bigr]q_2-(cq_2+d),
\]
\[
\psi_2=aq_2-bq_1q_2-bq_2^2-byq_2-cq_2-d.
\]

Для учета неопределенности вводятся гарантированные функции
\[
\Phi_i(q_1,q_2,y)=\psi_i(q_1,q_2,y)+\frac{y^2}{2},
\qquad i=1,2.
\]

\subsubsection{Общая схема Парето-гарантированного равновесия}
Для фиксированной ситуации $q=(q_1,q_2)$ рассматривается двухкритериальная задача по неопределенности
\[
\left\langle Y,\{\Phi_1(q,y),\Phi_2(q,y)\}\right\rangle.
\]
Сначала строится Парето-минимальная неопределенность
\[
y^P(q).
\]
Затем после подстановки $y=y^P(q)$ получается игра гарантий без неопределенности:
\[
\left\langle \{1,2\},\{X_i\}_{i=1,2},\{\Phi_i(q,y^P(q))\}_{i=1,2}\right\rangle.
\]
Ситуация $q^e=(q_1^e,q_2^e)$ является равновесием Нэша в этой игре, если
\[
\Phi_1(q_1^e,q_2^e,y^P(q^e))\ge
\Phi_1(q_1,q_2^e,y^P(q_1,q_2^e))\quad \forall q_1\ge 0,
\]
\[
\Phi_2(q_1^e,q_2^e,y^P(q^e))\ge
\Phi_2(q_1^e,q_2,y^P(q_1^e,q_2))\quad \forall q_2\ge 0.
\]

\subsubsection{Явный вид в статической модели}
Для модели Курно с импортом получается
\[
y^P(q_1,q_2)=\frac{b(q_1+q_2)}{2}.
\]
Условия первого порядка в игре гарантий дают симметричное равновесие
\[
q_1^e=q_2^e=\frac{a-c}{b(3+b)}.
\]
Соответствующая Парето-минимальная неопределенность:
\[
y^P(q^e)=\frac{a-c}{3+b}.
\]
Прибыль каждой фирмы:
\[
\psi_i^e=\frac{(a-c)^2}{b(3+b)^2}-d,
\qquad i=1,2.
\]
Гарантированные выигрыши:
\[
\Phi_i^e=\psi_i^e+\frac{(y^P(q^e))^2}{2}.
\]

\subsubsection{Многошаговый вариант}
В многошаговой версии состояние системы в момент $k$ задается вектором
\[
q(k)=(q_1(k),q_2(k)).
\]
В двухшаговой позиционной модели рассматриваются моменты $k=0,1$:
\[
q_1(k+1)=\frac{a-c}{2b}-\frac{q_2(k)}{2}-\frac{z(k)}{2}+u_1(k,q(k)),
\]
\[
q_2(k+1)=\frac{a-c}{2b}-\frac{q_1(k)}{2}-\frac{z(k)}{2}+u_2(k,q(k)),
\]
\[
q_1(0)=q_{10},\qquad q_2(0)=q_{20}.
\]
Здесь $u_i(k,q(k))$ --- позиционное управление $i$-й фирмы, $z(k)$ --- неопределенность, связанная с импортом.

Функции выигрыша:
\[
J_i(U,Z,q_0)=[q_i(2)]^2+\sum_{k=0}^{1}\left(q_i(k)-2u_i^2(k)+\frac{z^2(k)}{2}\right),
\qquad i=1,2.
\]

\subsubsection{Определение в многошаговой игре}
Тройка
\[
(U^e,J_1^e[q_0],J_2^e[q_0])
\]
называется гарантированным по Парето равновесием, если существует Парето-минимальная неопределенность $Z^P$, такая что:
\begin{enumerate}
  \item $Z^P$ минимальна по Парето в двухкритериальной задаче по неопределенности;
  \item $U^e=(U_1^e,U_2^e)$ является равновесием Нэша в игре без неопределенности после подстановки $Z=Z^P$;
  \item гарантированные выигрыши равны
  \[
  J_i^e[q_0]=J_i(U^e,Z^P,q_0),\qquad i=1,2.
  \]
\end{enumerate}

\subsubsection{Алгоритм построения}
Многошаговая игра решается обратной индукцией:
\[
k=2\rightarrow k=1\rightarrow k=0.
\]
В конечный момент задаются терминальные функции
\[
V_i^{(2)}(q)=q_i^2,
\qquad i=1,2.
\]
На шаге $k=1$ для фиксированных $q$ и $u$ сначала находится Парето-минимальное значение $z^P(1,q,u)$, затем при нем строятся равновесные управления
\[
u_1^e(1,q),\qquad u_2^e(1,q),
\]
после чего вычисляются функции продолжения $V_i^{(1)}(q)$.

На шаге $k=0$ процедура повторяется с учетом уже найденных функций $V_i^{(1)}(q)$: сначала строится $z^P(0,q,u)$, затем равновесные управления $u_i^e(0,q)$.

Итоговые стратегии:
\[
U_i^e=\{u_i^e(0,q),u_i^e(1,q)\},
\]
а равновесные гарантированные выигрыши равны
\[
J_i^e[q_0]=V_i^{(0)}(q_0),\qquad i=1,2.
\]

\subsubsection{Экономический смысл}
Импорт увеличивает совокупное предложение и снижает цену:
\[
p=a-b(q_1+q_2+y).
\]
Поэтому фирмы должны учитывать не только конкурента, но и внешний неблагоприятный фактор. Парето-гарантированное равновесие означает: сначала выбирается Парето-минимальная неопределенность импорта, затем при ней строится устойчивое по Нэшу поведение фирм.

\subsection{中文解释}
本题的顺序是固定的：
\[
\text{进口不确定性} \rightarrow \text{Pareto-minimal uncertainty} \rightarrow \text{Nash equilibrium in game of guarantees}.
\]
不能先求普通 Cournot Nash，再补充一个进口项。静态模型中价格为
\[
p=a-b(q_1+q_2+y),
\]
利润为
\[
\psi_i=pq_i-(cq_i+d).
\]
文献中的显式结果是
\[
y^P(q)=\frac{b(q_1+q_2)}{2},\quad
q_1^e=q_2^e=\frac{a-c}{b(3+b)},\quad
\psi_i^e=\frac{(a-c)^2}{b(3+b)^2}-d.
\]
多步版本不是连续时间微分博弈，而是离散动态规划问题，应使用反向归纳。

\subsection{与学生版相比的主要修正}
\begin{enumerate}
  \item 学生版中口语化表述必须改成规范定义。
  \item 必须区分利润 $\psi_i$ 与保证函数 $\Phi_i=\psi_i+y^2/2$。
  \item 静态 Cournot 公式不能直接替代多步动态规划算法。
  \item 构造顺序必须是 $Z^P\rightarrow U^e\rightarrow J_i^e$。
  \item 第 22 题是离散多步问题，不是连续时间 PMP 问题。
\end{enumerate}

\subsection{Краткая версия}
Статическая модель:
\[
q_1,q_2\ge0,\\ y\ge0,
\qquad p=a-b(q_1+q_2+y).
\]
Прибыль фирмы:
\[
\psi_i=[a-b(q_1+q_2+y)]q_i-(cq_i+d).
\]
Гарантированный критерий:
\[
\Phi_i=\psi_i+\frac{y^2}{2}.
\]
Сначала строится Парето-минимальная неопределенность:
\[
y^P(q)=\frac{b(q_1+q_2)}{2}.
\]
Затем строится равновесие Нэша в игре гарантий. В статической модели:
\[
q_1^e=q_2^e=\frac{a-c}{b(3+b)},\qquad
\psi_i^e=\frac{(a-c)^2}{b(3+b)^2}-d.
\]
В многошаговой версии:
\[
q_1(k+1)=\frac{a-c}{2b}-\frac{q_2(k)}{2}-\frac{z(k)}{2}+u_1(k,q(k)),
\]
\[
q_2(k+1)=\frac{a-c}{2b}-\frac{q_1(k)}{2}-\frac{z(k)}{2}+u_2(k,q(k)).
\]
Определение:
\[
Z^P \text{ минимальна по Парето},\quad
U^e \text{ --- Нэш при } Z=Z^P,
\]
\[
J_i^e[q_0]=J_i(U^e,Z^P,q_0).
\]
Алгоритм: обратная индукция от $k=2$ к $k=0$.

\newpage
\section{Билет 23. Официальная формулировка}
\textbf{23. Гарантированное по выигрышам и рискам равновесие в линейно-квадратичной игре двух лиц.}

\subsection{Полная версия ответа}
Рассматривается бескоалиционная игра двух лиц при неопределенности
\[
\Gamma_2=\left\langle \{1,2\},X_1,X_2,Y,\{f_i(x,y)\}_{i=1}^2\right\rangle,
\]
где
\[
X_1=X_2=\R^n,
\qquad Y=\R^m.
\]
Стратегии игроков: $x_1\in X_1$, $x_2\in X_2$; неопределенность: $y\in Y$.

\subsubsection{Линейно-квадратичные выигрыши}
Функции выигрыша имеют вид
\[
f_1(x_1,x_2,y)=x_1^\top A_1x_1+2x_1^\top x_2+y^\top D_1y+2a_1^\top x_1+2d_1^\top y,
\]
\[
f_2(x_1,x_2,y)=x_2^\top A_2x_2-2x_2^\top x_1+y^\top D_2y+2a_2^\top x_2+2d_2^\top y.
\]
Предполагается
\[
A_i=A_i^\top<0,
\qquad D_i=D_i^\top>0,
\qquad i=1,2.
\]
Условие $A_i<0$ обеспечивает существование максимума по собственной стратегии игрока, а условие $D_i>0$ используется при построении Парето-минимальной неопределенности.

\subsubsection{Лучшие выигрыши игроков}
Лучший выигрыш первого игрока при фиксированных $x_2,y$:
\[
f_1[x_2,y]=\max_{x_1\in X_1}f_1(x_1,x_2,y).
\]
Так как $A_1<0$, максимум находится из условия
\[
\frac{\partial f_1}{\partial x_1}=2A_1x_1+2x_2+2a_1=0,
\]
откуда
\[
x_1^*(x_2)=-A_1^{-1}(x_2+a_1).
\]
Следовательно,
\[
f_1[x_2,y]=-(x_2+a_1)^\top A_1^{-1}(x_2+a_1)+y^\top D_1y+2d_1^\top y.
\]

Аналогично
\[
f_2[x_1,y]=\max_{x_2\in X_2}f_2(x_1,x_2,y),
\]
\[
\frac{\partial f_2}{\partial x_2}=2A_2x_2-2x_1+2a_2=0,
\]
\[
x_2^*(x_1)=A_2^{-1}(x_1-a_2),
\]
\[
f_2[x_1,y]=-(x_1-a_2)^\top A_2^{-1}(x_1-a_2)+y^\top D_2y+2d_2^\top y.
\]

\subsubsection{Функции риска}
Риск игрока --- это потеря относительно лучшего ответа:
\[
R_1(x_1,x_2,y)=f_1[x_2,y]-f_1(x_1,x_2,y),
\]
\[
R_2(x_1,x_2,y)=f_2[x_1,y]-f_2(x_1,x_2,y).
\]
Всегда
\[
R_i(x_1,x_2,y)\ge 0.
\]
Равенство $R_i=0$ означает, что игрок $i$ выбрал свой лучший ответ.

\subsubsection{Определение гарантированного по выигрышам и рискам равновесия}
Пятерка
\[
(x^e,f_1^P,f_2^P,R_1^P,R_2^P)\in X\times\R^4
\]
называется гарантированным по выигрышам и рискам равновесием, если существует неопределенность $y^P\in Y$, такая что
\[
f_i^P=f_i(x^e,y^P),
\qquad R_i^P=R_i(x^e,y^P),
\qquad i=1,2,
\]
и выполнены два условия:
\begin{enumerate}
  \item $x^e=(x_1^e,x_2^e)$ является равновесием Нэша в игре при фиксированном $y=y^P$:
  \[
  f_1(x_1,x_2^e,y^P)\le f_1(x^e,y^P)\quad \forall x_1,
  \]
  \[
  f_2(x_1^e,x_2,y^P)\le f_2(x^e,y^P)\quad \forall x_2;
  \]
  \item $y^P$ минимальна по Парето в четырехкритериальной задаче
  \[
  \left\langle Y,\{f_1(x^e,y),-R_1(x^e,y),f_2(x^e,y),-R_2(x^e,y)\}\right\rangle.
  \]
\end{enumerate}

Знак минус перед рисками обязателен: выигрыши максимизируются, а риски минимизируются.

\subsubsection{Построение $y^P$}
В данной модели неопределенность входит в выигрыши только через
\[
y^\top D_i y+2d_i^\top y.
\]
Скаляризация четырех критериев приводит к квадратичной задаче по $y$ и дает
\[
y^P=-(D_1+D_2)^{-1}(d_1+d_2).
\]
В данной упрощенной модели $y^P$ не зависит от $x$, поскольку отсутствуют смешанные члены вида $x_i^\top C_i y$.

\subsubsection{Построение $x^e$}
При фиксированном $y=y^P$ условия Нэша имеют вид
\[
2A_1x_1^e+2x_2^e+2a_1=0,
\]
\[
2A_2x_2^e-2x_1^e+2a_2=0.
\]
То есть
\[
A_1x_1^e+x_2^e=-a_1,
\]
\[
-x_1^e+A_2x_2^e=-a_2.
\]
Эта система является самым надежным способом записи ответа, поскольку явные формулы зависят от выбранной конвенции знаков.

При используемой здесь конвенции знаков:
\[
x_1^e=(A_2A_1+E_n)^{-1}(a_2-A_2a_1),
\]
\[
x_2^e=-(A_1A_2+E_n)^{-1}(a_1+A_1a_2).
\]

\subsubsection{Гарантированные выигрыши и риски}
Гарантированные выигрыши:
\[
f_i^P=f_i(x^e,y^P),
\qquad i=1,2.
\]
В развернутом виде:
\[
f_1^P=(x_1^e)^\top A_1x_1^e+2(x_1^e)^\top x_2^e+(y^P)^\top D_1y^P+2a_1^\top x_1^e+2d_1^\top y^P,
\]
\[
f_2^P=(x_2^e)^\top A_2x_2^e-2(x_2^e)^\top x_1^e+(y^P)^\top D_2y^P+2a_2^\top x_2^e+2d_2^\top y^P.
\]
Используя условия Нэша, можно упростить:
\[
f_1^P=-(x_1^e)^\top A_1x_1^e+(y^P)^\top D_1y^P+2d_1^\top y^P,
\]
\[
f_2^P=-(x_2^e)^\top A_2x_2^e+(y^P)^\top D_2y^P+2d_2^\top y^P.
\]
Риски:
\[
R_1^P=f_1[x_2^e,y^P]-f_1(x^e,y^P),
\]
\[
R_2^P=f_2[x_1^e,y^P]-f_2(x^e,y^P).
\]
Так как $x^e$ является равновесием Нэша при $y=y^P$, каждый игрок выбирает лучший ответ. Поэтому в точной равновесной ситуации
\[
R_1^P=R_2^P=0.
\]

\subsection{中文解释}
第 23 题不是只看收益，也不是只看风险，而是同时看四个指标：
\[
f_1,\\ -R_1,\\ f_2,\\ -R_2.
\]
风险 $R_i$ 是后悔值：如果玩家知道对方策略和不确定性，本来可以选最佳反应；实际策略与最佳反应收益的差就是风险。

构造顺序是：
\[
y^P\rightarrow x^e\rightarrow f_i^P,R_i^P.
\]
先找四准则意义下的 Pareto-minimal uncertainty $y^P$，再在 $y=y^P$ 固定后求 Nash 均衡 $x^e$，最后计算保证收益和风险。

\subsection{与学生版相比的主要修正}
\begin{enumerate}
  \item 必须说明 $R_i$ 是后悔值：
  \[
  R_i=\text{best payoff}-\text{actual payoff}.
  \]
  \item 四准则方向必须写为 $\{f_1,-R_1,f_2,-R_2\}$，不能写成 $\{f_1,R_1,f_2,R_2\}$。
  \item Nash 条件是在固定 $y^P$ 后写的。
  \item $A_i<0$ 是为了保证对 $x_i$ 的最大值存在。
  \item 第二个玩家的符号容易出错；最稳妥的写法是先写线性方程组，再给显式公式。
\end{enumerate}

\subsection{Краткая версия}
Выигрыши:
\[
f_1=x_1^\top A_1x_1+2x_1^\top x_2+y^\top D_1y+2a_1^\top x_1+2d_1^\top y,
\]
\[
f_2=x_2^\top A_2x_2-2x_2^\top x_1+y^\top D_2y+2a_2^\top x_2+2d_2^\top y.
\]
Условия:
\[
A_i<0,
\qquad D_i>0.
\]
Лучшие выигрыши:
\[
f_1[x_2,y]=\max_{x_1}f_1(x_1,x_2,y),
\qquad
f_2[x_1,y]=\max_{x_2}f_2(x_1,x_2,y).
\]
Риски:
\[
R_1=f_1[x_2,y]-f_1(x_1,x_2,y),
\]
\[
R_2=f_2[x_1,y]-f_2(x_1,x_2,y).
\]
Гарантированное по выигрышам и рискам равновесие:
\[
(x^e,f_1^P,f_2^P,R_1^P,R_2^P).
\]
Существует $y^P$, такое что
\[
y^P=-(D_1+D_2)^{-1}(d_1+d_2),
\]
а $x^e$ является равновесием Нэша при $y=y^P$:
\[
A_1x_1^e+x_2^e=-a_1,
\]
\[
-x_1^e+A_2x_2^e=-a_2.
\]
Гарантированные выигрыши и риски:
\[
f_i^P=f_i(x^e,y^P),
\qquad
R_i^P=R_i(x^e,y^P).
\]
Так как $x^e$ --- Nash при $y^P$, то
\[
R_1^P=R_2^P=0.
\]
Главная фраза:
\[
\text{сначала Парето-минимальная неопределенность по }\{f_1,-R_1,f_2,-R_2\},
\text{ затем Нэш при фиксированном }y^P.
\]

\newpage
\section*{Итог блока G}
\addcontentsline{toc}{section}{Итог блока G}
\begin{tabular}{p{0.12\textwidth}p{0.78\textwidth}}
\toprule
Билет & Ключевой смысл \\
\midrule
21 & Нулевой риск существует тогда и только тогда, когда есть одна стратегия, оптимальная при всех неопределенностях. \\
22 & В дуополии Курно с импортом сначала строится Парето-минимальная неопределенность импорта, затем равновесие Нэша в игре гарантий. \\
23 & В ЛК-игре двух лиц учитываются одновременно выигрыши и риски: критерии имеют вид $\{f_1,-R_1,f_2,-R_2\}$; после выбора $y^P$ строится Nash. \\
\bottomrule
\end{tabular}

\section*{Использованные материалы}
\addcontentsline{toc}{section}{Использованные материалы}
\begin{enumerate}
  \item Официальный список вопросов кафедры оптимального управления ВМК МГУ на 2026 год.
  \item Студенческий конспект 2017 года по билетам государственного экзамена.
  \item Краткая версия ответов 2026 года.
  \item Жуковский В.И., Кудрявцев К.Н. \textit{Уравновешивание конфликтов при неопределенности. II. Аналог максимина}.
  \item Болтянский В.Г. \textit{Оптимальное управление дискретными системами}. М.: Наука, 1973.
  \item Пропой А.И. \textit{Элементы теории оптимальных дискретных процессов}. М.: Наука, 1973.
\end{enumerate}

\blocktitle{H}{Экономические модели, природные ресурсы и DICE}{Билеты 24--26}
\section*{Замечание об источниках и надежности}
\addcontentsline{toc}{section}{Замечание об источниках и надежности}
Для блока H использованы официальная программа 2026 года, студенческий конспект 2017 года, краткая версия 2026 года и следующие основные источники: Hritonenko--Yatsenko \textit{Mathematical Modeling in Economics, Ecology and the Environment}; Barro--Sala-i-Martin \textit{Economic Growth}; Nordhaus--Sztorc \textit{DICE 2013R: Introduction and User's Manual}; Понтрягин и др. \textit{Математическая теория оптимальных процессов}; Асеев--Кряжимский и Aseev по бесконечному горизонту. Книга Nordhaus \textit{Managing the Global Commons} в текущем наборе материалов отсутствует; для конкретной структуры модели DICE используется руководство Nordhaus--Sztorc DICE 2013R. Книга Seierstad--Sydsaeter в полном виде также отсутствует; для билета 25 основным предметным источником является Hritonenko--Yatsenko, глава о невозобновляемых ресурсах, а оптимизационная часть проверяется по Понтрягину и Асееву.

\section{Билет 24. Официальная формулировка}
\textbf{24. Математические модели технического прогресса (экзогенный и эндогенный прогресс). Задачи оптимизации подходы к решению.}

\subsection{Полная версия ответа}
Технический прогресс в экономических моделях означает рост выпуска, который нельзя объяснить только увеличением количества традиционных факторов производства: капитала $K$, труда $L$ и природных ресурсов. Формально это отражается либо явной зависимостью производственной функции от времени, либо введением отдельного технологического уровня $A(t)$, либо включением человеческого капитала $H(t)$, знаний, опыта, НИОКР и образовательного сектора.

Обычно различают
\[
\text{экзогенный технический прогресс}
\]
и
\[
\text{эндогенный технический прогресс}.
\]
Экзогенный прогресс задается вне модели. Его динамика не объясняется поведением агентов. Эндогенный прогресс возникает внутри модели как результат целенаправленной деятельности: инвестиций в образование, НИОКР, накопление человеческого капитала, обучение, распространение знаний и государственной политики.

\subsubsection*{1. Экзогенный технический прогресс}
В простейшей форме выпуск задается производственной функцией
\[
Y(t)=F(K(t),L(t),t),
\]
или, что более удобно в теории роста,
\[
Y(t)=F(K(t),A(t)L(t)),
\]
где $A(t)$ -- уровень технологии, а $A(t)L(t)$ -- эффективный труд. В модели экзогенного прогресса $A(t)$ задана извне, например
\[
\dot A(t)=gA(t), \qquad A(0)=A_0,
\]
то есть
\[
A(t)=A_0e^{gt}.
\]
Труд также часто считается экзогенным:
\[
\dot L(t)=nL(t), \qquad L(0)=L_0.
\]
Капитал накапливается по уравнению
\[
\dot K(t)=I(t)-\delta K(t),
\]
где $I(t)$ -- инвестиции, $\delta>0$ -- норма амортизации.

Если используется постоянная норма накопления
\[
s=\frac{I(t)}{Y(t)}, \qquad 0<s<1,
\]
то получаем модель Солоу--Свона с экзогенным техническим прогрессом:
\[
\dot K(t)=sF(K(t),A(t)L(t))-\delta K(t).
\]
В переменных на единицу эффективного труда
\[
k(t)=\frac{K(t)}{A(t)L(t)}
\]
при постоянной отдаче от масштаба получается фундаментальное уравнение:
\[
\dot k(t)=sf(k(t))-(\delta+n+g)k(t),
\]
где $f(k)=F(k,1)$. Экономический смысл: технический прогресс $g$ повышает эффективность труда, но одновременно увеличивает ``эффективную амортизацию'' капитала на единицу эффективного труда:
\[
\delta+n+g.
\]
В этой модели $s$, $n$, $g$, $\delta$ заданы, поэтому модель Солоу--Свона является главным образом моделью динамики и сравнительной статики, а не полной задачей оптимального управления.

\subsubsection*{2. Экзогенный прогресс в задаче оптимизации: модель Рамсея}
Если вместо заданной нормы накопления $s$ выбирается потребление $C(t)$ или доля накопления $s(t)$, то возникает задача оптимального управления. Например,
\[
\max_{C(\cdot)}\int_0^\infty e^{-\rho t}U(c(t))L(t)\,dt,
\]
где
\[
c(t)=\frac{C(t)}{L(t)}
\]
-- потребление на душу населения, $\rho>0$ -- коэффициент дисконтирования.

Ограничение:
\[
\dot K(t)=F(K(t),A(t)L(t))-C(t)-\delta K(t),
\]
\[
K(0)=K_0,\qquad C(t)\ge 0,\qquad K(t)\ge 0.
\]
Здесь технический прогресс по-прежнему экзогенен:
\[
A(t)=A_0e^{gt}.
\]
Переходя к переменным на единицу эффективного труда, можно записать:
\[
\dot k(t)=f(k(t))-c_e(t)-(\delta+n+g)k(t),
\]
где $c_e(t)=C(t)/(A(t)L(t))$.

Гамильтониан в current-value форме:
\[
\mathcal H=U(c)+\lambda\left[f(k)-c-(\delta+n+g)k\right].
\]
Условие максимума:
\[
\frac{\partial \mathcal H}{\partial c}=0
\quad\Longrightarrow\quad
U'(c)=\lambda.
\]
Сопряженное уравнение:
\[
\dot\lambda=\lambda\left(\rho+\delta+n+g-f'(k)\right),
\]
в current-value записи, с учетом выбранной нормировки. Условие трансверсальности на бесконечности обычно записывается в виде
\[
\lim_{t\to\infty}e^{-\rho t}\lambda(t)k(t)=0,
\]
или в эквивалентной форме для present-value сопряженной переменной.

Таким образом, в модели Рамсея технический прогресс может быть экзогенным, но задача уже является оптимизационной: выбирается траектория потребления и накопления.

\subsubsection*{3. Эндогенный технический прогресс}
В моделях эндогенного прогресса уровень технологии, человеческий капитал или запас знаний становятся внутренними переменными модели. Их динамика зависит от решений экономических агентов. Общая идея:
\[
Y(t)=F(K(t),L(t),A(t),H(t)),
\]
где $A(t)$ -- запас знаний или технологический уровень, $H(t)$ -- человеческий капитал. В отличие от экзогенной модели, здесь $A(t)$ и $H(t)$ не задаются заранее, а удовлетворяют собственным динамическим уравнениям.

\subsubsection*{4. Модель с человеческим капиталом}
Один из стандартных вариантов -- модель Узавы--Лукаса. Пусть
\[
u(t)\in[0,1]
\]
-- доля человеческого капитала, используемая в производстве. Тогда $1-u(t)$ -- доля, направляемая на образование и накопление человеческого капитала.

Производственная функция:
\[
Y(t)=K^\alpha(t)(u(t)H(t))^{1-\alpha},\qquad 0<\alpha<1.
\]
Динамика физического капитала:
\[
\dot K(t)=Y(t)-C(t)-\delta K(t).
\]
Динамика человеческого капитала:
\[
\dot H(t)=B(1-u(t))H(t),\qquad B>0.
\]
Задача оптимизации:
\[
\max_{C(\cdot),u(\cdot)}\int_0^\infty e^{-\rho t}U(C(t))\,dt
\]
при ограничениях
\[
C(t)\ge 0,\qquad 0\le u(t)\le 1,\qquad K(t)\ge 0,\quad H(t)\ge 0.
\]
Здесь управление $u(t)$ распределяет человеческий капитал между текущим производством и образованием. Если $u(t)$ велико, экономика больше производит сегодня; если $u(t)$ мало, больше ресурсов идет на накопление знаний и будущий рост.

\subsubsection*{5. Модели с НИОКР и накоплением знаний}
Другой класс моделей эндогенного прогресса описывает технологический уровень $A(t)$ как результат НИОКР:
\[
\dot A(t)=\Phi(A(t),L_A(t),A_F(t)),
\]
где $L_A(t)$ -- трудовые ресурсы, направленные в научно-исследовательский сектор, а $A_F(t)$ -- внешний технологический уровень, например уровень страны-лидера.

Простые варианты:
\[
\dot A(t)=\beta L_A(t)A(t),
\]
или с внешними знаниями:
\[
\dot A(t)=\beta L_A(t)(A_F(t)-A(t)).
\]
Разность $A_F(t)-A(t)$ характеризует технологический разрыв между лидером и последователем. Тогда задача оптимизации имеет вид
\[
\max \int_0^\infty e^{-\rho t}U(C(t))\,dt
\]
при ограничениях
\[
Y(t)=F(K(t),A(t)L_Y(t)),
\]
\[
\dot K(t)=Y(t)-C(t)-\delta K(t),
\]
\[
\dot A(t)=\Phi(A(t),L_A(t),A_F(t)),
\]
\[
L_Y(t)+L_A(t)=L(t).
\]
Управления -- $C(t)$, $L_A(t)$, $L_Y(t)$ или доля труда в НИОКР
\[
v(t)=\frac{L_A(t)}{L(t)}.
\]

\subsubsection*{6. Подходы к решению задач оптимизации}
Для моделей технического прогресса используются три основных подхода.

\paragraph{Принцип максимума Понтрягина.}
Для задачи
\[
\max_{u(\cdot)}\int_0^\infty e^{-\rho t}F_0(x(t),u(t))\,dt,
\qquad \dot x=f(x,u),
\]
строится гамильтониан в current-value форме
\[
H(x,u,\psi)=F_0(x,u)+\psi^\top f(x,u).
\]
Необходимые условия:
\[
\dot x=\frac{\partial H}{\partial \psi},\qquad
\dot\psi=\rho\psi-\frac{\partial H}{\partial x},
\]
\[
H(x^*,u^*,\psi)=\max_{u\in U}H(x^*,u,\psi).
\]
Для бесконечного горизонта обязательно требуется условие трансверсальности, например
\[
\lim_{t\to\infty}e^{-\rho t}\psi(t)x(t)=0,
\]
или более точная форма, зависящая от модели и состояния.

\paragraph{Метод Беллмана.}
Если модель стационарна, можно искать функцию ценности
\[
V(x)=\sup_{u(\cdot)}\int_0^\infty e^{-\rho t}F_0(x(t),u(t))\,dt.
\]
Тогда формально она удовлетворяет уравнению Гамильтона--Якоби--Беллмана:
\[
\rho V(x)=\max_{u\in U}\{F_0(x,u)+V_x(x)f(x,u)\}.
\]
Этот метод удобен для построения стационарных обратных связей $u^*=u^*(x)$.

\paragraph{Стационарные траектории и сбалансированный рост.}
В моделях роста часто ищут steady state или balanced growth path. Для переменных на единицу эффективного труда стационарное состояние определяется условиями
\[
\dot k=0,\qquad \dot c=0,
\]
или, в более сложных моделях,
\[
\dot k=0,\qquad \dot h=0,\qquad \dot A/A=\mathrm{const}.
\]
В эндогенных моделях особый интерес представляет возможность положительного долгосрочного темпа роста без экзогенно заданного $g$.

\subsection{Пояснение на китайском}
第24题的核心不是背很多模型名，而是区分两类逻辑：外生技术进步是技术从模型外给定；内生技术进步是技术由模型内的投资、教育、研发产生。

外生技术进步的典型形式是
\[
Y(t)=F(K(t),A(t)L(t)),\qquad \dot A=gA.
\]
这里 $A(t)$ 是技术水平，但它不是决策变量。Solow模型中若储蓄率 $s$ 固定，则
\[
\dot k=sf(k)-(\delta+n+g)k.
\]
Ramsey模型把消费或储蓄变成控制变量：
\[
\max\int_0^\infty e^{-\rho t}U(c(t))dt,
\qquad
\dot K=F(K,AL)-C-\delta K.
\]

内生技术进步中，技术或人力资本成为状态变量。例如
\[
Y=K^\alpha(uH)^{1-\alpha},\qquad \dot H=B(1-u)H.
\]
$u$ 是用于生产的人力资本比例，$1-u$ 是用于教育的人力资本比例。研发模型也可写成
\[
\dot A=\beta L_AA
\]
或
\[
\dot A=\beta L_A(A_F-A).
\]
考试中要特别区分：Solow通常是给定储蓄率的动态模型；Ramsey才是典型的最优控制模型。

\subsection{与学生版相比的主要修正}
\begin{enumerate}
\item 不能把所有含人力资本 $H$ 的模型都说成外生技术进步。如果 $H$ 由教育投资或资源分配决定，例如 $\dot H=B(1-u)H$，它应属于内生增长模型。
\item Solow--Swan 和 Ramsey 的性质必须区分。Solow通常给定 $s=\mathrm{const}$；Ramsey选择 $C(t)$ 或 $s(t)$。
\item 外生技术进步最好写成 $Y=F(K,A(t)L)$，这样单位有效劳动变量下的方程 $\dot k=sf(k)-(\delta+n+g)k$ 才清楚。
\item 无限时域问题必须写横截条件，不能只写 Hamiltonian 和一阶条件。
\item 技术进步分类要完整但不应堆概念，主体应围绕 экзогенный vs. эндогенный。
\end{enumerate}

\section{Билет 25. Официальная формулировка}
\textbf{25. Математические модели оптимальной добычи невозобновляемых ресурсов. Задачи оптимизации подходы к решению.}

\subsection{Полная версия ответа}
Невозобновляемый ресурс -- это ресурс, запас которого за рассматриваемый экономический период не восстанавливается естественным образом. Примеры: нефть, газ, уголь, руды, редкие металлы. Главная особенность таких моделей состоит в том, что текущая добыча уменьшает будущие возможности. Поэтому запас ресурса является фазовой переменной, а интенсивность добычи -- управлением.

Обозначим:
\[
R(t)\ge 0
\]
-- текущий запас ресурса в недрах,
\[
R(0)=R_0>0
\]
-- начальный запас,
\[
E(t)\ge 0
\]
-- интенсивность добычи,
\[
p(t)>0
\]
-- рыночная цена единицы добытого ресурса,
\[
C(E,R)
\]
-- затраты на добычу,
\[
r>0
\]
-- ставка дисконтирования.

Динамика запаса:
\[
\dot R(t)=-E(t).
\]
Обычно предполагается ограничение
\[
0\le E(t)\le E_{\max}.
\]
Экономический смысл: $E(t)>0$ означает добычу и продажу ресурса сейчас; $E(t)=0$ означает сохранение ресурса в недрах для будущей добычи.

\subsubsection*{1. Базовая задача оптимальной добычи}
Типичная задача фирмы состоит в максимизации дисконтированной прибыли:
\[
J(E,T)=\int_0^T e^{-rt}\left[p(t)E(t)-C(E(t),R(t))\right]dt\to \max_{E(\cdot),T}.
\]
Ограничения:
\[
\dot R(t)=-E(t),\qquad R(0)=R_0,
\]
\[
R(t)\ge 0,
\qquad
0\le E(t)\le E_{\max}.
\]
Возможны два варианта горизонта: $T$ задан заранее; либо $T$ свободен, и нужно определить оптимальный момент завершения разработки месторождения. В модели с полным исчерпанием часто задают терминальное условие
\[
R(T)=0.
\]
Если $T$ свободен, задача состоит в нахождении $E^*(t)$, $T^*$, $R^*(t)$ и $J^*$.

\subsubsection*{2. Экономические предположения об издержках}
Обычно предполагается
\[
C_E(E,R)\ge 0,
\]
то есть чем выше интенсивность добычи, тем больше издержки. Также часто предполагается
\[
C_R(E,R)\le 0.
\]
Это означает: чем больше оставшийся запас $R$, тем ниже издержки добычи. Когда ресурс истощается, добыча становится дороже.

В простых моделях возможны частные случаи:
\[
C(E,R)=0,
\]
\[
C(E,R)=C(E),
\]
\[
C(E,R)=cE.
\]

\subsubsection*{3. Принцип максимума Понтрягина}
Запишем гамильтониан в present-value форме:
\[
H(t,R,E,\psi)=e^{-rt}\left[p(t)E-C(E,R)\right]-\psi E.
\]
Здесь $\psi(t)$ -- сопряженная переменная. Экономически она является теневой ценой единицы ресурса в недрах в present-value измерении.

Сопряженное уравнение:
\[
\dot\psi(t)=-\frac{\partial H}{\partial R}=e^{-rt}C_R(E(t),R(t)).
\]
Условие максимума:
\[
E^*(t)\in \argmax_{0\le E\le E_{\max}}\left\{e^{-rt}[p(t)E-C(E,R^*(t))]-\psi(t)E\right\}.
\]
Если максимум внутренний, то
\[
\frac{\partial H}{\partial E}=0,
\]
следовательно,
\[
e^{-rt}\left[p(t)-C_E(E^*(t),R^*(t))\right]-\psi(t)=0.
\]
То есть
\[
\psi(t)=e^{-rt}\left[p(t)-C_E(E^*(t),R^*(t))\right].
\]
В current-value форме положим
\[
\lambda(t)=e^{rt}\psi(t).
\]
Тогда внутреннее условие принимает вид
\[
\lambda(t)=p(t)-C_E(E^*(t),R^*(t)).
\]
Сопряженное уравнение в current-value форме:
\[
\dot\lambda(t)=r\lambda(t)+C_R(E^*(t),R^*(t)).
\]
Если $C_R=0$, то
\[
\dot\lambda=r\lambda.
\]
Это форма правила Хотеллинга для чистой ренты.

\subsubsection*{4. Bang-bang структура при линейной прибыли}
Если $C(E,R)=0$, то гамильтониан линеен по $E$:
\[
H=e^{-rt}p(t)E-\psi E=(e^{-rt}p(t)-\psi(t))E.
\]
Оптимальное управление определяется знаком функции переключения
\[
\Delta(t)=e^{-rt}p(t)-\psi(t).
\]
Если $\Delta(t)>0$, то выгодно добывать максимально:
\[
E^*(t)=E_{\max}.
\]
Если $\Delta(t)<0$, то добычу выгодно отложить:
\[
E^*(t)=0.
\]
Если $\Delta(t)=0$, возможен особый режим или промежуточная добыча.

\subsubsection*{5. Условия на конце}
Если момент окончания $T$ фиксирован и требуется полное исчерпание ресурса,
\[
R(T)=0,
\]
то задача решается как задача с фиксированным концом по фазовой переменной. Если $T$ свободен и нет терминального выигрыша, дополнительно появляется условие трансверсальности по свободному времени:
\[
H(T)=0.
\]
Если конечный запас не фиксирован, то при отсутствии терминальной ценности остатка возникает условие
\[
\psi(T)=0.
\]
Но в типичной задаче разработки месторождения предполагается полное исчерпание $R(T)=0$, поэтому нельзя одновременно без пояснения писать $R(T)=0$ и $\psi(T)=0$: это разные варианты терминальных условий.

\subsubsection*{6. Правило Хотеллинга}
Классическая модель Хотеллинга рассматривает конкурентный рынок невозобновляемого ресурса. Основная идея: ресурс в недрах является капитальным активом. Поэтому его чистая цена должна приносить тот же темп доходности, что и другие активы.

Обозначим чистую цену, или ренту ресурса:
\[
\pi(t)=p(t)-c(t),
\]
где $c(t)$ -- предельные издержки добычи. Правило Хотеллинга:
\[
\frac{\dot\pi(t)}{\pi(t)}=r.
\]
Если издержки нулевые, $c(t)=0$, то
\[
\frac{\dot p(t)}{p(t)}=r,
\qquad
p(t)=p_0e^{rt}.
\]
Смысл: если цена ресурса растет медленнее, чем ставка процента, выгоднее добыть ресурс сейчас и вложить выручку в финансовый актив. Если цена растет быстрее, выгоднее оставить ресурс в недрах.

\subsubsection*{7. Модель Хотеллинга с заданным спросом}
Пусть спрос на ресурс имеет изоэластический вид
\[
E(t)=d_0p(t)^{-\alpha},\qquad \alpha>0.
\]
При правиле Хотеллинга
\[
p(t)=p_0e^{rt}.
\]
Тогда
\[
E(t)=d_0p_0^{-\alpha}e^{-\alpha rt}.
\]
Постоянная $p_0$ определяется из условия исчерпания начального запаса:
\[
\int_0^T E(t)\,dt=R_0.
\]
В результате оптимальная добыча убывает со временем:
\[
E(t)=\frac{r\alpha R_0}{1-e^{-r\alpha T}}e^{-r\alpha t}
\]
для конечного $T$. При бесконечном горизонте:
\[
E(t)=r\alpha R_0e^{-r\alpha t}.
\]

\subsubsection*{8. Модификации модели Хотеллинга}
Классическая модель Хотеллинга слишком упрощена. Поэтому вводятся модификации: издержки добычи $C(E,R)$; технологический прогресс в добыче; неопределенность спроса и цены; открытие новых запасов и разведка; наличие альтернативной технологии; связь добычи с общим экономическим ростом.

\subsubsection*{9. Модель Дасгупты--Хила}
Более общий класс моделей связывает невозобновляемый ресурс с производством и потреблением. Пусть $K(t)$ -- капитал, $R(t)$ -- запас ресурса, $E(t)$ -- добыча ресурса, $C(t)$ -- потребление, $I(t)$ -- инвестиции. Производственная функция:
\[
Y(t)=F(K(t),E(t)).
\]
Баланс выпуска:
\[
Y(t)=C(t)+I(t).
\]
Динамика капитала:
\[
\dot K(t)=I(t)-\delta K(t).
\]
Динамика ресурса:
\[
\dot R(t)=-E(t).
\]
Задача социального планировщика:
\[
\max_{C(\cdot),E(\cdot)}\int_0^\infty e^{-\rho t}U(C(t))\,dt
\]
при ограничениях
\[
C(t)\ge 0,
\quad
E(t)\ge 0,
\quad
K(t)\ge 0,
\quad
R(t)\ge 0.
\]
Экономический смысл: ресурс уже не просто продается как товар, а является фактором производства. Поэтому оптимальная политика должна учитывать компромисс между текущим потреблением, накоплением капитала и сохранением ресурса.

\subsubsection*{10. Подходы к решению}
Основные подходы: принцип максимума Понтрягина, метод Беллмана, анализ особых режимов и переключений, численное решение.

Для стационарной бесконечногоризонтной задачи можно ввести функцию ценности
\[
V(R)=\sup_{E(\cdot)}\int_0^\infty e^{-rt}\left[pE-C(E,R)\right]dt.
\]
Тогда формально:
\[
rV(R)=\max_{0\le E\le E_{\max}}\{pE-C(E,R)-V_R(R)E\}.
\]
Здесь $V_R(R)$ -- текущая теневая цена ресурса.

При сложных функциях $C(E,R)$, $p(t)$, многомерных состояниях или неопределенности задача решается численно: прямой дискретизацией, динамическим программированием, методами нелинейного программирования или shooting method для системы ПМП.

\subsection{Пояснение на китайском}
第25题的核心是：不可再生资源会越采越少，所以状态变量是资源存量 $R(t)$，控制变量是开采强度 $E(t)$，状态方程是
\[
\dot R(t)=-E(t).
\]
基本目标是最大化贴现利润：
\[
J=\int_0^T e^{-rt}[p(t)E(t)-C(E(t),R(t))]dt.
\]
Hamiltonian 为
\[
H=e^{-rt}[pE-C(E,R)]-\psi E.
\]
伴随方程：
\[
\dot\psi=e^{-rt}C_R(E,R).
\]
内部最优条件：
\[
e^{-rt}[p-C_E(E^*,R^*)]=\psi.
\]
经济意义是：当前开采一单位资源的贴现边际收益等于留在地下的一单位资源的影子价值。

Hotelling规则说明地下资源像一种资产。净价格
\[
\pi(t)=p(t)-c(t)
\]
应满足
\[
\frac{\dot\pi(t)}{\pi(t)}=r.
\]
若无开采成本，则 $p(t)=p_0e^{rt}$。考试建议按“模型 -- PMP -- 影子价格 -- Hotelling rule -- 经济解释”的顺序回答。

\subsection{与学生版相比的主要修正}
\begin{enumerate}
\item $C(t)$ 和 $C(E,R)$ 不能混写。正式答案中用 $C(E,R)$ 表示总开采成本；若写单位成本，应另记为 $c(t)$。
\item $R(T)=0$ 和 $\psi(T)=0$ 不能同时乱写。$R(T)=0$ 是终端状态固定；$\psi(T)=0$ 适用于终端剩余资源自由且无残值的情形。
\item Hotelling rule 是不可再生资源最优开采模型的核心，应明确写出 $\dot\pi/\pi=r$。
\item 要区分外生价格的企业局部问题和由 no-arbitrage 与市场需求决定价格的一般均衡模型。
\item $C_R\le 0$ 的经济意义是资源存量越大，开采成本越低；资源越少，开采越困难。
\end{enumerate}

\section{Билет 26. Официальная формулировка}
\textbf{26. Интегрированные модели климата и экономики. Модель DICE. Задачи оптимизации подходы к решению.}

\subsection{Полная версия ответа}
Интегрированные модели климата и экономики, или \textit{integrated assessment models} (IAM), описывают совместную динамику экономической системы и климатической системы. Их задача -- связать
\[
\text{производство}\to \text{энергопотребление}\to \text{выбросы CO}_2
\to \text{концентрация углерода}
\to \text{температура}
\to \text{климатический ущерб}
\to \text{экономическое благосостояние}.
\]
Модель DICE -- это агрегированная глобальная интегрированная модель климата и экономики. Аббревиатура означает
\[
\text{DICE}=\text{Dynamic Integrated model of Climate and the Economy}.
\]
В отличие от чисто климатических моделей, DICE содержит экономический оптимизационный блок. В отличие от обычных моделей экономического роста, DICE добавляет природный капитал климатической системы: концентрация парниковых газов рассматривается как отрицательный природный капитал, а снижение выбросов -- как инвестиции в сохранение климата.

\subsubsection*{1. Назначение IAM и DICE}
Интегрированные модели бывают двух типов: policy evaluation models и policy optimization models. DICE относится главным образом ко второму типу. В ней выбираются такие траектории сбережений и климатической политики, которые максимизируют функцию общественного благосостояния.

Цель DICE -- найти оптимальный компромисс между текущим потреблением и снижением будущего климатического ущерба. Снижение выбросов требует затрат сегодня, но уменьшает будущие потери от изменения климата.

\subsubsection*{2. Целевая функция}
В модели DICE общественное благосостояние задается как дисконтированная сумма полезностей потребления на душу населения:
\[
W=\sum_{t=1}^{T_{\max}}U(c(t),L(t))R(t)\to \max.
\]
Здесь
\[
c(t)=\frac{C(t)}{L(t)}
\]
-- потребление на душу населения, $L(t)$ -- население / труд, $R(t)$ -- дисконтирующий множитель.

Функция полезности имеет вид с постоянной эластичностью предельной полезности:
\[
U(c(t),L(t))=L(t)\frac{c(t)^{1-\alpha}}{1-\alpha},\qquad \alpha\ne 1.
\]
При $\alpha=1$ используется логарифмическая полезность:
\[
U(c,L)=L\ln c.
\]
Дисконтирующий множитель:
\[
R(t)=(1+\rho)^{-t},
\]
где $\rho$ -- чистая норма социального межвременного предпочтения.

\subsubsection*{3. Экономический блок}
Экономический блок основан на неоклассической модели роста. Основная производственная функция имеет кобб-дouglasовский вид:
\[
Y_{\mathrm{gross}}(t)=A(t)K(t)^\gamma L(t)^{1-\gamma}.
\]
Здесь $A(t)$ -- совокупная факторная производительность, $K(t)$ -- капитал, $L(t)$ -- труд / население, $0<\gamma<1$ -- эластичность выпуска по капиталу. В DICE технологический прогресс $A(t)$ и население $L(t)$ являются экзогенными траекториями, а капитал определяется оптимизацией потребления и инвестиций.

Баланс выпуска:
\[
Y(t)=C(t)+I(t),
\]
где $C(t)$ -- потребление, $I(t)$ -- инвестиции. Динамика капитала:
\[
K(t+1)=(1-\delta)^hK(t)+hI(t),
\]
где $h$ -- длина временного шага. В DICE 2013R используется шаг 5 лет. Часто управление задается через норму сбережений:
\[
s(t)=\frac{I(t)}{Y(t)}.
\]
Тогда
\[
C(t)=(1-s(t))Y(t),\qquad I(t)=s(t)Y(t).
\]

\subsubsection*{4. Климатический ущерб и издержки снижения выбросов}
Валовой выпуск уменьшается из-за климатического ущерба и затрат на снижение выбросов. В DICE используется функция ущерба, зависящая от повышения атмосферной температуры:
\[
\Omega(t)=\psi_1T_{\mathrm{AT}}(t)+\psi_2T_{\mathrm{AT}}^2(t).
\]
В DICE 2013R ущерб упрощенно задается квадратичной функцией температуры.

Функция затрат на снижение выбросов имеет вид
\[
\Lambda(t)=\theta_1(t)\mu(t)^{\theta_2}.
\]
Здесь $\mu(t)\in[0,1]$ -- доля сокращения выбросов, или emissions control rate. Если $\mu(t)=0$, выбросы не сокращаются; если $\mu(t)=1$, промышленные выбросы полностью устраняются, если это допускается модельными ограничениями. Функция abatement cost является выпуклой по $\mu(t)$, поэтому предельные затраты сокращения выбросов растут при увеличении контроля.

С учетом ущерба и издержек чистый выпуск можно записать как
\[
Y(t)=\frac{1-\Lambda(t)}{1+\Omega(t)}Y_{\mathrm{gross}}(t),
\]
или эквивалентно как выпуск после вычета ущерба и затрат на сокращение выбросов.

\subsubsection*{5. Блок выбросов}
Промышленные выбросы CO$_2$ задаются как функция выпуска, углеродной интенсивности и контроля выбросов:
\[
E_{\mathrm{ind}}(t)=\sigma(t)(1-\mu(t))Y_{\mathrm{gross}}(t).
\]
Здесь $\sigma(t)$ -- углеродная интенсивность выпуска, $\mu(t)$ -- доля сокращения выбросов. Полные выбросы включают промышленные выбросы и выбросы от землепользования:
\[
E(t)=E_{\mathrm{ind}}(t)+E_{\mathrm{land}}(t).
\]

\subsubsection*{6. Углеродный цикл}
DICE использует трехрезервуарную модель углеродного цикла:
\[
M_{\mathrm{AT}}(t),\quad M_{\mathrm{UP}}(t),\quad M_{\mathrm{LO}}(t),
\]
где $M_{\mathrm{AT}}$ -- углерод в атмосфере, $M_{\mathrm{UP}}$ -- углерод в верхнем океане и биосфере, $M_{\mathrm{LO}}$ -- углерод в глубоком океане.

Уравнения имеют вид линейной дискретной системы:
\[
M_{\mathrm{AT}}(t+1)=\phi_{11}M_{\mathrm{AT}}(t)+\phi_{21}M_{\mathrm{UP}}(t)+E(t),
\]
\[
M_{\mathrm{UP}}(t+1)=\phi_{12}M_{\mathrm{AT}}(t)+\phi_{22}M_{\mathrm{UP}}(t)+\phi_{32}M_{\mathrm{LO}}(t),
\]
\[
M_{\mathrm{LO}}(t+1)=\phi_{23}M_{\mathrm{UP}}(t)+\phi_{33}M_{\mathrm{LO}}(t).
\]
Смысл: выбросы сначала попадают в атмосферу, затем углерод перераспределяется между атмосферой, верхним океаном, биосферой и глубоким океаном.

\subsubsection*{7. Радиационное воздействие}
Концентрация CO$_2$ влияет на радиационный баланс Земли. Радиационное воздействие задается формулой
\[
F(t)=\eta\log_2\left(\frac{M_{\mathrm{AT}}(t)}{M_{\mathrm{AT}}^{1750}}\right)+F_{\mathrm{EX}}(t),
\]
где $M_{\mathrm{AT}}^{1750}$ -- доиндустриальный уровень углерода в атмосфере, $F_{\mathrm{EX}}(t)$ -- внешнее радиационное воздействие от других парниковых газов и факторов.

\subsubsection*{8. Температурный блок}
Климатическая система описывается двумя температурными переменными:
\[
T_{\mathrm{AT}}(t)
\]
-- средняя температура атмосферы / поверхности,
\[
T_{\mathrm{LO}}(t)
\]
-- температура глубокого океана.

Типичная форма уравнений:
\[
T_{\mathrm{AT}}(t+1)=T_{\mathrm{AT}}(t)+\xi_1\left[F(t)-\xi_2T_{\mathrm{AT}}(t)-\xi_3(T_{\mathrm{AT}}(t)-T_{\mathrm{LO}}(t))\right],
\]
\[
T_{\mathrm{LO}}(t+1)=T_{\mathrm{LO}}(t)+\xi_4(T_{\mathrm{AT}}(t)-T_{\mathrm{LO}}(t)).
\]
Первое уравнение описывает нагрев атмосферы под воздействием radiative forcing. Второе описывает медленный перенос тепла в глубокий океан.

\subsubsection*{9. Управления в модели DICE}
Основные управляющие переменные:
\[
s(t)
\]
-- норма сбережений, или investment rate,
\[
\mu(t)
\]
-- норма сокращения выбросов, или emissions control rate.

Иногда вместо $\mu(t)$ можно использовать углеродный налог или carbon price, поскольку в оптимуме carbon price равен предельным издержкам выбросов / social cost of carbon.

\subsubsection*{10. Задача оптимизации DICE}
В общей форме задача DICE:
\[
\max_{\{s(t),\mu(t)\}_{t=1}^{T}}\sum_{t=1}^{T}R(t)L(t)\frac{c(t)^{1-\alpha}}{1-\alpha}
\]
при ограничениях
\[
K(t+1)=(1-\delta)^hK(t)+hI(t),
\]
\[
Y_{\mathrm{gross}}(t)=A(t)K(t)^\gamma L(t)^{1-\gamma},
\]
\[
Y(t)=\frac{1-\Lambda(t)}{1+\Omega(t)}Y_{\mathrm{gross}}(t),
\]
\[
C(t)=Y(t)-I(t),
\]
\[
E_{\mathrm{ind}}(t)=\sigma(t)(1-\mu(t))Y_{\mathrm{gross}}(t),
\]
\[
E(t)=E_{\mathrm{ind}}(t)+E_{\mathrm{land}}(t),
\]
\[
M(t+1)=\Phi M(t)+BE(t),
\]
\[
F(t)=\eta\log_2\frac{M_{\mathrm{AT}}(t)}{M_{\mathrm{AT}}^{1750}}+F_{\mathrm{EX}}(t),
\]
\[
T(t+1)=\Psi T(t)+\Xi F(t),
\]
\[
0\le s(t)\le 1,\qquad 0\le \mu(t)\le \bar\mu(t).
\]
Здесь $M(t)$ -- вектор запасов углерода, $T(t)$ -- вектор температур.

\subsubsection*{11. Подходы к решению}
DICE является дискретной нелинейной оптимизационной задачей большой размерности. На практике строится конечномерная задача
\[
\max_{s_1,\ldots,s_T,\mu_1,\ldots,\mu_T}W
\]
при системе нелинейных ограничений. Обычно используются GAMS / NLP solver / CONOPT, а также Excel Solver.

Теоретически можно записать дискретный аналог принципа максимума или условия Каруша--Куна--Таккера. Состояния:
\[
K,\ M_{\mathrm{AT}},M_{\mathrm{UP}},M_{\mathrm{LO}},T_{\mathrm{AT}},T_{\mathrm{LO}}.
\]
Управления:
\[
s,\ \mu.
\]
Сопряженные переменные имеют экономическую интерпретацию: теневая цена капитала, теневая цена атмосферного углерода, теневая цена температуры, social cost of carbon.

Формально можно ввести функцию ценности
\[
V(K,M,T,t)
\]
и записать дискретное уравнение Беллмана:
\[
V_t(K,M,T)=\max_{s,\mu}\left\{U(c,L)R(t)+V_{t+1}(K',M',T')\right\}.
\]
Однако из-за высокой размерности состояния этот подход сложен численно. Поэтому в DICE обычно применяется прямое нелинейное программирование.

DICE также используется для сравнения сценариев: baseline, optimal, temperature-limited, low discounting, Copenhagen Accord.

\subsubsection*{12. Экономический смысл решения}
Оптимальное решение DICE определяет $s^*(t)$ -- сколько выпуска направлять на инвестиции; $\mu^*(t)$ -- какую долю выбросов сокращать; $C^*(t)$ -- траекторию потребления; $K^*(t)$ -- траекторию капитала; $E^*(t)$ -- траекторию выбросов; $T_{\mathrm{AT}}^*(t)$ -- траекторию глобальной температуры; SCC$(t)$ -- social cost of carbon.

Экономический компромисс: высокая $\mu(t)$ снижает будущий климатический ущерб, но требует текущих затрат; высокая $s(t)$ снижает текущее потребление, но повышает будущий капитал; высокая текущая эмиссия увеличивает сегодняшнее производство, но ухудшает климат в будущем.

\subsection{Пояснение на китайском}
第26题的核心是：DICE把经济增长模型和气候模型连在一起，并选择最优储蓄率和减排率。DICE 的全称是 Dynamic Integrated model of Climate and the Economy。

它的逻辑链是：资本、劳动、技术 $\to$ 产出 $\to$ 排放 $\to$ 碳循环 $\to$ 辐射强迫 $\to$ 温度 $\to$ 经济损害 $\to$ 福利最大化。

经济部分：
\[
Y_{\mathrm{gross}}=AK^\gamma L^{1-\gamma},
\qquad
K(t+1)=(1-\delta)^hK(t)+hI(t).
\]
控制变量之一是储蓄率
\[
s(t)=I(t)/Y(t).
\]

排放部分：
\[
E_{\mathrm{ind}}=\sigma(t)(1-\mu(t))Y_{\mathrm{gross}}(t),
\]
其中 $\mu(t)$ 是减排率。气候部分使用三层碳循环 $M_{\mathrm{AT}},M_{\mathrm{UP}},M_{\mathrm{LO}}$，再通过 radiative forcing 影响温度 $T_{\mathrm{AT}},T_{\mathrm{LO}}$。

损害函数和减排成本分别为
\[
\Omega=\psi_1T_{\mathrm{AT}}+\psi_2T_{\mathrm{AT}}^2,
\qquad
\Lambda=\theta_1(t)\mu(t)^{\theta_2}.
\]
目标函数是社会福利最大化：
\[
W=\sum_t R(t)L(t)\frac{c(t)^{1-\alpha}}{1-\alpha}.
\]
考试中最重要的是明确两个控制变量：$s(t)$ 和 $\mu(t)$。

\subsection{与学生版相比的主要修正}
\begin{enumerate}
\item DICE 不是主要以连续时间 PMP 形式实现的模型，而是离散非线性规划模型。DICE 2013R 使用5年一个 period。
\item 控制变量必须明确写成 $s(t)$ 和 $\mu(t)$。
\item 要区分 gross output 和 net output：$Y_{\mathrm{gross}}=AK^\gamma L^{1-\gamma}$，净产出受损害和减排成本影响。
\item 碳循环不是一个简单状态变量，而是三储层系统 $M_{\mathrm{AT}},M_{\mathrm{UP}},M_{\mathrm{LO}}$。
\item 损害函数随温度升高而增加，减排成本随 $\mu$ 增加而增加；两者都降低当前净产出，但 $\mu$ 同时减少未来排放和损害。
\end{enumerate}

\section*{Итоговая схема блока H}
\addcontentsline{toc}{section}{Итоговая схема блока H}
\begin{longtable}{p{0.14\textwidth}p{0.35\textwidth}p{0.42\textwidth}}
\toprule
Билет & Главная модель & Главный вывод \\
\midrule
24 & Технический прогресс: экзогенный $A(t)$ и эндогенный $A,H$ & Экзогенный прогресс задается извне; эндогенный возникает из образования, НИОКР и накопления знаний. \\
25 & Добыча ресурса: $\dot R=-E$ & Оптимальная добыча сравнивает текущую предельную ренту с теневой ценой ресурса; правило Хотеллинга: $\dot\pi/\pi=r$. \\
26 & DICE IAM: экономика + климат & Оптимально выбираются $s(t)$ и $\mu(t)$, связывая рост, выбросы, углеродный цикл, температуру, ущерб и благосостояние. \\
\bottomrule
\end{longtable}

\blocktitle{I}{Неантагонистические дифференциальные игры, бесконечный горизонт и конкурентная реклама}{Билеты 27--29}
\textbf{Замечание об источниках и надежности.} Для билетов 27--28 используются стандартные результаты теории оптимального управления и дифференциальных игр: принцип максимума Понтрягина для программных стратегий и уравнения Беллмана--Гамильтона--Якоби для позиционных стратегий. Для бесконечного горизонта используются материалы Асеева и Асеева--Кряжимского по бесконечногоризонтному оптимальному управлению, current-value переменным и условиям трансверсальности на бесконечности. Для билета 29 конкретная модель конкурентной рекламы берется по студенческому конспекту и краткой версии 2026 года; текст Sorger (1989) в текущем наборе материалов отсутствует, поэтому специальные утверждения, зависящие от оригинальной статьи, следует считать требующими последующей сверки. Формулы PMP и HJB ниже выводятся непосредственно из указанной модели.

\section{Билет 27}

\textbf{Официальная формулировка:} \emph{Неантагонистические дифференциальные игры. Равновесие по Нэшу в программных и позиционных стратегиях.}

\subsection{Полная версия ответа}

Рассмотрим неантагонистическую дифференциальную игру $N$ лиц на конечном промежутке времени $[t_0,T]$. Состояние системы обозначим
\[
 x(t)\in\R^n,
\]
управление $i$-го игрока:
\[
 u_i(t)\in U_i\subset\R^{m_i},\qquad i=1,\ldots,N.
\]
Динамика задается системой
\[
 \dot x(t)=f(t,x(t),u_1(t),\ldots,u_N(t)),\qquad x(t_0)=x_0.
\]
Функционал выигрыша $i$-го игрока имеет вид
\[
 J_i(u_1,\ldots,u_N)=\int_{t_0}^{T}F_i(t,x(t),u_1(t),\ldots,u_N(t))\,dt+q_i(x(T)).
\]
Здесь $F_i$ -- текущий выигрыш, $q_i$ -- терминальный выигрыш.

Игра называется \textbf{неантагонистической}, если выигрыши игроков не сводятся к одной функции с противоположными знаками. В частности, вообще говоря,
\[
 J_i+J_j\ne0.
\]
Это не означает отсутствия конфликта: интересы игроков могут частично совпадать и частично противоречить друг другу.

\subsubsection{Равновесие по Нэшу}

Ситуация
\[
 u^*=(u_1^*,\ldots,u_N^*)
\]
называется равновесием по Нэшу, если ни один игрок не может увеличить свой выигрыш односторонним отклонением при фиксированных стратегиях остальных игроков:
\[
 J_i(u_i,u_{-i}^*)\le J_i(u_i^*,u_{-i}^*)
 \qquad \forall u_i\in\mathcal U_i,\quad i=1,\ldots,N.
\]
Здесь
\[
 u_{-i}^*=(u_1^*,\ldots,u_{i-1}^*,u_{i+1}^*,\ldots,u_N^*).
\]

\subsubsection{Программные стратегии}

Программная стратегия -- это управление, заданное как функция времени:
\[
 u_i=u_i(t).
\]
Она может зависеть от начальной позиции $(t_0,x_0)$, но не зависит от текущего состояния $x(t)$ в процессе игры. Игрок заранее выбирает весь план действий
\[
 u_i(\cdot):[t_0,T]\to U_i.
\]

Равновесие по Нэшу в программных стратегиях определяется неравенствами
\[
 J_i(u_i,u_{-i}^*)\le J_i(u_i^*,u_{-i}^*)
 \qquad \forall u_i(\cdot),\quad i=1,\ldots,N.
\]

Для вывода необходимых условий для каждого игрока фиксируются равновесные стратегии остальных игроков и рассматривается обычная задача оптимального управления по $u_i$.

Гамильтониан $i$-го игрока:
\[
 H_i(t,x,u_1,\ldots,u_N,\psi_i)=F_i(t,x,u_1,\ldots,u_N)+\langle\psi_i,f(t,x,u_1,\ldots,u_N)\rangle.
\]
У каждого игрока свой функционал, поэтому у каждого игрока своя сопряженная переменная $\psi_i(t)$.

Если $u^*$ является программным равновесием по Нэшу и выполнены стандартные условия регулярности, то для каждого $i$ существует $\psi_i(t)$ такое, что вдоль равновесной траектории $x^*(t)$:
\[
 \dot x^*(t)=f(t,x^*(t),u_1^*(t),\ldots,u_N^*(t)),
\]
\[
 \dot\psi_i(t)=-\frac{\partial H_i}{\partial x}(t,x^*(t),u_1^*(t),\ldots,u_N^*(t),\psi_i(t)),
\]
\[
 \psi_i(T)=q_{i,x}(x^*(T)).
\]
Условие максимума:
\[
 u_i^*(t)\in\argmax_{v_i\in U_i}H_i(t,x^*(t),u_1^*(t),\ldots,u_{i-1}^*(t),v_i,u_{i+1}^*(t),\ldots,u_N^*(t),\psi_i(t)).
\]
Это условие означает: каждый игрок максимизирует свой гамильтониан только по своему управлению $u_i$, считая управления остальных игроков фиксированными. В данной записи используется конвенция максимизации выигрышей; при минимизации затрат знак условий меняется.

\subsubsection{Позиционные стратегии}

Позиционная стратегия -- это стратегия обратной связи:
\[
 u_i=\varphi_i(t,x).
\]
Она зависит от времени и текущего состояния системы. Если игроки используют позиционные стратегии $\varphi^*=(\varphi_1^*,\ldots,\varphi_N^*)$, то равновесная траектория определяется замкнутой системой
\[
 \dot x^*(t)=f(t,x^*(t),\varphi_1^*(t,x^*(t)),\ldots,\varphi_N^*(t,x^*(t))).
\]

Позиционное равновесие по Нэшу означает, что ни один игрок не может увеличить свой выигрыш, заменив свою обратную связь $\varphi_i^*$ на другую допустимую стратегию $\varphi_i$, если остальные игроки сохраняют $\varphi_{-i}^*$.

Для позиционных стратегий особенно важно понятие \textbf{состоятельного} равновесия. Оно означает, что равновесие сохраняется не только из начальной позиции $(t_0,x_0)$, но и из любой промежуточной позиции $(t,x)$. Поэтому оно естественно описывается методом динамического программирования.

Пусть существуют функции выигрыша, или функции Беллмана,
\[
 V_i(t,x),\qquad i=1,\ldots,N.
\]
Терминальные условия:
\[
 V_i(T,x)=q_i(x).
\]
Тогда для состоятельного позиционного равновесия должны выполняться уравнения Гамильтона--Якоби:
\[
 -V_{i,t}(t,x)=\max_{v_i\in U_i}\Bigl\{F_i(t,x,\varphi_{-i}^*(t,x),v_i)+V_{i,x}(t,x)f(t,x,\varphi_{-i}^*(t,x),v_i)\Bigr\}.
\]
Равновесная стратегия $i$-го игрока достигает этот максимум:
\[
 \varphi_i^*(t,x)\in\argmax_{v_i\in U_i}\Bigl\{F_i(t,x,\varphi_{-i}^*(t,x),v_i)+V_{i,x}(t,x)f(t,x,\varphi_{-i}^*(t,x),v_i)\Bigr\}.
\]

\subsubsection{Итог}

В программных стратегиях Nash equilibrium задается заранее выбранными управлениями $u_i(t)$ и проверяется через односторонние отклонения. Необходимые условия имеют вид принципа максимума Понтрягина для каждого игрока. В позиционных стратегиях $u_i=\varphi_i(t,x)$; состоятельное позиционное равновесие описывается системой HJB для всех игроков.

\subsection{中文解释}

第 27 题核心是区分两类 Nash 均衡：程序策略 Nash 和位置策略 Nash。程序策略是
\[
 u_i=u_i(t),
\]
即赛前给定一整条控制计划，不根据当前状态实时调整。对每个玩家固定其他人的策略后，可以写一个 PMP：每个玩家有自己的 Hamiltonian 和自己的伴随变量 $\psi_i$。

位置策略是反馈：
\[
 u_i=\varphi_i(t,x).
\]
它根据当前状态调整。强形式是“состоятельное позиционное равновесие”，即从任何中间状态开始仍然构成 Nash 均衡，因此用 Bellman/HJB 方程描述。

最重要的对应关系是：
\[
 \text{program Nash} \Rightarrow \text{PMP},\qquad \text{feedback Nash} \Rightarrow \text{HJB system}.
\]

\subsection{与学生版相比的主要修正}

\begin{enumerate}
\item “Неантагонистическая”不是“没有冲突”，而是不是零和游戏；一般 $J_i+J_j\ne0$。
\item PMP 中每个玩家都有自己的伴随变量 $\psi_i$，不能只写一个公共的 $\psi$。
\item PMP 条件一般是必要条件，不是自动充分条件；若要充分性，需要额外凹性/凸性条件。
\item 位置均衡最好强调“состоятельное”，即任意子博弈仍为 Nash。
\item 有限时域 HJB 必须带终端条件 $V_i(T,x)=q_i(x)$。
\end{enumerate}

\newpage
\section{Билет 28}

\textbf{Официальная формулировка:} \emph{Дифференциальные игры с бесконечной продолжительностью. Равновесие по Нэшу в программных и позиционных стратегиях.}

\subsection{Полная версия ответа}

Рассмотрим неантагонистическую дифференциальную игру $N$ лиц на бесконечном промежутке времени
\[
 [t_0,+\infty).
\]
Состояние системы:
\[
 x(t)\in X\subset\R^n.
\]
Управление $i$-го игрока:
\[
 u_i(t)\in U_i\subset\R^{m_i},\qquad i=1,\ldots,N.
\]
Динамика:
\[
 \dot x(t)=f(t,x(t),u_1(t),\ldots,u_N(t)),\qquad x(t_0)=x_0.
\]
Выигрыш $i$-го игрока:
\[
 J_i(u_1,\ldots,u_N)=\int_{t_0}^{+\infty}e^{-r(t-t_0)}F_i(t,x(t),u_1(t),\ldots,u_N(t))\,dt,
\]
где $r>0$ -- коэффициент дисконтирования.

Дисконтирование отражает меньшую текущую ценность будущих выигрышей. Однако само по себе условие $r>0$ не гарантирует сходимость функционала при любом $F_i$. Обычно дополнительно предполагается ограниченность или допустимый рост $F_i$, чтобы интеграл был конечным.

\subsubsection{Равновесие в программных стратегиях}

Программная стратегия -- это управление $u_i=u_i(t)$, зависящее от времени и начальной позиции, но не зависящее от текущего состояния. Набор
\[
 u^*=(u_1^*,\ldots,u_N^*)
\]
называется равновесием по Нэшу в программных стратегиях, если
\[
 J_i(u_i,u_{-i}^*)\le J_i(u_i^*,u_{-i}^*)\qquad \forall u_i(\cdot),\quad i=1,\ldots,N.
\]

Для каждого игрока фиксируются стратегии остальных игроков и получается задача оптимального управления на бесконечном горизонте.

Present-value гамильтониан:
\[
 \mathcal H_i(t,x,u,\psi_i)=e^{-r(t-t_0)}F_i(t,x,u_1,\ldots,u_N)+\langle\psi_i,f(t,x,u_1,\ldots,u_N)\rangle.
\]
Необходимые условия PMP:
\[
 \dot x^*(t)=f(t,x^*(t),u_1^*(t),\ldots,u_N^*(t)),
\]
\[
 \dot\psi_i(t)=-\frac{\partial\mathcal H_i}{\partial x}(t,x^*(t),u^*(t),\psi_i(t)),
\]
\[
 u_i^*(t)\in\argmax_{v_i\in U_i}\mathcal H_i(t,x^*(t),u_1^*(t),\ldots,u_{i-1}^*(t),v_i,u_{i+1}^*(t),\ldots,u_N^*(t),\psi_i(t)).
\]

На бесконечном горизонте нет терминального условия вида $\psi_i(T)=q_{i,x}(x(T))$. Вместо него используются условия трансверсальности на бесконечности. Типичная форма:
\[
 \lim_{t\to+\infty}\langle\psi_i(t),x^*(t)\rangle=0,
\]
или более точная модельная форма, например условие для $\langle \psi_i(t),x(t)-x^*(t)\rangle$.

\subsubsection{Current-value переменные}

В экономических задачах удобно перейти от present-value сопряженной переменной $\psi_i(t)$ к current-value переменной $p_i(t)$:
\[
 p_i(t)=e^{r(t-t_0)}\psi_i(t),\qquad \psi_i(t)=e^{-r(t-t_0)}p_i(t).
\]
Current-value гамильтониан:
\[
 H_i(t,x,u,p_i)=F_i(t,x,u_1,\ldots,u_N)+\langle p_i,f(t,x,u_1,\ldots,u_N)\rangle.
\]
Тогда условия принимают вид
\[
 u_i^*(t)\in\argmax_{v_i\in U_i}H_i(t,x^*(t),u_{-i}^*(t),v_i,p_i(t)),
\]
\[
 \dot p_i(t)=r p_i(t)-\frac{\partial H_i}{\partial x}(t,x^*(t),u^*(t),p_i(t)).
\]
Член $r p_i(t)$ является характерным отличием current-value записи бесконечногоризонтной задачи.

\subsubsection{Позиционные стратегии}

Позиционная стратегия имеет вид
\[
 u_i=\varphi_i(t,x)
\]
или, в автономной стационарной задаче,
\[
 u_i=\varphi_i(x).
\]
Набор $\varphi^*=(\varphi_1^*,\ldots,\varphi_N^*)$ называется равновесием по Нэшу в позиционных стратегиях, если ни один игрок не может увеличить свой функционал, заменив только свою обратную связь $\varphi_i^*$, при условии что остальные игроки сохраняют $\varphi_{-i}^*$.

Равновесная траектория задается замкнутой системой
\[
 \dot x^*(t)=f(t,x^*(t),\varphi_1^*(t,x^*(t)),\ldots,\varphi_N^*(t,x^*(t))).
\]

Для состоятельного позиционного равновесия вводятся функции ценности игроков
\[
 V_i(t,x),\qquad i=1,\ldots,N.
\]
В автономном стационарном случае обычно ищут $V_i(x)$ и $\varphi_i^*(x)$. Стационарная система Беллмана имеет вид
\[
 rV_i(x)=\max_{v_i\in U_i}\Bigl\{F_i(x,\varphi_{-i}^*(x),v_i)+V_{i,x}(x)f(x,\varphi_{-i}^*(x),v_i)\Bigr\}.
\]
Равновесная стратегия достигает максимум:
\[
 \varphi_i^*(x)\in\argmax_{v_i\in U_i}\Bigl\{F_i(x,\varphi_{-i}^*(x),v_i)+V_{i,x}(x)f(x,\varphi_{-i}^*(x),v_i)\Bigr\}.
\]

\subsubsection{Отличия от конечного горизонта}

\begin{enumerate}
\item Нет конечного момента $T$ и терминального выигрыша $q_i(x(T))$ в стандартной формулировке.
\item Требуется сходимость бесконечного интеграла $J_i$.
\item В current-value PMP появляется член $r p_i$.
\item Вместо конечного терминального условия появляются условия трансверсальности на бесконечности.
\item В стационарном HJB возникает член $rV_i$.
\end{enumerate}

\subsection{中文解释}

第 28 题是第 27 题的无限时域版本。核心变化是 $T$ 变成 $+\infty$。因此收益泛函变成无穷积分：
\[
 J_i=\int_{t_0}^{\infty}e^{-r(t-t_0)}F_i\,dt.
\]
贴现因子不能自动保证收敛，还需要 $F_i$ 的增长条件。

程序策略仍然是 $u_i(t)$，用 PMP。为了经济解释，常把 present-value 伴随变量 $\psi_i$ 改成 current-value 变量：
\[
 p_i=e^{r(t-t_0)}\psi_i.
\]
然后伴随方程里出现
\[
 \dot p_i=rp_i-H_{i,x}.
\]

位置策略通常是 stationary feedback：
\[
 u_i=\varphi_i(x).
\]
对应的 Bellman 方程是
\[
 rV_i=\max\{F_i+V_{i,x}f\}.
\]
这里的 $rV_i$ 是无限时域 HJB 的关键标志。

\subsection{与学生版相比的主要修正}

\begin{enumerate}
\item 贴现因子 $r>0$ 不能简单说成“自动保证泛函收敛”。
\item $\psi_i$ 和 $p_i$ 的关系必须写对：$p_i=e^{r(t-t_0)}\psi_i$。
\item 无限时域 PMP 必须写无穷远横截条件。
\item HJB 中必须有 $rV_i$，不能照搬有限时域的 $-V_t$ 形式。
\item 自治问题中位置均衡通常写成 stationary feedback $u_i=\varphi_i(x)$。
\end{enumerate}

\newpage
\section{Билет 29}

\textbf{Официальная формулировка:} \emph{Модель конкурентной рекламы с двумя участниками.}

\subsection{Полная версия ответа}

Рассматривается рынок фиксированного размера, на котором конкурируют две фирмы. Переменная
\[
 x(t)\in[0,1]
\]
обозначает долю рынка первой фирмы в момент времени $t$. Тогда $1-x(t)$ -- доля рынка второй фирмы.

Управления игроков:
\[
 u_1(t)\ge0,
 \qquad
 u_2(t)\ge0
\]
-- расходы фирм на рекламу.

Динамика доли рынка первой фирмы задается уравнением
\[
 \dot x(t)=u_1(t)\sqrt{1-x(t)}-u_2(t)\sqrt{x(t)},\qquad x(0)=x_0\in[0,1].
\]
Смысл уравнения:
\begin{itemize}
\item реклама первой фирмы увеличивает ее долю рынка;
\item чем меньше доля первой фирмы, тем больше оставшийся рынок $1-x$, на который она может воздействовать;
\item реклама второй фирмы уменьшает долю первой фирмы;
\item эффект рекламы второй фирмы зависит от текущей доли первой фирмы $x$.
\end{itemize}

Если $u_1,u_2\ge0$, то отрезок $[0,1]$ является естественным фазовым множеством: при $x=0$ имеем $\dot x=u_1\ge0$, а при $x=1$ имеем $\dot x=-u_2\le0$.

\subsubsection{Функционалы выигрыша}

Первая фирма максимизирует
\[
 J_1=\int_0^T\left(q_1x(t)-\frac{c_1}{2}u_1^2(t)\right)e^{-rt}\,dt+e^{-rT}S_1x(T).
\]
Вторая фирма максимизирует
\[
 J_2=\int_0^T\left(q_2(1-x(t))-\frac{c_2}{2}u_2^2(t)\right)e^{-rt}\,dt+e^{-rT}S_2(1-x(T)).
\]
Здесь
\[
 r,q_i,c_i,S_i>0,
 \qquad i=1,2.
\]
Параметры имеют экономический смысл: $q_i$ -- доходность рыночной доли фирмы $i$, $c_i$ -- коэффициент стоимости рекламы, $S_i$ -- терминальная ценность рыночной доли, $r$ -- ставка дисконтирования.

\subsubsection{Программное равновесие}

Программная стратегия -- это управление $u_i=u_i(t)$, не зависящее от текущего состояния. Набор $(u_1^*,u_2^*)$ называется равновесием по Нэшу в программных стратегиях, если
\[
 J_1(u_1,u_2^*)\le J_1(u_1^*,u_2^*)\qquad \forall u_1,
\]
\[
 J_2(u_1^*,u_2)\le J_2(u_1^*,u_2^*)\qquad \forall u_2.
\]

Гамильтониан первой фирмы:
\[
 H_1=e^{-rt}\left(q_1x-\frac{c_1}{2}u_1^2\right)+\Lambda_1\left(u_1\sqrt{1-x}-u_2\sqrt{x}\right).
\]
Гамильтониан второй фирмы:
\[
 H_2=e^{-rt}\left(q_2(1-x)-\frac{c_2}{2}u_2^2\right)+\Lambda_2\left(u_1\sqrt{1-x}-u_2\sqrt{x}\right).
\]

Условия максимума:
\[
 \frac{\partial H_1}{\partial u_1}=-c_1e^{-rt}u_1+\Lambda_1\sqrt{1-x}=0,
\]
откуда
\[
 u_1^*(t)=\frac{\Lambda_1(t)}{c_1}e^{rt}\sqrt{1-x^*(t)}.
\]
Для второго игрока:
\[
 \frac{\partial H_2}{\partial u_2}=-c_2e^{-rt}u_2-\Lambda_2\sqrt{x}=0,
\]
поэтому
\[
 u_2^*(t)=-\frac{\Lambda_2(t)}{c_2}e^{rt}\sqrt{x^*(t)}.
\]
Если явно накладываются ограничения $u_i(t)\ge0$, корректнее использовать проекционную форму:
\[
 u_1^*(t)=\max\left\{0,\frac{\Lambda_1(t)}{c_1}e^{rt}\sqrt{1-x^*(t)}\right\},
\]
\[
 u_2^*(t)=\max\left\{0,-\frac{\Lambda_2(t)}{c_2}e^{rt}\sqrt{x^*(t)}\right\}.
\]

Сопряженные уравнения:
\[
 \dot\Lambda_1=-q_1e^{-rt}+\Lambda_1\left(\frac{u_1^*}{2\sqrt{1-x^*}}+\frac{u_2^*}{2\sqrt{x^*}}\right),
\]
\[
 \dot\Lambda_2=q_2e^{-rt}+\Lambda_2\left(\frac{u_1^*}{2\sqrt{1-x^*}}+\frac{u_2^*}{2\sqrt{x^*}}\right).
\]
Терминальные условия:
\[
 \Lambda_1(T)=e^{-rT}S_1,
\]
\[
 \Lambda_2(T)=-e^{-rT}S_2.
\]

Вместе с уравнением состояния и условием $x(0)=x_0$ это дает краевую задачу для программного равновесия.

\subsubsection{Позиционное равновесие}

Позиционная стратегия имеет вид
\[
 u_i=\varphi_i(t,x).
\]
Ищем состоятельное позиционное равновесие через функции Беллмана. Удобно использовать линейный по $x$ вид:
\[
 V_1(t,x)=e^{-rt}(A_1(t)x+B_1(t)),
\]
\[
 V_2(t,x)=e^{-rt}(A_2(t)(1-x)+B_2(t)).
\]
Тогда
\[
 V_{1x}=e^{-rt}A_1(t),\qquad V_{2x}=-e^{-rt}A_2(t).
\]

Из HJB-условия максимума по $u_1$ получаем
\[
 \varphi_1^*(t,x)=\frac{A_1(t)}{c_1}\sqrt{1-x}.
\]
Для второго игрока:
\[
 \varphi_2^*(t,x)=\frac{A_2(t)}{c_2}\sqrt{x}.
\]

Подстановка обратных связей в уравнения Беллмана дает систему для коэффициентов:
\[
 \dot A_1=rA_1-q_1+\frac{A_1^2}{2c_1}+\frac{A_1A_2}{c_2},
\]
\[
 \dot B_1=rB_1-\frac{A_1^2}{2c_1},
\]
\[
 \dot A_2=rA_2-q_2+\frac{A_1A_2}{c_1}+\frac{A_2^2}{2c_2},
\]
\[
 \dot B_2=rB_2-\frac{A_2^2}{2c_2}.
\]
Терминальные условия:
\[
 A_1(T)=S_1,\\quad B_1(T)=0,
\]
\[
 A_2(T)=S_2,
\quad B_2(T)=0.
\]

Равновесная динамика при позиционных стратегиях:
\[
 \dot x=\varphi_1^*(t,x)\sqrt{1-x}-\varphi_2^*(t,x)\sqrt{x},
\]
то есть
\[
 \dot x=\frac{A_1(t)}{c_1}(1-x)-\frac{A_2(t)}{c_2}x.
\]
Это линейное уравнение по $x$.

\subsubsection{Экономический смысл}

Если первая фирма увеличивает рекламу $u_1$, ее доля рынка растет. Но эффект рекламы уменьшается при $x\to1$, потому что оставшийся рынок $1-x$ становится мал. Если вторая фирма увеличивает рекламу $u_2$, доля первой фирмы падает. Квадратичные затраты $\frac{c_i}{2}u_i^2$ выражают возрастающую предельную стоимость рекламной активности.

Программное равновесие строится через PMP и зависит от начальной траектории. Позиционное равновесие задает обратную связь $\varphi_i(t,x)$ и потому реагирует на изменение состояния в процессе игры.

\subsection{中文解释}

第 29 题是广告竞争微分博弈。状态变量 $x(t)$ 是第一家企业的市场份额，第二家企业的市场份额是 $1-x(t)$。控制 $u_1,u_2$ 是广告投入。

状态方程：
\[
 \dot x=u_1\sqrt{1-x}-u_2\sqrt{x}.
\]
第一家广告使 $x$ 增加，第二家广告使 $x$ 减少。因为模型含有 $\sqrt{x}$ 和 $\sqrt{1-x}$，必须有 $x\in[0,1]$。

程序 Nash 用 PMP。关键公式是
\[
 u_1^*=\frac{\Lambda_1}{c_1}e^{rt}\sqrt{1-x^*},\qquad
 u_2^*=-\frac{\Lambda_2}{c_2}e^{rt}\sqrt{x^*}.
\]
其中第二个式子的负号很重要，因为状态方程中 $u_2$ 前面是负号，且 $\Lambda_2(T)=-e^{-rT}S_2$。

位置 Nash 用 Bellman 方法。设
\[
 V_1=e^{-rt}(A_1x+B_1),\qquad V_2=e^{-rt}(A_2(1-x)+B_2).
\]
则反馈策略为
\[
 \varphi_1^*=\frac{A_1}{c_1}\sqrt{1-x},\qquad
 \varphi_2^*=\frac{A_2}{c_2}\sqrt{x}.
\]

\subsection{与学生版相比的主要修正}

\begin{enumerate}
\item 需要明确 $x(t)\in[0,1]$，否则平方根无意义，也不符合市场份额解释。
\item 程序 Nash 的 $u_2^*$ 应有负号：$u_2^*=-\frac{\Lambda_2}{c_2}e^{rt}\sqrt{x^*}$。
\item PMP 对应程序策略，HJB 对应位置策略，不能混用。
\item $V_1$ 应按 $x$ 写，$V_2$ 应按 $1-x$ 写，不能把两个企业的市场份额方向写反。
\item Sorger (1989) 原文当前缺失；具体闭环结论以后如找到原文需再核对，但本答案中的 PMP 和 HJB 公式可由给定模型直接推出。
\end{enumerate}

\end{document}
